% Integration activities / Evaluation / Remediation (Chapter 2) — Anglais 1ère A — Cours signé OnBuch+
% Programme officiel MINESEC. Compiler avec : tectonic ENA26-integration-evaluation-remediation-26.tex
\newcommand{\DOCMATIERE}{Anglais}
\newcommand{\DOCNIVEAU}{1ère A}
\newcommand{\DOCMODULE}{Chapter 2 : Economic Life and Occupations}
\newcommand{\DOCLECON}{ENA26}
\newcommand{\DOCTITRE}{Integration activities / Evaluation / Remediation (Chapter 2)}
% OnBuch+ — préambule « cours signés » ENRICHI (compilé avec tectonic / XeLaTeX)
% Système visuel « Pop » d'OnBuch+ : fond crème (jamais blanc), bordures encre
% épaisses, ombres « dures » décalées, titres Archivo Black en capitales.
% Version MATHÉMATIQUES : graphes (pgfplots/tikz), boîtes pédagogiques
% (définition, propriété, méthode, attention, le savais-tu, exercice résolu,
% exercice, TP, bilan), page de garde et signature OnBuch+ ancrée en bas.
%
% Macros attendues dans main.tex AVANT \input{preamble} :
%   \DOCMATIERE  \DOCNIVEAU  \DOCMODULE  \DOCLECON (numéro)  \DOCTITRE
\documentclass[11pt]{article}

\usepackage{fontspec}
\usepackage[english,french]{babel} % français = langue principale ; l'anglais via \eng{..} / textebox
\usepackage[a4paper,margin=2cm,top=3.4cm,bottom=2.8cm,headheight=1.4cm,footskip=1.3cm]{geometry}
\usepackage{amsmath,amssymb,amsfonts}
\usepackage{mathtools} % \xmapsto, \coloneqq... souvent utilisés par les modèles
\usepackage{cancel} % simplifications barrées (bilans de forces)
% \abs{z} (module, valeur absolue) et \norm{u} (norme), utilisés par les modèles
\providecommand{\abs}{}\let\abs\relax\DeclarePairedDelimiter{\abs}{\lvert}{\rvert}
\providecommand{\norm}{}\let\norm\relax\DeclarePairedDelimiter{\norm}{\lVert}{\rVert}
\usepackage[table]{xcolor} % [table] : fournit aussi \rowcolors (alternance de lignes)
\usepackage{pagecolor}
\usepackage{tikz}
\usetikzlibrary{arrows.meta,positioning,calc,shapes.geometric,shapes.misc,shapes.arrows,decorations.pathmorphing,decorations.pathreplacing,decorations.markings,patterns,shadows,mindmap,trees,fit,backgrounds,matrix,intersections,angles,quotes}
\usepackage{pgfplots}
\usepackage[version=4]{mhchem} % équations nucléaires \ce{^{235}_{92}U}
\usepackage[european,siunitx]{circuitikz} % schémas électriques
\pgfplotsset{compat=1.18}
\usepgfplotslibrary{fillbetween,groupplots} % aires entre courbes (calcul intégral)
\usepackage{enumitem}
\usepackage{fancyhdr}
% [most] charge toutes les bibliothèques tcolorbox (listings, minted, raster,
% documentation...) et gonfle inutilement la mémoire TeX sur un document long
% (~300 boîtes) : on ne charge que ce qui est réellement utilisé.
\usepackage[skins,breakable]{tcolorbox}
\usepackage{graphicx}
\usepackage{colortbl}
\usepackage{booktabs}
\usepackage{tabularx}
\usepackage{diagbox} % case d'en-tête diagonale des tableaux à double entrée
% Colonnes extensibles centrée / gauche / droite (C, L, R), souvent utilisées par les modèles
\newcolumntype{C}{>{\centering\arraybackslash}X}
\newcolumntype{L}{>{\raggedright\arraybackslash}X}
\newcolumntype{R}{>{\raggedleft\arraybackslash}X}
\usepackage{array}
\usepackage{multicol}
\usepackage{multirow}
\usepackage{siunitx}
% parse-numbers=false : accepte aussi \SI{2,5 \times 10^{-3}}{...}
\sisetup{locale=FR,output-decimal-marker={,},group-minimum-digits=4,parse-numbers=false}
\DeclareSIUnit\ohm{\ensuremath{\symup{\Omega}}} % U+2126 absent de la police
\DeclareSIUnit\division{div} % réglages d'oscilloscope (ms/div, V/div)
\DeclareSIUnit\tr{tr} % tours (vitesse de rotation, ex. \SI{1500}{\tr\per\minute})
\DeclareSIUnit\molar{M} % concentration molaire (ex. \SI{3}{\molar}, \SI{1}{\micro\molar})
\DeclareSIUnit\an{an} % année (le modèle confond parfois avec \year, un compteur TeX) — ex. \SI{5}{\kilo\meter\per\an}
\DeclareSIUnit\FCFA{FCFA} % franc CFA (ex. \SI{500}{\FCFA\per\kilo\gram})
\DeclareSIUnit\degreeBrix{\textdegree Brix} % teneur en sucre d'un sirop/jus (réfractométrie)
\DeclareSIUnit\month{mois} % le modèle utilise parfois \month, un compteur TeX, comme unité de durée
\usepackage{qrcode}
% \degree : commande fréquemment utilisée par les modèles pour un angle ou une
% température en dehors de \SI{}{} (non fournie par les paquets chargés).
\providecommand{\degree}{^{\circ}}
% \qed (fin de démonstration), utilisé par les modèles sans amsthm
\providecommand{\qed}{\hfill$\square$}
% \barre : double barre des valeurs interdites dans un tableau de variations (tkz-tab)
\providecommand{\barre}{\ensuremath{\|}}
% environnement proof (amsthm non chargé) utilisé par les modèles
\ifdefined\proof\else\newenvironment{proof}[1][Démonstration]{\par\noindent\textit{#1.}\ }{\qed\par}\fi
\usepackage{lastpage}
\usepackage{hyperref}
\hypersetup{hidelinks}

% ============================================================
% Palette « Pop » OnBuch+ — mêmes hex que la classe P (plus_ui.dart)
% ============================================================
\definecolor{poporange}{HTML}{F59321}
\definecolor{popdark}{HTML}{E07A0C}
\definecolor{popcream}{HTML}{FFF6E8}
\definecolor{popink}{HTML}{1C1714}
\definecolor{poppurple}{HTML}{7A5AE0}
\definecolor{popgreen}{HTML}{2E9E6A}
\definecolor{poppink}{HTML}{E0457B}
\definecolor{popblue}{HTML}{2F7DE1}
\definecolor{popgold}{HTML}{C98A12}
\definecolor{popmuted}{HTML}{8A7F76}
\definecolor{popline}{HTML}{EADFD2}
\definecolor{popgray}{HTML}{8A7F76}
\colorlet{popgrey}{popgray}
% Alias tolérants (le modèle nomme parfois les couleurs différemment)
\colorlet{popwhite}{white}
\colorlet{popblack}{popink}
\colorlet{popred}{poppink}
\colorlet{popyellow}{popgold}
% Teintes claires (fonds de tableaux/encadrés) — le modèle nomme parfois
% les couleurs avec un suffixe L (ex. popblueL) sans les avoir définies.
\colorlet{poporangeL}{poporange!20}
\colorlet{popdarkL}{popdark!20}
\colorlet{poppurpleL}{poppurple!20}
\colorlet{popgreenL}{popgreen!20}
\colorlet{poppinkL}{poppink!20}
\colorlet{popblueL}{popblue!20}
\colorlet{popgoldL}{popgold!20}
\colorlet{popredL}{popred!20}
\colorlet{popgrayL}{popgray!20}
% Teintes claires pour fonds de tableaux / zones
\colorlet{popblueL}{popblue!10!white}
\colorlet{poporangeL}{poporange!14!white}
\colorlet{popgreenL}{popgreen!12!white}
\colorlet{poppurpleL}{poppurple!12!white}
\colorlet{poppinkL}{poppink!12!white}
\colorlet{popgoldL}{popgold!14!white}

\pagecolor{popcream}

% ============================================================
% Typographie
% ============================================================
\defaultfontfeatures{Ligatures=TeX}
\setmainfont{PlusJakartaSans-Medium.ttf}[Path=../../../fonts/,BoldFont=PlusJakartaSans-Bold.ttf]
% Symboles, API et emojis absents de la police principale : repli sur DejaVu Sans (caractères actifs)
\newfontface\symfont{DejaVuSans.ttf}[Path=../../../fonts/]
\makeatletter
\newcommand\symactive[1]{\catcode#1=13 \begingroup\lccode`\~=#1 \lowercase{\endgroup\def~}{{\symfont\char#1}}}
\newcommand\symblank[1]{\catcode#1=13 \begingroup\lccode`\~=#1 \lowercase{\endgroup\def~}{}}
\newcommand\symhyph[1]{\catcode#1=13 \begingroup\lccode`\~=#1 \lowercase{\endgroup\def~}{\mbox{-}}}
\newcount\symc
\newcommand\symrange[2]{\symc=#1 \@whilenum\symc<#2 \do{\expandafter\symactive\expandafter{\the\symc}\advance\symc by 1 }}
\newcommand\symblankrange[2]{\symc=#1 \@whilenum\symc<#2 \do{\expandafter\symblank\expandafter{\the\symc}\advance\symc by 1 }}
\makeatother
\symrange{"0250}{"02FF}\symactive{"2713}\symactive{"2714}\symactive{"2717}\symactive{"2718}\symactive{"2610}\symactive{"2611}\symactive{"2612}
\symhyph{"2011}\symhyph{"2E1A}\symblankrange{"1F300}{"1FAFF}\symblank{"FE0F}
% Maths en sans-serif (Fira Math) avec chiffres et lettres droites de Plus
% Jakarta, pour que les formules se fondent dans le texte.
\usepackage[math-style=ISO,bold-style=ISO]{unicode-math}
\setmathfont{FiraMath-Regular.otf}[Path=../../../fonts/]
\setmathfont{PlusJakartaSans-Medium.ttf}[Path=../../../fonts/,range={up/num,up/{Latin,latin},\mathup}]
\usepackage{icomma} % virgule décimale sans espace en mode math
% Caractères mathématiques Unicode absents des polices de texte (Plus Jakarta,
% Archivo Black) : rendus par leur équivalent mathématique, en texte comme en
% formule — sinon ils s'affichent en carré vide (ex. « Arithmétique dans ℤ »).
\usepackage{newunicodechar}
\newunicodechar{ℝ}{\ensuremath{\mathbb{R}}}
\newunicodechar{ℤ}{\ensuremath{\mathbb{Z}}}
\newunicodechar{ℂ}{\ensuremath{\mathbb{C}}}
\newunicodechar{ℕ}{\ensuremath{\mathbb{N}}}
\newunicodechar{ℚ}{\ensuremath{\mathbb{Q}}}
\newunicodechar{𝔻}{\ensuremath{\mathbb{D}}}
\newunicodechar{α}{\ensuremath{\alpha}}
\newunicodechar{β}{\ensuremath{\beta}}
\newunicodechar{γ}{\ensuremath{\gamma}}
\newunicodechar{δ}{\ensuremath{\delta}}
\newunicodechar{ε}{\ensuremath{\varepsilon}}
\newunicodechar{θ}{\ensuremath{\theta}}
\newunicodechar{λ}{\ensuremath{\lambda}}
\newunicodechar{μ}{\ensuremath{\mu}}
\newunicodechar{σ}{\ensuremath{\sigma}}
\newunicodechar{τ}{\ensuremath{\tau}}
\newunicodechar{φ}{\ensuremath{\varphi}}
\newunicodechar{ψ}{\ensuremath{\psi}}
\newunicodechar{ω}{\ensuremath{\omega}}
\newunicodechar{Σ}{\ensuremath{\Sigma}}
% majuscules produites par \MakeUppercase dans les titres (α-aminés -> Α)
\newunicodechar{Α}{\ensuremath{\alpha}}
\newunicodechar{Β}{\ensuremath{\beta}}
\newunicodechar{Γ}{\ensuremath{\Gamma}}
\newunicodechar{Δ}{\ensuremath{\Delta}}
\newunicodechar{Ε}{\ensuremath{\varepsilon}}
\newunicodechar{Θ}{\ensuremath{\Theta}}
\newunicodechar{Λ}{\ensuremath{\Lambda}}
\newunicodechar{Μ}{\ensuremath{\mu}}
\newunicodechar{Τ}{\ensuremath{\tau}}
\newunicodechar{Φ}{\ensuremath{\Phi}}
\newunicodechar{Ψ}{\ensuremath{\Psi}}
\newunicodechar{Ω}{\ensuremath{\Omega}}
\newunicodechar{↦}{\ensuremath{\mapsto}}
\newunicodechar{∈}{\ensuremath{\in}}
\newunicodechar{∉}{\ensuremath{\notin}}
\newunicodechar{∧}{\ensuremath{\wedge}}
\newunicodechar{⇔}{\ensuremath{\Leftrightarrow}}
\newunicodechar{⟺}{\ensuremath{\Longleftrightarrow}}
\newunicodechar{⇒}{\ensuremath{\Rightarrow}}
\newunicodechar{⟹}{\ensuremath{\Longrightarrow}}
\newunicodechar{∘}{\ensuremath{\circ}}
\newunicodechar{∪}{\ensuremath{\cup}}
\newunicodechar{∩}{\ensuremath{\cap}}
\newunicodechar{≡}{\ensuremath{\equiv}}
\newunicodechar{⊂}{\ensuremath{\subset}}
\newunicodechar{⊥}{\ensuremath{\perp}}
\newunicodechar{∃}{\ensuremath{\exists}}
\newunicodechar{∀}{\ensuremath{\forall}}
\newunicodechar{∛}{\ensuremath{\sqrt[3]{\,}}}
\newunicodechar{∜}{\ensuremath{\sqrt[4]{\,}}}
\newunicodechar{⃗}{\ensuremath{{}^{\to}}}
\newunicodechar{ⁿ}{\ifmmode^{n}\else\textsuperscript{\ensuremath{n}}\fi}
\newunicodechar{ˣ}{\ifmmode^{x}\else\textsuperscript{\ensuremath{x}}\fi}
\newunicodechar{ᵇ}{\ifmmode^{b}\else\textsuperscript{\ensuremath{b}}\fi}
\newunicodechar{ʳ}{\ifmmode^{r}\else\textsuperscript{\ensuremath{r}}\fi}
\newunicodechar{ᵘ}{\ifmmode^{u}\else\textsuperscript{\ensuremath{u}}\fi}
\newunicodechar{ʸ}{\ifmmode^{y}\else\textsuperscript{\ensuremath{y}}\fi}
\newunicodechar{ᵃ}{\ifmmode^{a}\else\textsuperscript{\ensuremath{a}}\fi}
\newunicodechar{ᵉ}{\ifmmode^{e}\else\textsuperscript{\ensuremath{e}}\fi}
\newunicodechar{ᵏ}{\ifmmode^{k}\else\textsuperscript{\ensuremath{k}}\fi}
\newunicodechar{ᵖ}{\ifmmode^{p}\else\textsuperscript{\ensuremath{p}}\fi}
\newunicodechar{ᴺ}{\ifmmode^{N}\else\textsuperscript{\ensuremath{N}}\fi}
\newunicodechar{ᶜ}{\ifmmode^{c}\else\textsuperscript{\ensuremath{c}}\fi}
\newunicodechar{ʰ}{\ifmmode^{h}\else\textsuperscript{\ensuremath{h}}\fi}
\newunicodechar{ᵠ}{\ifmmode^{q}\else\textsuperscript{\ensuremath{q}}\fi}
\newunicodechar{ₙ}{\ifmmode_{n}\else\textsubscript{\ensuremath{n}}\fi}
\newunicodechar{ₐ}{\ifmmode_{a}\else\textsubscript{\ensuremath{a}}\fi}
\newunicodechar{ᵢ}{\ifmmode_{i}\else\textsubscript{\ensuremath{i}}\fi}
\newunicodechar{ᵣ}{\ifmmode_{r}\else\textsubscript{\ensuremath{r}}\fi}
\newunicodechar{ₖ}{\ifmmode_{k}\else\textsubscript{\ensuremath{k}}\fi}
\newunicodechar{ᵦ}{\ifmmode_{\beta}\else\textsubscript{\ensuremath{\beta}}\fi}
\newunicodechar{ₓ}{\ifmmode_{x}\else\textsubscript{\ensuremath{x}}\fi}
\newfontfamily\popdisplay{ArchivoBlack-Regular.ttf}[Path=../../../fonts/]
\newfontfamily\poplogo{SpaceGrotesk-Bold.ttf}[Path=../../../fonts/]
\newfontfamily\popsemi{PlusJakartaSans-SemiBold.ttf}[Path=../../../fonts/]
\newfontfamily\popxbold{PlusJakartaSans-ExtraBold.ttf}[Path=../../../fonts/]
\newfontfamily\popxboldls{PlusJakartaSans-ExtraBold.ttf}[Path=../../../fonts/,LetterSpace=3.0]

\setlist[enumerate]{leftmargin=1.7em,itemsep=5pt,topsep=4pt}
\setlist[itemize]{leftmargin=1.7em,itemsep=4pt,topsep=4pt,label={\color{poporange}\textbullet}}
\setlength{\parindent}{0pt}
\setlength{\parskip}{5pt}
\renewcommand{\arraystretch}{1.25}

% pgfplots : style Pop par défaut
\pgfplotsset{
  every axis/.append style={axis line style={popink,line width=1pt},tick style={popink},
    grid style={popline,line width=0.5pt},label style={font=\small},tick label style={font=\footnotesize}},
  popcurve/.style={poporange,line width=1.8pt,no markers,smooth},
}
% Même style disponible dans un \draw TikZ simple (hors pgfplots)
\tikzset{popcurve/.style={poporange,line width=1.8pt}}
\tikzset{popgrid/.style={popline,very thin}}
\colorlet{popcurve}{poporange} % le nom du style est parfois utilisé comme couleur

% ============================================================
% Pill / squiggle
% ============================================================
\newcommand{\poppill}[3]{%
  \tikz[baseline=-0.55ex]\node[draw=popink,line width=1.2pt,rounded corners=2.6pt,%
    fill=#2,inner xsep=6.5pt,inner ysep=3pt]{%
    \popxboldls\fontsize{7}{7}\selectfont\color{#3}\MakeUppercase{#1}};%
}
\newcommand{\popsquiggle}[1]{%
  \tikz[baseline=-0.4ex]{
    \draw[#1,line width=2.6pt,line cap=round]
      plot [smooth] coordinates {(0,0.09) (0.34,0.01) (0.68,0.21) (1.02,0.01) (1.36,0.10)};
  }%
}
% Mot-clé surligné (marqueur orange sous le texte)
\newcommand{\cle}[1]{{\popxbold\color{popdark}#1}}

% ============================================================
% En-tête / pied
% ============================================================
\pagestyle{fancy}
\fancyhf{}
\renewcommand{\headrulewidth}{0pt}
\renewcommand{\footrulewidth}{0pt}
\lhead{\raisebox{-0.15ex}{\poplogo\fontsize{13}{13}\selectfont\color{popink}OnBuch\color{poporange}+}}
\rhead{\poppill{\DOCMATIERE\ \textbullet\ \DOCNIVEAU\ \textbullet\ Leçon \DOCLECON}{white}{popink}}
\lfoot{\popxbold\fontsize{7.6}{7.6}\selectfont\color{popmuted}\MakeUppercase{Cours signé OnBuch+}\ \textbullet\ {\popsemi\DOCTITRE}}
\rfoot{\tikz[baseline=-0.6ex]\node[rounded corners=8.5pt,fill=popink,minimum height=17pt,minimum width=17pt,inner xsep=5pt,inner sep=0]{\popxbold\fontsize{8}{8}\selectfont\color{popcream}\thepage\,/\,\pageref*{LastPage}};}

% ============================================================
% Titres
% ============================================================
\newcommand{\chaptitre}[2]{%
  \begin{tcolorbox}[enhanced,colback=poporange,colframe=popink,boxrule=2.6pt,arc=11pt,
      drop shadow=popink,left=15pt,right=15pt,top=12pt,bottom=11pt,before skip=3pt,after skip=18pt]
    \poppill{#2}{popink}{popcream}\par\vspace{9pt}
    {\raggedright\hyphenpenalty=10000\exhyphenpenalty=10000\popdisplay\fontsize{22}{24}\selectfont\color{popcream}\MakeUppercase{#1}\par}
  \end{tcolorbox}
}
% Section numérotée : pastille orange avec le numéro + titre Archivo
\newcounter{popsec}
\newcommand{\coursec}[1]{%
  \stepcounter{popsec}\setcounter{popsub}{0}%
  \par\vspace{16pt}\noindent
  \begin{minipage}{\linewidth}
  \tikz[baseline=-0.7ex]\node[draw=popink,line width=1.6pt,fill=poporange,rounded corners=4pt,minimum size=22pt,inner sep=0]{\popdisplay\fontsize{12}{12}\selectfont\color{popcream}\thepopsec};%
  \hspace{8pt}{\popdisplay\fontsize{15}{17}\selectfont\color{popink}\MakeUppercase{#1}}\par
  \vspace{4pt}\noindent{\color{poporange}\rule{36pt}{3.6pt}}
  \end{minipage}\par\nopagebreak\vspace{8pt}
}
\newcounter{popsub}
\newcommand{\courssub}[1]{%
  \stepcounter{popsub}%
  \par\vspace{9pt}\noindent{\popxbold\fontsize{11.5}{13}\selectfont\color{popink}{\color{poporange}\thepopsec.\thepopsub}\hspace{6pt}#1}\par\nopagebreak\vspace{4pt}
}

% ============================================================
% Boîtes pédagogiques — même recette : fond blanc, bordure épaisse colorée,
% ombre dure encre, badge plein. Titre optionnel [#1].
% ============================================================
\tcbset{popbase/.style={breakable,enhanced,colback=white,boxrule=2.3pt,arc=9pt,
  shadow={3pt}{-3pt}{0pt}{popink},left=14pt,right=14pt,top=11pt,bottom=11pt,before skip=11pt,after skip=15pt,
  fontupper=\popsemi\fontsize{10.6}{14.5}\selectfont\color{popink},
  fontlower=\popsemi\fontsize{10.6}{14.5}\selectfont\color{popink}}}
% Coche dessinée (le glyphe ✓ n'existe pas dans les polices du document)
\newcommand{\popcheck}[1]{\tikz[baseline=-0.5ex]\draw[#1,line width=1.6pt,line cap=round,line join=round](-0.13,0.0)--(-0.04,-0.09)--(0.13,0.1);}
\providecommand{\coche}{\popcheck{popgreen}}
\providecommand{\resultat}[1]{\par\smallskip\noindent\fcolorbox{popgreen}{popgreenL}{\parbox{\dimexpr\linewidth-2\fboxsep-2\fboxrule\relax}{\textbf{Résultat.} #1}}\par\smallskip}
\newcommand{\popboxhead}[3]{\poppill{#1}{#2}{white}\ifx\relax#3\relax\else\hspace{6pt}{\popxbold\fontsize{10.6}{12}\selectfont\color{popink}#3}\fi\par\vspace{7pt}}

\newtcolorbox{aretenir}[1][]{popbase,colframe=popgreen,before upper={\popboxhead{À retenir}{popgreen}{#1}}}
% Texte en anglais (document de lecture, dialogue, extrait) : cadre
% d'étude. Un titre optionnel (source, auteur) s'affiche dans l'en-tête.
\newtcolorbox{textebox}[1][]{popbase,colframe=popblue,before upper={\popboxhead{Text}{popblue}{#1}\begin{otherlanguage}{english}},after upper={\end{otherlanguage}}}
% Marques ✓ / ✗ (forme correcte / fautive) souvent utilisées par les modèles
\providecommand{\cross}{\textcolor{poppink}{\ensuremath{\times}}}
\providecommand{\xmark}{\cross}\providecommand{\crossmark}{\cross}\providecommand{\cmark}{\ensuremath{\checkmark}}
% Ligne de réponse / justification à compléter (inventée par les modèles)
\providecommand{\justification}[1]{\par\noindent\textit{Justification~:} \dotfill\par}
% \textipa (package tipa non chargé) : la police ne couvre pas l'API -> simple passage
\providecommand{\textipa}[1]{#1}
% Soulignement (ulem non chargé) : \uline, \ul
\providecommand{\uline}[1]{\underline{#1}}\providecommand{\ul}[1]{\underline{#1}}
% \glossaire{\item ...} inventée par les modèles : liste de glossaire
\providecommand{\glossaire}[1]{\begin{itemize}#1\end{itemize}}
% \alert (beamer) inventée par les modèles : texte mis en évidence
\providecommand{\alert}[1]{\textbf{\textcolor{poppink}{#1}}}
% variantes ASCII d'ordinaux écrites en macro
\providecommand{\iere}{\textsuperscript{re}}\providecommand{\ieme}{\textsuperscript{e}}\providecommand{\ere}{\textsuperscript{re}}\providecommand{\eme}{\textsuperscript{e}}\providecommand{\er}{\textsuperscript{er}}\providecommand{\ie}{\textsuperscript{e}}
% Ordinaux français écrits en macro par les modèles : 1\ière, 3\ième
\newcommand{\ière}{\textsuperscript{re}}\newcommand{\ième}{\textsuperscript{e}}
% \exemple (inventée par les modèles) : étiquette « Exemple : »
\providecommand{\exemple}{\textbf{Exemple~:}\ }
% \square utilisable aussi hors mode math (cases à cocher)
\AtBeginDocument{\let\squareMath\square\renewcommand{\square}{\ensuremath{\squareMath}}}
% Mot ou phrase en anglais à l'intérieur d'un paragraphe français
\newcommand{\eng}[1]{\ifmmode\text{\textit{#1}}\else\begingroup\selectlanguage{english}\itshape #1\endgroup\fi}
\let\esp\eng % alias toléré
% flèches d'intonation inventées par les modèles
\providecommand{\textsearrow}{}\renewcommand{\textsearrow}{\ensuremath{\searrow}}\providecommand{\textnearrow}{}\renewcommand{\textnearrow}{\ensuremath{\nearrow}}
% \fr{..} : mot français (inventé par les modèles)
\providecommand{\fr}[1]{\textit{#1}}
\newtcolorbox{exemplebox}[1][]{popbase,colframe=poporange,before upper={\popboxhead{Exemple}{poporange}{#1}}}
\newtcolorbox{definition}[1][]{popbase,colframe=popblue,before upper={\popboxhead{Définition}{popblue}{#1}}}
\newtcolorbox{propriete}[1][]{popbase,colframe=poppurple,before upper={\popboxhead{Propriété}{poppurple}{#1}}}
\newtcolorbox{methode}[1][]{popbase,colframe=popgold,before upper={\popboxhead{Méthode}{popgold}{#1}}}
\newtcolorbox{attention}[1][]{popbase,colframe=poppink,before upper={\popboxhead{Attention}{poppink}{#1}}}
\newtcolorbox{experience}[1][]{popbase,colframe=popblue,colback=popblueL,before upper={\popboxhead{Expérience}{popblue}{#1}}}
% « Le savais-tu ? » : carte violette pleine (contexte, culture, Cameroun)
\newtcolorbox{savaistu}[1][]{popbase,colback=poppurple,colframe=popink,
  fontupper=\popsemi\fontsize{10.4}{14.2}\selectfont\color{white},
  before upper={\poppill{Le savais-tu\,?}{white}{poppurple}\ifx\relax#1\relax\else\hspace{6pt}{\popxbold\color{white}#1}\fi\par\vspace{7pt}}}
% Exercice résolu : énoncé en haut, \tcblower puis la solution sur fond crème
\newcounter{popexo}\newcounter{popexr}
\newtcolorbox{exoresolu}[1][]{popbase,colframe=poporange,colbacklower=popcream,
  segmentation style={popink,line width=1.2pt,dashed},
  before upper={\stepcounter{popexr}\popboxhead{Exercice résolu \thepopexr}{poporange}{#1}},
  before lower={\poppill{Solution}{popink}{popcream}\par\vspace{6pt}}}
% Exercice (non résolu en place) — la correction est dans la partie Corrigés
\newtcolorbox{exercice}[1][]{popbase,colframe=popink,
  before upper={\stepcounter{popexo}\popboxhead{Exercice \thepopexo}{popink}{#1}}}
% Corrigé
\newtcolorbox{corrige}[1][]{popbase,colframe=popgreen,colback=popgreenL,
  before upper={\popboxhead{Corrigé}{popgreen}{#1}}}
% Objectifs / prérequis / bilan
\newtcolorbox{objectifs}{popbase,colframe=popink,colback=poporangeL,
  before upper={\popboxhead{Objectifs de la leçon}{popink}{}}}
\newtcolorbox{prerequis}{popbase,colframe=popink,colback=popblueL,
  before upper={\popboxhead{Prérequis}{popblue}{}}}
\newtcolorbox{bilan}{popbase,colframe=popink,colback=white,boxrule=2.8pt,
  before upper={\popboxhead{Fiche bilan}{poporange}{L'essentiel à connaître pour l'examen}}}
% Figure encadrée avec légende
\newtcolorbox{popfigure}[1][]{enhanced,breakable=false,colback=white,colframe=popline,boxrule=1.4pt,arc=8pt,
  left=10pt,right=10pt,top=10pt,bottom=8pt,before skip=10pt,after skip=12pt,halign=center,
  lower separated=false,halign lower=center,
  fontlower=\popsemi\fontsize{9}{11.5}\selectfont\color{popmuted},
  #1}
\newcommand{\legende}[1]{\par\vspace{6pt}{\popsemi\fontsize{9}{11.5}\selectfont\color{popmuted}#1}}

% ============================================================
% Signature OnBuch+
% ============================================================
\newcommand{\poplogoinline}[1]{{\poplogo\fontsize{#1}{#1}\selectfont\color{popink}OnBuch\color{poporange}+}}
\newcommand{\signatureonbuch}{%
  \begin{tcolorbox}[enhanced,colback=popink,colframe=popink,boxrule=2.2pt,arc=10pt,
      drop shadow=poporange,left=14pt,right=14pt,top=10pt,bottom=10pt,before skip=0pt,after skip=0pt]
    \tikz[baseline=-0.7ex]\node[circle,fill=popgreen,draw=popcream,line width=1.3pt,minimum size=22pt,inner sep=0]{\popcheck{white}};%
    \hspace{9pt}\begin{minipage}[c]{0.62\linewidth}
      {\popxbold\fontsize{12}{13}\selectfont\color{popcream}Signé\ \textbullet\ {\poplogo OnBuch\color{poporange}+}}\par\vspace{2pt}
      {\raggedright\popsemi\fontsize{8}{10.5}\selectfont\color{popline}\MakeUppercase{Équipe pédagogique OnBuch+}\\\MakeUppercase{Conforme au programme officiel MINESEC}\par}
    \end{minipage}\hfill
    {\poplogo\fontsize{22}{22}\selectfont\color{popcream}OnBuch\color{poporange}+}
  \end{tcolorbox}%
}

% ============================================================
% Page de garde
% #1 titre · #2 matière · #3 niveau · #4 module · #5 n° leçon · #6 durée ·
% #7 description / objectifs (texte ou itemize)
% ============================================================
\newcommand{\pagegarde}[7]{%
  \thispagestyle{empty}
  \newgeometry{a4paper,margin=1.7cm,top=1.6cm,bottom=1.4cm}%
  \begin{tikzpicture}[remember picture,overlay]
    \fill[poporange!18] ([xshift=-3.2cm,yshift=-3.6cm]current page.north east) circle (4.6cm);
    \fill[poppurple!14] ([xshift=2.4cm,yshift=7.6cm]current page.south west) circle (3.8cm);
  \end{tikzpicture}%
  {\poplogo\fontsize{30}{30}\selectfont\color{popink}OnBuch\color{poporange}+}\hfill
  \raisebox{0.6ex}{\poppill{Cours signé \textbullet\ Premium}{popink}{popcream}}\par
  \vspace{1pt}\popsquiggle{poppurple}\par
  \vspace{8pt}
  \begin{tcolorbox}[enhanced,colback=poporange,colframe=popink,boxrule=2.8pt,arc=13pt,
      drop shadow=popink,left=18pt,right=18pt,top=12pt,bottom=13pt,after skip=10pt]
    \poppill{#2 \textbullet\ #3}{popink}{popcream}\hspace{5pt}\poppill{#4}{white}{popink}\par\vspace{8pt}
    {\popdisplay\fontsize{14}{14}\selectfont\color{popink}LEÇON #5}\par\vspace{4pt}
    {\raggedright\hyphenpenalty=10000\exhyphenpenalty=10000\popdisplay\fontsize{25}{27}\selectfont\color{popcream}\MakeUppercase{#1}\par}
    \vspace{8pt}
    {\popxbold\fontsize{9.5}{9.5}\selectfont\color{popink}\MakeUppercase{Durée conseillée : #6}}
  \end{tcolorbox}
  \begin{tcolorbox}[enhanced,colback=white,colframe=popink,boxrule=2.2pt,arc=10pt,
      drop shadow=popink,left=16pt,right=16pt,top=10pt,bottom=10pt,after skip=8pt]
    \poppill{Dans cette leçon}{poppurple}{white}\par\vspace{5pt}
    \popsemi\fontsize{9.3}{12.6}\selectfont\color{popink}#7
  \end{tcolorbox}
  \noindent\begin{minipage}[t]{0.487\linewidth}
    \begin{tcolorbox}[enhanced,colback=popgreen,colframe=popink,boxrule=2.2pt,arc=10pt,
        drop shadow=popink,left=11pt,right=11pt,top=10pt,bottom=10pt,height=3.1cm,valign=center]
      \tcbox[colback=white,colframe=popink,boxrule=1.3pt,arc=5pt,boxsep=0pt,nobeforeafter,
        left=3pt,right=3pt,top=3pt,bottom=3pt,valign=center]{\qrcode[height=1.9cm]{https://play.google.com/store/apps/details?id=com.luvvix.onbuch&hl=fr}}%
      \hspace{8pt}\begin{minipage}[c]{0.5\linewidth}
        {\popdisplay\fontsize{11}{12.5}\selectfont\color{white}\MakeUppercase{Télécharge OnBuch}}\par\vspace{4pt}
        {\raggedright\popsemi\fontsize{8.4}{11}\selectfont\color{white}Gratuit \textbullet\ Google Play\\Tuteur IA, annales, concours, résultats\par}
      \end{minipage}
    \end{tcolorbox}
  \end{minipage}\hfill
  \begin{minipage}[t]{0.487\linewidth}
    \begin{tcolorbox}[enhanced,colback=poppurple,colframe=popink,boxrule=2.2pt,arc=10pt,
        drop shadow=popink,left=11pt,right=11pt,top=10pt,bottom=10pt,height=3.1cm,valign=center]
      \tikz[baseline=-0.7ex]\node[circle,fill=white,draw=popink,line width=1.3pt,minimum size=1.6cm,inner sep=0]{\popdisplay\fontsize{24}{24}\selectfont\color{poppurple}+};%
      \hspace{8pt}\begin{minipage}[c]{0.6\linewidth}
        {\popdisplay\fontsize{11}{12.5}\selectfont\color{white}\MakeUppercase{Passe à OnBuch+}}\par\vspace{4pt}
        {\raggedright\popsemi\fontsize{8.4}{11}\selectfont\color{white}Tous les cours signés, Tuteur IA illimité, fiches PDF — onglet OnBuch+ de l'app\par}
      \end{minipage}
    \end{tcolorbox}
  \end{minipage}
  \vspace{8pt}
  \popcontenu
  \vfill
  \signatureonbuch
  \restoregeometry
  \newpage
}

% Rangée « Ce cours contient » (page de garde)
\newcommand{\poptile}[3]{% #1 couleur #2 gros texte #3 légende
  \begin{tcolorbox}[enhanced,colback=#1,colframe=popink,boxrule=2pt,arc=9pt,shadow={2.5pt}{-2.5pt}{0pt}{popink},
      left=6pt,right=6pt,top=9pt,bottom=9pt,halign=center,valign=center,height=1.75cm,width=\linewidth,nobeforeafter]
    {\popdisplay\fontsize{12.5}{13}\selectfont\color{white}\MakeUppercase{#2}}\par\vspace{3pt}
    {\centering\popsemi\fontsize{7.8}{9.5}\selectfont\color{white}#3\par}
  \end{tcolorbox}}
\newcommand{\popcontenu}{%
  \par\noindent{\popxboldls\fontsize{7.5}{7.5}\selectfont\color{popmuted}CE COURS SIGNÉ CONTIENT}\par\vspace{7pt}
  \noindent\begin{minipage}[t]{0.235\linewidth}\poptile{popblue}{Cours}{complet, illustré, pas à pas}\end{minipage}\hfill
  \begin{minipage}[t]{0.235\linewidth}\poptile{poporange}{Méthodes}{et exercices résolus}\end{minipage}\hfill
  \begin{minipage}[t]{0.235\linewidth}\poptile{poppink}{Exercices}{type examen + corrigés}\end{minipage}\hfill
  \begin{minipage}[t]{0.235\linewidth}\poptile{popgreen}{Fiche bilan}{l'essentiel à retenir}\end{minipage}\par}

% Sommaire
\newtcolorbox{sommaire}{enhanced,colback=white,colframe=popink,boxrule=2.2pt,arc=10pt,
  drop shadow=popink,left=18pt,right=18pt,top=13pt,bottom=13pt,before skip=6pt,after skip=20pt,
  before upper={\poppill{Sommaire}{poppurple}{white}\par\vspace{9pt}\popsemi\fontsize{10.6}{14.5}\selectfont\color{popink}}}

% Signature de fin de cours — poussée tout en bas de la dernière page
\newcommand{\findecours}{%
  \par\vfill\nopagebreak
  \begin{center}\popsquiggle{poporange}\end{center}\vspace{4pt}
  \signatureonbuch
}
\AtBeginDocument{\def\llbracket{\mathopen{[\mkern-3.5mu[}}\def\rrbracket{\mathclose{]\mkern-3.5mu]}}} % intervalles d'entiers (glyphe ⟦ absent de la police)
% Glyphes absents de Fira Math : redéfinis à partir de glyphes disponibles
\AtBeginDocument{%
  \def\setminus{\mathbin{\backslash}}\let\smallsetminus\setminus
  \def\vdots{\mathord{\vbox{\baselineskip3.2pt\lineskiplimit0pt\kern4pt\hbox{.}\hbox{.}\hbox{.}}}}%
  \def\ddots{\mathinner{\mkern1mu\raise6.5pt\hbox{.}\mkern2mu\raise3.6pt\hbox{.}\mkern2mu\raise0.7pt\hbox{.}\mkern1mu}}%
  \def\triangle{\mathord{\scalebox{0.9}{$\symup{\Delta}$}}}%
  \def\nabla{\mathord{\raisebox{\depth}{\scalebox{1}[-1]{$\symup{\Delta}$}}}}%
  \def\checkmark{\popcheck{popdark}}%
  \def\star{\mathbin{\ast}}\def\bigstar{\mathord{\ast}}%
  \def\diamond{\mathbin{\scalebox{0.6}{$\Diamond$}}}%
  \def\bigcup{\mathop{\mathchoice{\raisebox{-0.3em}{\scalebox{1.7}{$\cup$}}}{\scalebox{1.25}{$\cup$}}{\cup}{\cup}}\displaylimits}%
  \def\bigcap{\mathop{\mathchoice{\raisebox{-0.3em}{\scalebox{1.7}{$\cap$}}}{\scalebox{1.25}{$\cap$}}{\cap}{\cap}}\displaylimits}%
  \def\lesseqqgtr{\mathrel{\lessgtr}}\def\bigoplus{\mathop{\oplus}\displaylimits}%
}
\providecommand{\ssi}{si et seulement si} % abréviation fréquente
\AtBeginDocument{\providecommand{\wideparen}{\overparen}} % arc de cercle

%% FIGFIX-BEGIN v5 — correctifs globaux des figures (docs/onbuch-generation/tools/figfix)
\usepackage{adjustbox}\usepackage{etoolbox}\tcbuselibrary{raster}
% toute figure TikZ/pgfplots plus large que la ligne est réduite à la largeur disponible
\BeforeBeginEnvironment{tikzpicture}{\begin{adjustbox}{max width=\linewidth}}
\AfterEndEnvironment{tikzpicture}{\end{adjustbox}}
% légendes pgfplots lisibles par-dessus les courbes
\pgfplotsset{every axis/.append style={legend style={fill=white,fill opacity=0.92,text opacity=1,draw=black!15,inner sep=3pt}}}
% étiquettes d'arêtes (syntaxe edge["..."]) avec fond blanc
\tikzset{every edge quotes/.append style={fill=white,inner sep=1.2pt}}
%% FIGFIX-END
\begin{document}

\pagegarde{Integration activities / Evaluation / Remediation (Chapter 2)}{Anglais}{1ère A}{Chapter 2 : Economic Life and Occupations}{ENA26}{2 h}{Cette leçon de révision clôt le chapitre 2 « Economic Life and Occupations ». Elle synthétise le vocabulaire et la grammaire du chapitre, propose des exercices progressifs de type examen avec auto-évaluation, et offre des pistes de remédiation pour consolider les acquis.
\vspace{4pt}
\begin{multicols}{2}\small
\begin{itemize}
\item À la fin de cette leçon, je sais mobiliser le vocabulaire des métiers, de l'argent et de la vie économique dans des situations variées
\item je sais réutiliser correctement les points de grammaire du chapitre (present simple vs present continuous, comparatifs et superlatifs, quantifieurs)
\item je sais comprendre un texte ou un dialogue sur la vie économique et répondre à des questions de type examen
\item je sais rédiger un court texte (lettre de candidature, paragraphe descriptif) sur un métier ou une activité économique
\item je sais mener un dialogue sur les professions, les salaires et les projets professionnels
\item je sais m'auto-évaluer à l'aide d'une grille et identifier mes points forts et mes difficultés
\end{itemize}\end{multicols}}

\begin{sommaire}
\begin{multicols}{2}
\begin{enumerate}
\item Warm-up : Chapter 2 mind-map
\item Vocabulary synthesis : economic life
\item Grammar synthesis : key points
\item Reading \& listening practice (exam-type)
\item Speaking \& writing consolidation
\item Evaluation, self-assessment \& remediation
\item Activité d'intégration
\item Méthodes et exercices résolus
\item Exercices
\item Corrigés des exercices
\item Fiche bilan
\end{enumerate}
\end{multicols}
\end{sommaire}

\begin{objectifs}
\begin{itemize}
\item À la fin de cette leçon, je sais mobiliser le vocabulaire des métiers, de l'argent et de la vie économique dans des situations variées
\item je sais réutiliser correctement les points de grammaire du chapitre (present simple vs present continuous, comparatifs et superlatifs, quantifieurs)
\item je sais comprendre un texte ou un dialogue sur la vie économique et répondre à des questions de type examen
\item je sais rédiger un court texte (lettre de candidature, paragraphe descriptif) sur un métier ou une activité économique
\item je sais mener un dialogue sur les professions, les salaires et les projets professionnels
\item je sais m'auto-évaluer à l'aide d'une grille et identifier mes points forts et mes difficultés
\item je sais remédier à mes erreurs grâce à des exercices ciblés
\item je sais gérer mon temps et mes stratégies face à une épreuve de type Bac
\end{itemize}
\end{objectifs}

\begin{prerequis}
\begin{itemize}
\item Vocabulaire des métiers et occupations vu dans le chapitre (job, occupation, salary, employer, employee...)
\item Present simple et present continuous (habitudes vs actions en cours)
\item Comparatifs et superlatifs des adjectifs
\item Quantifieurs : much, many, a lot of, few, little
\item Structures de base de la lettre formelle vues en Seconde
\end{itemize}
\end{prerequis}

\newpage

\coursec{Warm-up : Chapter 2 mind-map}

\courssub{Opening situation : Amina's challenge}

\begin{textebox}[The end-of-term announcement]
It is the end of the term at GBHS Bamenda. The English teacher enters the classroom, closes her register and announces: “Next week, you will have a test on everything we studied about jobs, money and the economy.” Amina's heart beats faster. She remembers the lessons on occupations, salaries and banks. She can name a carpenter, a nurse and a bank clerk. She knows that “to earn a living” means “gagner sa vie”. But can she really use all these words and rules together, in one single test? Like Amina, every student needs a moment to stop, revise, test themselves and fill the gaps. This lesson is exactly that moment: integration, evaluation and remediation on Chapter 2.
\end{textebox}

\textbf{Glossaire utile :}

\begin{center}
\begin{tabularx}{\linewidth}{|l|l|X|}
\hline
\rowcolor{popblueL} English word & Nature & Traduction française \\
\hline
register & nom & cahier d'appel, registre \\
\hline
heart beats faster & expression & le cœur bat plus vite \\
\hline
to fill the gaps & expression & combler les lacunes \\
\hline
remediation & nom & remédiation, rattrapage \\
\hline
\end{tabularx}
\end{center}

\textbf{Problématique de la leçon.} How can we bring together everything we learned in Chapter 2 — jobs, money, banking, grammar and skills — and check that we are really ready for the test? La réponse tient en trois verbes que nous allons vivre aujourd'hui : \cle{revise} (réviser), \cle{evaluate} (évaluer) et \cle{remediate} (remédier). Tout commence par un retour collectif sur nos acquis : la carte mentale.

\courssub{Collective brainstorm : “What did we learn in Chapter 2?”}

Avant d'ouvrir le cahier, fermons-le ! Le professeur écrit au tableau la question \eng{What did we learn in Chapter 2?} et chaque élève cite, de mémoire, un mot, une expression ou une règle du chapitre. C'est la technique du \cle{brainstorming} : on dit tout ce qui vient, sans craindre l'erreur, puis on classe ensuite.

Voici le type de réponses attendues dans une classe de 1ère A :

\begin{itemize}
\item \eng{a farmer, a tailor, a bank clerk, a nurse, a teacher} — des métiers ;
\item \eng{salary, wages, income, savings, a loan, an account} — l'argent et la banque ;
\item \eng{to earn a living, to save money, to borrow money} — des expressions clés ;
\item \eng{the present simple for habits, comparatives, countable and uncountable nouns} — la grammaire ;
\item \eng{reading a job advertisement, writing a letter of application} — les compétences.
\end{itemize}

\begin{experience}[Class activity : the four-branch mind-map]
Le professeur dessine au tableau un nœud central : \eng{Chapter 2: Economic Life and Occupations}. Quatre branches partent du centre : \textbf{Jobs \& Occupations}, \textbf{Money \& Banking}, \textbf{Grammar points}, \textbf{Skills}. Chaque élève, à tour de rôle, cite un mot ou une règle et dit dans quelle branche le placer, en justifiant : \eng{“A loan goes in Money and Banking because you get it from a bank.”} La classe valide ou corrige. Ensuite, chacun recopie la carte mentale complétée dans son cahier : ce sera sa fiche de révision personnelle.
\end{experience}

\begin{popfigure}
\begin{center}\tcbox[colback=popblue,colframe=popblue,coltext=white,arc=9pt,boxsep=4pt,left=6pt,right=6pt]{\bfseries\small Chapter 2 Economic Life and Occupations}\end{center}
\vspace{-2pt}
\begin{tcbraster}[raster columns=2,raster equal height=rows,raster column skip=6pt,raster row skip=6pt,enhanced,blankest]
\begin{tcolorbox}[enhanced,colback=white,colframe=popgreen,boxrule=0.9pt,arc=3pt,title={Vocabulary},fonttitle=\bfseries\scriptsize,coltitle=white,colbacktitle=popgreen,left=4pt,right=4pt,top=2pt,bottom=2pt,boxsep=1pt,toptitle=1.5pt,bottomtitle=1.5pt,halign=left]
\begin{itemize}[nosep,leftmargin=*,itemsep=1pt,font=\scriptsize]
\item jobs \& occupations
\item money \& banking
\end{itemize}
\end{tcolorbox}
\begin{tcolorbox}[enhanced,colback=white,colframe=poporange,boxrule=0.9pt,arc=3pt,title={Grammar},fonttitle=\bfseries\scriptsize,coltitle=white,colbacktitle=poporange,left=4pt,right=4pt,top=2pt,bottom=2pt,boxsep=1pt,toptitle=1.5pt,bottomtitle=1.5pt,halign=left]
\begin{itemize}[nosep,leftmargin=*,itemsep=1pt,font=\scriptsize]
\item present simple
\item comparatives countable / uncountable
\end{itemize}
\end{tcolorbox}
\begin{tcolorbox}[enhanced,colback=white,colframe=poppurple,boxrule=0.9pt,arc=3pt,title={Skills},fonttitle=\bfseries\scriptsize,coltitle=white,colbacktitle=poppurple,left=4pt,right=4pt,top=2pt,bottom=2pt,boxsep=1pt,toptitle=1.5pt,bottomtitle=1.5pt,halign=left]
\begin{itemize}[nosep,leftmargin=*,itemsep=1pt,font=\scriptsize]
\item reading a job advert
\item writing a letter
\end{itemize}
\end{tcolorbox}
\begin{tcolorbox}[enhanced,colback=white,colframe=poppink,boxrule=0.9pt,arc=3pt,title={Culture},fonttitle=\bfseries\scriptsize,coltitle=white,colbacktitle=poppink,left=4pt,right=4pt,top=2pt,bottom=2pt,boxsep=1pt,toptitle=1.5pt,bottomtitle=1.5pt,halign=left]
\begin{itemize}[nosep,leftmargin=*,itemsep=1pt,font=\scriptsize]
\item earning a living
\item work in Cameroon
\end{itemize}
\end{tcolorbox}
\end{tcbraster}

\legende{Carte mentale de révision du chapitre 2 : un nœud central et quatre branches (Vocabulary, Grammar, Skills, Culture) à compléter en classe.}
\end{popfigure}

\courssub{How to revise efficiently : la méthode en trois étapes}

Réviser ne veut pas dire relire passivement son cahier pendant des heures. La recherche en didactique le confirme : on retient ce qu'on a \emph{testé}, pas ce qu'on a simplement relu.

\begin{methode}[The three-step revision method]
\begin{enumerate}
\item \textbf{List what you know.} Fermez le cahier et écrivez, de mémoire, tout ce que vous savez sur le chapitre (mots, règles, expressions). C'est exactement ce que fait la carte mentale : elle rend visible ce qui est déjà en mémoire.
\item \textbf{Test yourself.} Répondez à un quiz, faites un exercice de type examen, chronométré et sans aide. Notez précisément ce qui échoue : un mot oublié, une terminaison en \eng{-s} omise, une préposition fausse.
\item \textbf{Rework only what fails.} Ne relisez pas tout le chapitre : retravaillez \emph{uniquement} les points qui ont échoué au test, puis retestez-les. C'est le principe de la remédiation ciblée.
\end{enumerate}
\end{methode}

\begin{aretenir}[Objectifs du chapitre 2 et déroulé de la leçon]
À la fin du chapitre 2, vous devez savoir : nommer les métiers et les activités économiques, parler d'argent et de banque, utiliser les points de grammaire du chapitre (present simple, comparatifs, noms dénombrables et indénombrables) et produire des écrits simples (lettre, paragraphe descriptif). Aujourd'hui : \textbf{révision} (mind-map et synthèses) $\rightarrow$ \textbf{évaluation} (exercices de type examen) $\rightarrow$ \textbf{remédiation} (retravail des points faibles).
\end{aretenir}

\courssub{First diagnostic : job or work?}

Le premier piège que la carte mentale fait remonter est presque toujours le même : la confusion entre \eng{job} et \eng{work}. Observons d'abord deux phrases :

\eng{My uncle has a new job at the bank in Douala.} « Mon oncle a un nouveau poste à la banque de Douala. »

\eng{My uncle has a lot of work this week.} « Mon oncle a beaucoup de travail cette semaine. »

\begin{propriete}[Job (countable) vs work (uncountable)]
\begin{itemize}
\item \cle{job} est un nom \textbf{dénombrable} : il désigne un poste précis, un emploi. On peut dire \eng{a job}, \eng{two jobs}, \eng{many jobs}. Formes : \eng{a job / jobs}.
\item \cle{work} est un nom \textbf{indénombrable} : il désigne le travail en général, l'activité. On ne dit jamais \eng{a work} ni \eng{works} au sens de « travail ». Pour quantifier : \eng{some work}, \eng{a lot of work}, \eng{a piece of work}.
\item \eng{work} existe aussi comme \textbf{verbe} : \eng{She works in a hospital.} « Elle travaille dans un hôpital. »
\end{itemize}
\end{propriete}

\begin{center}
\begin{tabularx}{\linewidth}{|X|X|X|}
\hline
\rowcolor{popblueL} Français & English & Remarque \\
\hline
un emploi, un poste & a job & dénombrable : \eng{two jobs} \\
\hline
du travail (en général) & work & indénombrable : jamais \eng{a work} \\
\hline
chercher du travail & to look for a job / work & les deux sont possibles \\
\hline
travailler (verbe) & to work & \eng{He works hard.} \\
\hline
un travail (une tâche) & a piece of work & pour rendre \eng{work} dénombrable \\
\hline
\end{tabularx}
\end{center}

\begin{attention}[Erreur fréquente des francophones]
Ne dites jamais \eng{I have a work} ni \eng{I have many works} pour « j'ai du travail / beaucoup de travail ». Dites \eng{I have work} ou \eng{I have a lot of work}. À l'inverse, « Il a trouvé un emploi » se traduit \eng{He found a job}, et non \eng{He found a work}. Attention aussi au faux ami : \eng{a work of art} signifie « une œuvre d'art », pas « un travail d'art ».
\end{attention}

\begin{savaistu}[“To earn a living” : gagner sa vie]
Dans le Cameroun anglophone — à Bamenda, Buea ou Kumba — on entend partout l'expression \eng{to earn a living} : \eng{She earns a living by selling vegetables at the market.} « Elle gagne sa vie en vendant des légumes au marché. » C'est l'expression clé du chapitre : elle relie le monde des métiers (\eng{occupations}) à celui de l'argent (\eng{earnings}). On dit aussi \eng{to make a living}, avec le même sens. Retenez la structure : \eng{to earn a living by + verbe en -ing}.
\end{savaistu}

\courssub{Mini-quiz : five diagnostic questions}

Pour terminer la mise en route, le professeur interroge oralement la classe avec cinq questions rapides. Chacun note ses réponses sur une ardoise ou une feuille : c'est le premier \cle{diagnostic} de la séance, celui qui dira quelle branche de la carte mentale est solide et laquelle devra être retravaillée.

\begin{exercice}[Oral mini-quiz — write your answers]
\begin{enumerate}
\item Name three jobs you can find in a hospital.
\item Is \eng{money} countable or uncountable? Give the correct sentence: “He earns \ldots money.”
\item Complete: “A nurse \ldots (work) in a clinic.”
\item Translate: « Mon père gagne sa vie comme chauffeur de taxi. »
\item What is the comparative of \eng{good} and of \eng{bad}?
\end{enumerate}
\end{exercice}

\begin{corrige}[Mini-quiz]
\begin{enumerate}
\item \eng{a doctor, a nurse, a surgeon} (aussi : \eng{a midwife, a pharmacist}).
\item \eng{Money} est indénombrable : \eng{He earns a lot of money.} (jamais \eng{many money}).
\item \eng{A nurse works in a clinic.} — troisième personne du singulier : \eng{-s} obligatoire au present simple.
\item \eng{My father earns a living as a taxi driver.}
\item \eng{good} $\rightarrow$ \eng{better} ; \eng{bad} $\rightarrow$ \eng{worse} (comparatifs irréguliers).
\end{enumerate}
\textbf{Auto-évaluation :} 5/5 : excellent, passez à la suite ; 3--4/5 : revoyez la branche faible de la carte mentale ; 0--2/5 : refaites la synthèse de vocabulaire et de grammaire avant l'évaluation.
\end{corrige}

\begin{exoresolu}[Placing words on the mind-map]
Énoncé — Place each item in the correct branch of the mind-map (Jobs \& Occupations / Money \& Banking / Grammar / Skills): \eng{a loan}, \eng{the present simple}, \eng{a tailor}, \eng{writing a letter of application}, \eng{savings}, \eng{a mechanic}.
\tcblower
Solution détaillée — \eng{a loan} (« un prêt ») appartient à \textbf{Money \& Banking}, car on l'obtient dans une banque. \eng{The present simple} est un temps : branche \textbf{Grammar}. \eng{A tailor} (« un tailleur ») et \eng{a mechanic} (« un mécanicien ») sont des métiers : branche \textbf{Jobs \& Occupations}. \eng{Savings} (« l'épargne ») concerne l'argent déposé à la banque : \textbf{Money \& Banking}. Enfin, \eng{writing a letter of application} (« rédiger une lettre de candidature ») est une compétence de production écrite : branche \textbf{Skills}. Justifier chaque choix à voix haute, comme dans l'activité de classe, est la meilleure façon de mémoriser.
\end{exoresolu}

\begin{aretenir}[Ce qu'il faut retenir de cette mise en route]
La carte mentale est votre plan de bataille pour la révision : quatre branches (vocabulaire, grammaire, compétences, culture) qui résument tout le chapitre 2. Réviser efficacement, c'est \textbf{lister} ce qu'on sait, \textbf{tester}, puis \textbf{retravailler} uniquement ce qui échoue. Le mini-quiz a déjà révélé vos points forts et vos lacunes : gardez-les en tête, car les sections suivantes vont les traiter un par un.
\end{aretenir}

\coursec{Vocabulary synthesis : economic life}

Dans la section précédente, nous avons construit ensemble une carte mentale du chapitre 2 : autour du thème central \eng{Economic Life and Occupations}, vous avez replacé les grandes familles de mots rencontrées dans les leçons précédentes. Il est temps maintenant de passer de la carte au contenu : cette section reprend, classe et approfondit tout le \cle{vocabulaire} de la vie économique. L'objectif est double : enrichir votre lexique actif (pour parler et écrire) et sécuriser votre lexique passif (pour comprendre un texte ou un enregistrement le jour de l'évaluation).

\courssub{Observation : un texte pour faire émerger le vocabulaire}

Avant toute liste de mots, observons le vocabulaire \emph{en situation}. Lisez le texte suivant, rédigé dans le style d'un portrait de presse, et soulignez mentalement tous les mots liés à l'argent et au travail.

\begin{textebox}[A trader's life in Douala — reading text]
My uncle is a trader in Douala. He runs a small business near the Mboppi market and employs five people: two shop assistants, a driver, a tailor and an accountant. Every month he pays their wages on time, because he knows that a happy employee works better. His own income is not fixed: some months are good, others are difficult, especially during the rainy season when customers stay at home.

Ten years ago, my uncle was a civil servant in Yaoundé. He earned a regular salary, but he wanted to be his own boss. So he applied for a loan at a microfinance bank, opened a savings account and started his business with two employees only. At the beginning, the interest on the loan was high, and he was afraid of falling into debt. He made a strict budget: every franc was counted, and he saved money whenever he could.

Today, his business is successful. He pays taxes to the state, he has repaid his loan, and he is proud to be an entrepreneur. Last week, a young man came with his CV to apply for a job as a carpenter. My uncle interviewed him and offered him a six-month contract. “In Cameroon,” my uncle says, “a small business can change a whole family's life. But you must work hard, save money and never spend more than you earn.”
\end{textebox}

\textbf{Glossaire (mots difficiles)}\par
\begin{center}
\begin{tabularx}{\linewidth}{|l|l|X|}
\hline
\rowcolor{popblueL} English word & Nature & Traduction française \\
\hline
trader & nom & commerçant(e) \\
\hline
to run a business & expression & gérer / diriger une entreprise \\
\hline
to employ & verbe & employer (du personnel) \\
\hline
wages & nom pluriel & salaire (payé à l'heure ou à la semaine) \\
\hline
income & nom & revenu \\
\hline
loan & nom & prêt, emprunt \\
\hline
interest & nom & intérêts (d'un prêt) \\
\hline
debt & nom & dette \\
\hline
budget & nom & budget \\
\hline
to repay & verbe & rembourser \\
\hline
contract & nom & contrat \\
\hline
\end{tabularx}
\end{center}

\textbf{Questions de compréhension (entraînement type examen)}\par
\begin{enumerate}
\item Where does the uncle work, and how many people does he employ?
\item Why did he leave his job as a civil servant?
\item What did he do to get money to start his business?
\item Find in the text two words that mean “money you receive regularly”.
\item What advice does the uncle give at the end of the text?
\end{enumerate}

\begin{corrige}[Questions de compréhension]
\begin{enumerate}
\item \eng{He works in Douala, near the Mboppi market, and he employs five people.}
\item \eng{He left because he wanted to be his own boss.}
\item \eng{He applied for a loan at a microfinance bank and opened a savings account.}
\item \eng{salary} et \eng{wages} (on accepte aussi \eng{income} dans le sens de « revenu régulier »).
\item \eng{He advises people to work hard, save money and never spend more than they earn.}
\end{enumerate}
\end{corrige}

\courssub{Champ lexical 1 : jobs and occupations}

Le premier champ lexical du chapitre concerne les \cle{métiers}. Observez comment le texte les utilise : \eng{a trader}, \eng{a driver}, \eng{a tailor}, \eng{a civil servant}, \eng{a carpenter}. Remarquez la présence systématique de l'article indéfini \eng{a/an} devant le métier quand on dit ce que quelqu'un fait : \eng{He is a teacher.} (Il est professeur.) — alors que le français dit « il est professeur » sans article.

\begin{definition}[Jobs and occupations — les métiers du chapitre]
Un \cle{métier} (\eng{job} ou \eng{occupation}) désigne l'activité professionnelle d'une personne. \eng{Occupation} est le terme formel, utilisé dans les formulaires administratifs : \eng{Occupation: nurse.} \eng{Job} est le terme courant de la conversation.
\end{definition}

\begin{center}
\begin{tabularx}{\linewidth}{|X|X|X|}
\hline
\rowcolor{popblueL} English & Français & Exemple en contexte \\
\hline
farmer & agriculteur, fermier & \eng{My grandfather is a farmer in Bafoussam.} \\
\hline
trader & commerçant & \eng{The trader sells fabric at the market.} \\
\hline
teacher & enseignant, professeur & \eng{She is a teacher at Government High School.} \\
\hline
nurse & infirmier / infirmière & \eng{The nurse works at the district hospital.} \\
\hline
driver & chauffeur, conducteur & \eng{The taxi driver knows every street in Bamenda.} \\
\hline
tailor & tailleur, couturier & \eng{The tailor made my uniform in two days.} \\
\hline
carpenter & menuisier, charpentier & \eng{The carpenter is repairing our door.} \\
\hline
civil servant & fonctionnaire & \eng{Her father is a civil servant in the ministry.} \\
\hline
entrepreneur & entrepreneur & \eng{Young entrepreneurs create jobs in Buea.} \\
\hline
\end{tabularx}
\end{center}

\begin{definition}[Entrepreneur]
An \cle{entrepreneur} is a person who starts and runs their own business, taking financial risks to make a profit. (Un entrepreneur : une personne qui crée et dirige sa propre entreprise, en prenant des risques financiers pour réaliser un bénéfice.)
\end{definition}

\begin{savaistu}[D'où vient le mot « entrepreneur » ?]
Le mot \eng{entrepreneur} vient du français « entreprendre » ! Les Anglais l'ont emprunté au XVIII\textsuperscript{e} siècle et l'utilisent tel quel, sans le traduire. Prononcez-le à l'anglaise : « on-tre-pre-neur » avec l'accent sur la dernière syllabe. C'est un des rares cas où un mot français est utilisé directement en anglais économique, comme \eng{résumé} (CV) en anglais américain.
\end{savaistu}

\courssub{Champ lexical 2 : money and banking}

Le deuxième champ lexical concerne l'\cle{argent} et la \cle{banque}. Reprenons les phrases du texte : \eng{he applied for a loan}, \eng{opened a savings account}, \eng{the interest on the loan was high}, \eng{he made a strict budget}, \eng{he saved money}.

\begin{center}
\begin{tabularx}{\linewidth}{|X|X|X|}
\hline
\rowcolor{popblueL} English & Français & Exemple en contexte \\
\hline
salary & salaire (mensuel, fixe) & \eng{She earns a good salary as a nurse.} \\
\hline
wages & salaire (à l'heure / à la semaine) & \eng{He pays the workers' wages every Friday.} \\
\hline
income & revenu (toutes sources) & \eng{His income comes from his shop and his farm.} \\
\hline
loan & prêt, emprunt & \eng{He applied for a loan to buy a motorbike.} \\
\hline
savings & économies, épargne & \eng{She keeps her savings in the bank.} \\
\hline
account & compte (bancaire) & \eng{I want to open an account, please.} \\
\hline
interest & intérêts & \eng{The bank charges 10\% interest on the loan.} \\
\hline
debt & dette & \eng{He is in debt because of the hospital bills.} \\
\hline
budget & budget & \eng{Our family budget includes school fees.} \\
\hline
\end{tabularx}
\end{center}

\begin{attention}[Salary ou wages ? Une distinction que le français ne fait pas]
En français, « salaire » couvre tout. En anglais, il faut choisir :
\begin{itemize}
\item \eng{salary} : rémunération \textbf{fixe}, versée \textbf{chaque mois}, typiquement pour les employés qualifiés et les fonctionnaires. \eng{A teacher receives a monthly salary.}
\item \eng{wages} (presque toujours au pluriel) : rémunération calculée \textbf{à l'heure, au jour ou à la semaine}, typiquement pour les ouvriers et employés payés à la tâche. \eng{The carpenter receives his wages every Saturday.}
\item \eng{income} : terme général = tout l'argent qui entre (salaire + commerce + loyers...). \eng{Farmers' income depends on the harvest.}
\end{itemize}
Ne dites donc pas \eng{my father receives wages every month} pour un fonctionnaire : dites \eng{a monthly salary}.
\end{attention}

\begin{exemplebox}[Exemples traduits]
\eng{She earns a good salary.} = Elle gagne un bon salaire.\par
\eng{He applied for a loan.} = Il a demandé un prêt.\par
\eng{They opened a savings account at the credit union.} = Ils ont ouvert un compte d'épargne à la coopérative de crédit.\par
\eng{We must save money for the school fees.} = Nous devons économiser de l'argent pour les frais de scolarité.
\end{exemplebox}

\courssub{Champ lexical 3 : the workplace}

Le troisième champ lexical décrit le \cle{monde du travail} : les personnes, l'embauche, les documents.

\begin{center}
\begin{tabularx}{\linewidth}{|X|X|X|}
\hline
\rowcolor{popblueL} English & Français & Exemple en contexte \\
\hline
employer & employeur & \eng{The employer pays the workers at the end of the month.} \\
\hline
employee & employé(e), salarié(e) & \eng{The company has fifty employees.} \\
\hline
colleague & collègue & \eng{My colleagues and I work in the same office.} \\
\hline
boss & patron, chef & \eng{My boss is strict but fair.} \\
\hline
interview & entretien (d'embauche) & \eng{She has a job interview tomorrow morning.} \\
\hline
application & candidature & \eng{He sent his application letter last week.} \\
\hline
CV (curriculum vitae) & CV & \eng{Bring your CV to the interview.} \\
\hline
contract & contrat & \eng{She signed a two-year contract.} \\
\hline
\end{tabularx}
\end{center}

\begin{propriete}[Les collocations clés du chapitre]
En anglais, certains verbes vont obligatoirement avec certains noms : ce sont des \cle{collocations}. Il faut les apprendre comme des blocs, jamais mot à mot :
\begin{itemize}
\item \eng{to earn a salary / money} = gagner un salaire / de l'argent (jamais \eng{win a salary} !)
\item \eng{to apply for a job / a loan} = postuler à un emploi / demander un prêt (attention à la préposition \eng{for})
\item \eng{to open an account} = ouvrir un compte
\item \eng{to save money} = économiser de l'argent
\item \eng{to run a business} = gérer une entreprise (jamais \eng{make a business} dans ce sens)
\item \eng{to pay taxes} = payer des impôts
\item \eng{to look for a job} = chercher un emploi
\item \eng{to sign a contract} = signer un contrat
\end{itemize}
\end{propriete}

\begin{attention}[L'erreur n°1 des francophones : « a work »]
En français, « un travail » est comptable. En anglais, \eng{work} est \textbf{indénombrable} quand il signifie « travail, emploi » : on ne peut jamais dire \eng{a work} ni \eng{works}.\par
✗ \eng{I am looking for a work.}\par
✓ \eng{I am looking for a job.} (Je cherche un emploi.)\par
✓ \eng{I have a lot of work today.} (J'ai beaucoup de travail aujourd'hui.)\par
Astuce : \eng{job} = comptable (un poste précis) ; \eng{work} = indénombrable (l'activité en général). Comparez : \eng{He found a job at the bank.} / \eng{His work is very interesting.}
\end{attention}

\begin{attention}[Economic ou economical ?]
Deux adjectifs, deux sens différents :
\begin{itemize}
\item \eng{economic} = qui concerne l'\textbf{économie} (la science, la politique) : \eng{the economic life of the country}, \eng{economic growth}.
\item \eng{economical} = qui permet d'\textbf{économiser}, qui est économe : \eng{This motorbike is very economical — it uses little fuel.}
\end{itemize}
Dans le titre de notre chapitre, c'est bien \eng{Economic Life} (la vie économique), pas \eng{economical life}.
\end{attention}

\courssub{La famille du mot « employ » : construire le vocabulaire}

Une stratégie puissante pour multiplier votre vocabulaire consiste à apprendre les \cle{familles de mots} : à partir d'une même racine, l'anglais fabrique des verbes, des noms de personnes, des noms abstraits et des adjectifs grâce à des suffixes. Observons la famille d'\eng{employ}, omniprésente dans notre texte (\eng{employs five people}, \eng{a happy employee}, \eng{the employer}).

\begin{center}
\begin{tabular}{|l|l|l|}
\hline
\rowcolor{popblueL} Catégorie & Mot anglais & Traduction \\
\hline
Verbe & to employ & employer \\
\hline
Nom de personne (celui qui embauche) & employer & employeur \\
\hline
Nom de personne (celui qui est embauché) & employee & employé, salarié \\
\hline
Nom abstrait & employment & emploi (le fait d'avoir un travail) \\
\hline
Adjectif (avec travail) & employed & employé, qui a un emploi \\
\hline
Adjectif (sans travail) & unemployed & au chômage \\
\hline
Nom abstrait (chômage) & unemployment & chômage \\
\hline
\end{tabular}
\end{center}

\begin{propriete}[Les suffixes -er / -ee : qui fait l'action, qui la subit]
\begin{itemize}
\item Le suffixe \cle{-er} (ou \cle{-or}) désigne celui qui \textbf{fait} l'action : \eng{employ} → \eng{employer} (celui qui emploie) ; \eng{teach} → \eng{teacher} ; \eng{drive} → \eng{driver} ; \eng{trade} → \eng{trader}.
\item Le suffixe \cle{-ee} désigne celui qui \textbf{subit} l'action : \eng{employ} → \eng{employee} (celui qui est employé) ; \eng{train} → \eng{trainee} (stagiaire) ; \eng{interview} → \eng{interviewee} (la personne interviewée).
\end{itemize}
Prononciation : \eng{employer} s'accentue sur « ploi » (« èm-ploi-eur »), \eng{employee} s'accentue sur la fin (« èm-ploi-î »).
\end{propriete}

\begin{popfigure}
\begin{tikzpicture}[
  root/.style={rectangle, rounded corners, draw=popblue, fill=popblueL, align=center, font=\bfseries, minimum width=2.6cm, minimum height=0.9cm},
  branch/.style={rectangle, rounded corners, draw=poporange, fill=poporangeL, align=center, minimum width=2.9cm, minimum height=0.8cm, font=\small},
  lbl/.style={font=\scriptsize\itshape, text=popmuted}
]
\node[root, align=center] (r){to employ\\ (employer)};
\node[branch, above left=1.1cm and 2.6cm of r, align=center] (a){employer\\ employeur};
\node[branch, above right=1.1cm and 2.6cm of r, align=center] (b){employee\\ employé};
\node[branch, below left=1.1cm and 2.6cm of r, align=center] (c){employment\\ emploi};
\node[branch, below right=1.1cm and 2.6cm of r, align=center] (d){unemployed / unemployment\\ chômé / chômage};
\node[lbl, above=0.15cm of a] {suffixe -er : celui qui fait};
\node[lbl, above=0.15cm of b] {suffixe -ee : celui qui subit};
\node[lbl, below=0.15cm of c] {suffixe -ment : nom abstrait};
\node[lbl, below=0.15cm of d] {préfixe un- : négation};
\draw[-{Stealth}, thick, popblue] (r) -- (a);
\draw[-{Stealth}, thick, popblue] (r) -- (b);
\draw[-{Stealth}, thick, popblue] (r) -- (c);
\draw[-{Stealth}, thick, popblue] (r) -- (d);
\end{tikzpicture}
\legende{Carte mentale de la famille du mot « employ » : une racine, quatre directions (personne qui agit, personne qui subit, nom abstrait, négation).}
\end{popfigure}

\begin{exemplebox}[La famille en action]
\eng{The employer employs twenty people.} = L'employeur emploie vingt personnes.\par
\eng{Every employee receives a contract.} = Chaque employé reçoit un contrat.\par
\eng{Employment is a major issue for young graduates.} = L'emploi est un enjeu majeur pour les jeunes diplômés.\par
\eng{Unemployment is high in many towns.} = Le chômage est élevé dans beaucoup de villes.\par
\eng{Many unemployed youths become entrepreneurs.} = Beaucoup de jeunes au chômage deviennent entrepreneurs.
\end{exemplebox}

\courssub{Exercice résolu et entraînement}

\begin{exoresolu}[Complétez avec le mot correct de la famille d'« employ » ou du champ lexical]
Énoncé — Complétez chaque phrase avec le mot qui convient : \emph{employer, employee, unemployment, wages, loan, salary, job, account}.
\begin{enumerate}
\item My sister is a nurse; she earns a good monthly \ldots .
\item The market woman asked the bank for a \ldots to buy more goods.
\item The factory \ldots pays his workers every Friday.
\item After the interview, she finally got the \ldots !
\item Many young people leave school and face \ldots .
\item The driver receives his \ldots every Saturday evening.
\item To save money, my father opened a savings \ldots .
\item A good \ldots arrives at work on time and respects the boss.
\end{enumerate}
\tcblower
Solution détaillée :
\begin{enumerate}
\item \eng{salary} — un salaire mensuel fixe (profession qualifiée) ; \eng{wages} serait faux ici car le paiement est mensuel.
\item \eng{loan} — « demander un prêt » = \eng{to apply for / ask for a loan}.
\item \eng{employer} — celui qui paie les travailleurs, celui qui fait l'action d'employer (suffixe \eng{-er}).
\item \eng{job} — \eng{work} est indénombrable et ne peut pas désigner un « poste » précis ; \eng{to get the job} = obtenir le poste.
\item \eng{unemployment} — nom abstrait avec le préfixe négatif \eng{un-} ; « faire face au chômage ».
\item \eng{wages} — paiement hebdomadaire d'un travailleur manuel : \eng{wages} au pluriel.
\item \eng{account} — collocation obligatoire : \eng{to open a (savings) account}.
\item \eng{employee} — celui qui subit l'action d'être employé (suffixe \eng{-ee}).
\end{enumerate}
\end{exoresolu}

\begin{exercice}[À vous de jouer — préparation à l'évaluation]
\begin{enumerate}
\item Traduisez en anglais : « Mon oncle est commerçant à Douala. Il gagne un bon revenu et paie des impôts. »
\item Corrigez l'erreur : \eng{She is looking for a work in a bank.}
\item Choisissez le bon mot : \eng{The government wants to reduce (unemployment / employer) among young people.}
\item Trouvez la collocation correcte : to \ldots a business (gérer une entreprise) ; to \ldots for a job (postuler) ; to \ldots money (économiser).
\item Classez ces mots dans les trois champs lexicaux de la leçon : \emph{nurse, budget, contract, tailor, interest, CV, farmer, debt}.
\end{enumerate}
\end{exercice}

\begin{corrige}[Exercice « À vous de jouer »]
\begin{enumerate}
\item \eng{My uncle is a trader in Douala. He earns a good income and pays taxes.} (article \eng{a} devant le métier ; \eng{income} pour le revenu global ; \eng{to pay taxes}.)
\item \eng{She is looking for a job in a bank.} (\eng{work} est indénombrable : jamais \eng{a work}.)
\item \eng{unemployment} (le gouvernement veut réduire le chômage ; \eng{employer} est une personne).
\item \eng{run} a business ; \eng{apply} for a job ; \eng{save} money.
\item Jobs and occupations : \eng{nurse, tailor, farmer} ; money and banking : \eng{budget, interest, debt} ; the workplace : \eng{contract, CV}.
\end{enumerate}
\end{corrige}

\begin{aretenir}[L'essentiel du vocabulaire à retenir pour l'évaluation]
\begin{itemize}
\item Trois champs lexicaux : les métiers (\eng{farmer, trader, tailor, civil servant, entrepreneur}), l'argent et la banque (\eng{salary, wages, income, loan, savings, account, interest, debt, budget}), le monde du travail (\eng{employer, employee, colleague, interview, application, CV, contract}).
\item \eng{salary} = mensuel et fixe ; \eng{wages} = à l'heure ou à la semaine ; \eng{income} = revenu total.
\item \eng{work} est indénombrable : on cherche \eng{a job}, jamais \eng{a work}.
\item \eng{economic} (de l'économie) ≠ \eng{economical} (économe).
\item Collocations à connaître par cœur : \eng{earn a salary, apply for a job, open an account, save money, run a business, pay taxes}.
\item Famille d'\eng{employ} : \eng{employer} (-er, celui qui fait) / \eng{employee} (-ee, celui qui subit) / \eng{employment} / \eng{unemployed} / \eng{unemployment}.
\end{itemize}
\end{aretenir}

Ce socle lexical étant désormais consolidé, nous pouvons passer à la synthèse grammaticale du chapitre : les points de grammaire qui vous permettront d'employer ce vocabulaire dans des phrases correctes, à l'écrit comme à l'oral.

\coursec{Synthèse grammaticale : points clés}

Dans la section précédente, nous avons revu tout le vocabulaire de la vie économique : les métiers (\eng{occupations}), les banques, les salaires, le commerce. Mais un vocabulaire riche ne suffit pas : pour parler correctement du monde du travail, il faut maîtriser trois outils grammaticaux essentiels du chapitre 2 : le \cle{present simple} face au \cle{present continuous}, les \cle{comparatifs et superlatifs}, et les \cle{quantifieurs}. C'est l'objet de cette synthèse. À l'examen, ces trois points reviennent systématiquement, aussi bien en grammaire qu'en rédaction.

\courssub{Point 1 — Present simple vs present continuous}

\textbf{Observation.} Lisons d'abord ce court texte sur une jeune Camerounaise, puis observons les verbes en gras.

\begin{textebox}[A week in the life of a bank clerk]
Claudine \textbf{works} in a bank in Douala. Every morning, she \textbf{leaves} home at 6.30 and \textbf{takes} a bendskin to Bonanjo. The bank \textbf{opens} at 8 a.m. and \textbf{closes} at 4 p.m. Claudine usually \textbf{serves} customers and \textbf{counts} money. She \textbf{likes} her job because she \textbf{meets} many people every day.

But this week, things are different. The bank \textbf{is training} its young workers, so Claudine \textbf{is attending} a computer course. Right now, she \textbf{is sitting} in front of a screen and she \textbf{is learning} how to use new software. Her colleague Bertrand \textbf{is replacing} her at the counter. “I \textbf{am enjoying} this course,” she says, “but I \textbf{miss} my customers!”
\end{textebox}

\textbf{Glossaire :} \emph{clerk} (n.) employé(e) de bureau ; \emph{bendskin} (n.) moto-taxi (Cameroun) ; \emph{to serve} (v.) servir ; \emph{counter} (n.) guichet ; \emph{to miss} (v.) manquer (à quelqu'un) ; \emph{software} (n., indénombrable) logiciel.

\textbf{Questions de compréhension :}
\begin{enumerate}
\item What time does the bank open and close?
\item Why is Claudine not serving customers this week?
\item Which verbs describe her habits? Which verbs describe this week only?
\end{enumerate}

\textbf{Analyse.} Les verbes du premier paragraphe (\eng{works, leaves, takes, opens, serves}) décrivent des \emph{habitudes} et des \emph{faits permanents} : c'est le \cle{present simple}. Les verbes du deuxième paragraphe (\eng{is training, is attending, is sitting, is learning}) décrivent des actions \emph{en cours} ou \emph{temporaires} : c'est le \cle{present continuous}.

\begin{propriete}[Le present simple : habitudes et vérités générales]
\textbf{Forme :} sujet + verbe (avec \textbf{-s} à la 3\textsuperscript{e} personne du singulier : \eng{he works, she opens}). Négation et question avec \eng{do/does} : \eng{She does not work on Saturdays. Does she work in Douala?}

\textbf{Emplois :}
\begin{itemize}
\item les habitudes et routines : \eng{She works in a bank.} = Elle travaille dans une banque.
\item les vérités générales et faits permanents : \eng{Banks open at 8 a.m.} = Les banques ouvrent à 8 h.
\item les horaires officiels : \eng{The market closes at 6 p.m.}
\end{itemize}

\textbf{Marqueurs typiques :} \eng{always, usually, often, sometimes, never, every day, every week, on Mondays}.
\end{propriete}

\begin{propriete}[Le present continuous : action en cours ou temporaire]
\textbf{Forme :} sujet + \eng{be} (am/is/are) + verbe en \textbf{-ing} : \eng{I am revising, she is attending, they are working}.

\textbf{Emplois :}
\begin{itemize}
\item une action en cours au moment où l'on parle : \eng{I am revising for my test.} = Je révise pour mon contrôle (en ce moment).
\item une situation temporaire : \eng{She is attending a training course this week.} = Elle suit une formation cette semaine (ce n'est pas son activité habituelle).
\end{itemize}

\textbf{Marqueurs typiques :} \eng{now, right now, at the moment, today, this week, this month}.

\textbf{Orthographe du -ing :} \eng{make} $\rightarrow$ \eng{making} (e muet supprimé) ; \eng{sit} $\rightarrow$ \eng{sitting}, \eng{run} $\rightarrow$ \eng{running} (consonne doublée après voyelle courte accentuée).
\end{propriete}

\begin{popfigure}
\begin{center}
\begin{tikzpicture}[scale=1]
% Present simple : points répétés
\draw[-{Stealth}, thick, popdark] (0,0) -- (10.5,0);
\foreach \x in {1,2.5,4,5.5,7,8.5} {\fill[popblue] (\x,0) circle (0.09);}
\node[below, popblue, align=center] at (5,-0.15) {\textbf{Present simple} : actions répétées};
\node[above, popmuted] at (1,0.1) {\small Lun};
\node[above, popmuted] at (4,0.1) {\small Mer};
\node[above, popmuted] at (7,0.1) {\small Ven};
% Present continuous : flèche autour de now
\draw[-{Stealth}, thick, popdark] (0,-2.2) -- (10.5,-2.2);
\draw[very thick, poporange, -{Stealth}] (3.8,-2.2) -- (6.2,-2.2);
\fill[poporange] (5,-2.2) circle (0.11);
\node[above, poporange] at (5,-2.05) {\textbf{now}};
\node[below, poporange, align=center] at (5,-2.35) {\textbf{Present continuous} : action en cours};
\end{tikzpicture}
\end{center}
\legende{Le present simple : des points répétés sur la ligne du temps ; le present continuous : une flèche qui encadre le moment présent.}
\end{popfigure}

\begin{attention}[Les verbes d'état ne prennent pas le continuous]
Les verbes qui expriment un \emph{état} et non une action — \eng{know, want, like, love, hate, need, believe, understand, own, prefer} — s'emploient presque toujours au present simple, même pour parler du moment présent :

\eng{I want a job.} ✓ (jamais \eng{I am wanting} ✗)

\eng{She knows the manager.} ✓ (jamais \eng{She is knowing} ✗)

\eng{Do you need some money?} ✓ (jamais \eng{Are you needing} ✗)

En français, on dit « je veux » sans hésiter ; en anglais, c'est la même logique : pas de forme en -ing pour ces verbes.
\end{attention}

\begin{exemplebox}[Simple ou continuous ? À vous de justifier]
\begin{itemize}
\item \eng{My father sells cocoa in Kumba.} = Mon père vend du cacao à Kumba. (habitude, métier permanent)
\item \eng{This month, he is selling his cocoa to a new company.} = Ce mois-ci, il vend son cacao à une nouvelle entreprise. (situation temporaire)
\item \eng{Prices usually rise in December.} = Les prix augmentent généralement en décembre. (vérité générale, marqueur \eng{usually})
\item \eng{Look! The price of fuel is rising again.} = Regarde ! Le prix du carburant augmente encore. (en ce moment, marqueur \eng{Look!})
\end{itemize}
\end{exemplebox}

\begin{exoresolu}[Choisissez le bon temps]
\textbf{Énoncé.} Mettez les verbes entre parenthèses au present simple ou au present continuous.

1. Aminatou (work) \dots\ in a bakery in Bamenda. 2. Listen! The manager (speak) \dots\ to a customer. 3. Farmers (harvest) \dots\ cocoa every year in October. 4. This week, we (study) \dots\ economic life at school. 5. I (not / understand) \dots\ this question.

\tcblower
\textbf{Solution détaillée.}
\begin{enumerate}
\item \eng{Aminatou \textbf{works} in a bakery in Bamenda.} — métier permanent $\rightarrow$ present simple, -s à la 3\textsuperscript{e} personne.
\item \eng{Listen! The manager \textbf{is speaking} to a customer.} — \eng{Listen!} signale une action en cours $\rightarrow$ present continuous.
\item \eng{Farmers \textbf{harvest} cocoa every year in October.} — \eng{every year} = habitude $\rightarrow$ present simple.
\item \eng{This week, we \textbf{are studying} economic life at school.} — \eng{this week} = période temporaire $\rightarrow$ present continuous.
\item \eng{I \textbf{do not understand} this question.} — \eng{understand} est un verbe d'état $\rightarrow$ present simple, jamais de continuous.
\end{enumerate}
\end{exoresolu}

\courssub{Point 2 — Comparatifs et superlatifs}

\textbf{Observation.} Comparons trois métiers : \eng{A nurse earns less than a doctor. A teacher's job is harder than it looks. Trading is the most common occupation here.} = Une infirmière gagne moins qu'un médecin. Le métier d'enseignant est plus dur qu'il n'y paraît. Le commerce est le métier le plus courant ici. Remarquez : \emph{less than}, \emph{harder than}, \emph{the most common}.

\begin{propriete}[Formation des comparatifs et superlatifs]
\begin{itemize}
\item \textbf{Adjectifs courts} (1 syllabe, parfois 2) : comparatif en \textbf{-er + than}, superlatif en \textbf{the + -est} : \eng{cheap $\rightarrow$ cheaper than $\rightarrow$ the cheapest}. Attention au doublage : \eng{big $\rightarrow$ bigger $\rightarrow$ the biggest} ; \eng{happy $\rightarrow$ happier $\rightarrow$ the happiest} (y $\rightarrow$ i).
\item \textbf{Adjectifs longs} (2 syllabes et plus) : \textbf{more ... than} / \textbf{the most ...} : \eng{expensive $\rightarrow$ more expensive than $\rightarrow$ the most expensive}.
\item \textbf{Irréguliers} (à apprendre par cœur) : \eng{good $\rightarrow$ better $\rightarrow$ the best} ; \eng{bad $\rightarrow$ worse $\rightarrow$ the worst} ; \eng{far $\rightarrow$ further/farther $\rightarrow$ the furthest/farthest}.
\end{itemize}
\end{propriete}

\begin{center}
\begin{tabularx}{\linewidth}{|X|X|X|X|}
\hline
\rowcolor{popblueL} Type & Adjectif & Comparatif & Superlatif \\
\hline
Court & cheap (bon marché) & cheaper than & the cheapest \\
\hline
Court (y $\rightarrow$ i) & busy (occupé) & busier than & the busiest \\
\hline
Long & difficult & more difficult than & the most difficult \\
\hline
Long & common (courant) & more common than & the most common \\
\hline
Irrégulier & good & better than & the best \\
\hline
Irrégulier & bad & worse than & the worst \\
\hline
\end{tabularx}
\end{center}

\begin{propriete}[Autres structures de comparaison]
\begin{itemize}
\item \textbf{Égalité :} \eng{as ... as} : \eng{A driver earns as much as a secretary.} = Un chauffeur gagne autant qu'un secrétaire. Négation : \eng{not as ... as} : \eng{Farming is not as easy as people think.}
\item \textbf{Infériorité :} \eng{less ... than} : \eng{A nurse earns less than a doctor.} = Une infirmière gagne moins qu'un médecin.
\item \textbf{Proportion :} \eng{the + comparatif ..., the + comparatif ...} : \eng{The more you save, the richer you become.} = Plus tu épargnes, plus tu deviens riche. \eng{The harder you work, the more you earn.} = Plus tu travailles dur, plus tu gagnes.
\end{itemize}
\end{propriete}

\begin{exemplebox}[Comparatifs dans la vie économique]
\begin{itemize}
\item \eng{Trading is the most common occupation here.} = Le commerce est le métier le plus courant ici.
\item \eng{Yaoundé is bigger than Buea, but Douala is the biggest economic centre.} = Yaoundé est plus grande que Buea, mais Douala est le plus grand centre économique.
\item \eng{A bank loan is more expensive than family help.} = Un prêt bancaire coûte plus cher qu'une aide familiale.
\item \eng{This year, the harvest is better than last year.} = Cette année, la récolte est meilleure que l'année dernière.
\end{itemize}
\end{exemplebox}

\begin{attention}[Pièges des francophones avec les comparaisons]
\begin{itemize}
\item Ne mélangez jamais les deux systèmes : \eng{more cheaper} ✗ $\rightarrow$ \eng{cheaper} ✓ ; \eng{the most biggest} ✗ $\rightarrow$ \eng{the biggest} ✓.
\item \eng{Good} ne fait jamais \eng{gooder} ✗ : c'est \eng{better / the best}.
\item Après un comparatif, c'est toujours \textbf{than} (et non \emph{that} ni \emph{then}) : \eng{richer than} ✓.
\item Le superlatif exige toujours \textbf{the} : \eng{the best worker} ✓, jamais \eng{a best worker} ✗.
\end{itemize}
\end{attention}

\courssub{Point 3 — Quantifieurs : many, much, a lot of, (a) few, (a) little}

\textbf{Observation.} \eng{Many people work in the informal sector. How much money do they earn? A lot of traders sell food. Few workers have a contract, but a few have insurance. There is little water today, but we still have a little hope.}

\begin{propriete}[Le choix du quantifieur dépend du nom]
\begin{itemize}
\item \eng{many} + nom \textbf{dénombrable pluriel} : \eng{many workers, many customers, many jobs}.
\item \eng{much} + nom \textbf{indénombrable} (surtout en négation et question) : \eng{much money, much time, much information}. \eng{How much does it cost?}
\item \eng{a lot of} + les deux (surtout à l'affirmative) : \eng{a lot of people, a lot of money}.
\item \eng{(a) few} + dénombrable ; \eng{(a) little} + indénombrable. Sans \eng{a}, le sens est négatif (« presque pas ») ; avec \eng{a}, le sens est positif (« quelques / un peu ») : \eng{Few customers came} = peu de clients sont venus (décevant) ; \eng{A few customers came} = quelques clients sont venus (positif).
\end{itemize}
\end{propriete}

\begin{center}
\begin{tabularx}{\linewidth}{|X|X|X|}
\hline
\rowcolor{popblueL} Français & English & Remarque \\
\hline
beaucoup de personnes & many people / a lot of people & people est dénombrable pluriel \\
\hline
beaucoup d'argent & much money / a lot of money & money est indénombrable \\
\hline
combien de métiers ? & how many jobs? & dénombrable \\
\hline
combien de temps ? & how much time? & indénombrable \\
\hline
quelques clients & a few customers & sens positif \\
\hline
un peu d'espoir & a little hope & sens positif \\
\hline
peu d'informations & little information & information : jamais de « s » \\
\hline
\end{tabularx}
\end{center}

\begin{attention}[« Much people » n'existe pas !]
\eng{Much people} ✗ $\rightarrow$ \eng{many people} ✓ : \eng{people} (« les gens ») est un pluriel dénombrable en anglais. De même, les noms indénombrables anglais suivants ne prennent \textbf{jamais de « s »}, contrairement au français : \eng{information} (des informations), \eng{money} (de l'argent), \eng{news} (des nouvelles), \eng{advice} (des conseils), \eng{work} (du travail), \eng{furniture} (des meubles). On dit : \eng{The news is good.} ✓ (jamais \eng{The news are good} ✗). Enfin, « beaucoup de » ne se traduit pas par un mot unique : choisissez \eng{many} ou \eng{much} selon le nom qui suit.
\end{attention}

\begin{exoresolu}[Many, much, a few ou a little ?]
\textbf{Énoncé.} Complétez avec \eng{many}, \eng{much}, \eng{a few} ou \eng{a little}.

1. How \dots\ money does a taxi driver earn per day? 2. There are \dots\ banks in Yaoundé. 3. I need \dots\ information about this job. 4. \dots\ students want to become nurses. 5. We have \dots\ time before the shop closes — let's hurry!

\tcblower
\textbf{Solution détaillée.}
\begin{enumerate}
\item \eng{How \textbf{much} money does a taxi driver earn per day?} — \eng{money} est indénombrable $\rightarrow$ \eng{much} (question).
\item \eng{There are \textbf{many} banks in Yaoundé.} — \eng{banks} est dénombrable pluriel $\rightarrow$ \eng{many}.
\item \eng{I need \textbf{a little} information about this job.} — \eng{information} est indénombrable et le sens est positif (« un peu ») $\rightarrow$ \eng{a little}.
\item \eng{\textbf{Many} students want to become nurses.} — \eng{students} est dénombrable $\rightarrow$ \eng{many}.
\item \eng{We have \textbf{a little} time before the shop closes.} — \eng{time} indénombrable, sens positif (« il nous reste un peu de temps ») $\rightarrow$ \eng{a little}.
\end{enumerate}
\end{exoresolu}

\begin{aretenir}[L'essentiel de la synthèse grammaticale]
\begin{itemize}
\item Habitude ou vérité générale $\rightarrow$ \textbf{present simple} (\eng{She works in a bank}) ; action en cours ou temporaire $\rightarrow$ \textbf{present continuous} (\eng{She is attending a course this week}). Verbes d'état : jamais de -ing.
\item Comparatif : adjectif court \textbf{-er than}, adjectif long \textbf{more ... than} ; superlatif : \textbf{the -est} / \textbf{the most} ; irréguliers : \eng{good $\rightarrow$ better $\rightarrow$ the best}, \eng{bad $\rightarrow$ worse $\rightarrow$ the worst}.
\item Quantifieurs : \eng{many} + dénombrable, \eng{much} + indénombrable, \eng{a lot of} + les deux ; \eng{information, money, news} : jamais de « s ».
\end{itemize}
\end{aretenir}

Ces trois points de grammaire étant désormais consolidés, nous allons les mettre à l'épreuve dans la section suivante, à travers des exercices de compréhension écrite et orale de type examen.

\coursec{Reading \& listening practice (exam-type)}

Après avoir synthétisé les points de grammaire clés du chapitre (le présent simple et le présent continu, le prétérit, les comparatifs, l'expression du but avec \eng{to} et \eng{in order to}), il est temps de vérifier que vous savez les \emph{mobiliser} dans des épreuves de type examen. Aux évaluations de fin de séquence comme au baccalauréat, la compréhension de l'écrit (\cle{reading}) et de l'oral (\cle{listening}) représente une part essentielle de la note. Cette section vous entraîne sur un texte et un dialogue inédits, dans les conditions réelles : questions variées, justifications exigées, gestion du temps.

\courssub{Reading practice : « From Hawker to Shop Owner »}

\textbf{Étape 1 : découvrir le texte.}\par

Lisez le texte suivant une première fois, sans dictionnaire, pour en dégager l'idée générale. Il s'agit d'un article de presse original, rédigé dans le style des journaux anglophones camerounais.

\begin{textebox}[From Hawker to Shop Owner — The Buea Weekly Herald]
BUEA — Five years ago, Manka Eposi was a hawker at the Buea central market. Every morning, she carried a heavy basket of vegetables on her head and walked from street to street, calling out to customers. The work was tiring, and her daily profit was small — sometimes less than two thousand francs.

But Manka had a plan. Every evening, before going home, she saved a small part of her profit in a wooden box under her bed. “My friends laughed at me,” she remembers. “They said the money was too little. But I knew that small drops of water make a river.”

After three years of discipline, Manka had saved enough money to rent a small stall inside the market. Her business grew quickly because she was honest and her prices were fair. Customers trusted her, and many became regular clients.

Last year, Manka took a bigger step. She went to a microfinance bank in Molyko and asked for a loan. The bank manager studied her savings record and agreed to lend her five hundred thousand francs. With this money, she opened her own shop, where she now sells food products, soap and cooking oil.

Today, Manka employs two assistants and pays their salaries every month. “I am proud of my journey,” she says with a smile. “My advice to young people is simple: start small, save regularly, and never be ashamed of honest work.”

Manka's story shows that economic success does not always require a rich family or a big diploma. Sometimes, it only requires patience, hard work and good management.
\end{textebox}

\textbf{Glossaire.}\par

\begin{center}
\begin{tabularx}{\linewidth}{|l|l|X|}
\hline
\rowcolor{popblueL} English & Français & Exemple tiré du texte \\
\hline
\cle{hawker} (n.) & vendeur / vendeuse ambulant(e) & \eng{Manka was a hawker at the central market.} \\
\hline
\cle{stall} (n.) & étal, stand (au marché) & \eng{She rented a small stall inside the market.} \\
\hline
\cle{loan} (n.) & prêt (bancaire) & \eng{She asked the bank for a loan.} \\
\hline
\cle{savings} (n. plur.) & économies, épargne & \eng{The manager studied her savings record.} \\
\hline
\cle{profit} (n.) & bénéfice, gain & \eng{Her daily profit was small.} \\
\hline
\cle{customer} (n.) & client(e) & \eng{Customers trusted her.} \\
\hline
\cle{to employ} (v.) & employer, embaucher & \eng{Manka employs two assistants.} \\
\hline
\cle{salary} (n.) & salaire & \eng{She pays their salaries every month.} \\
\hline
\end{tabularx}
\end{center}

\begin{attention}[Faux amis et pièges lexicaux]
\begin{itemize}
\item \eng{profit} se prononce « pro-fit » et signifie \emph{bénéfice} ; ne pas confondre avec le verbe français « profiter (de) », qui se dit \eng{to benefit from} ou \eng{to take advantage of}.
\item \eng{savings} est presque toujours au pluriel : \eng{her savings} = « ses économies ». On ne dit pas \emph{a saving} pour « une économie ».
\item \eng{to employ} = « employer » (donner du travail), mais \eng{an employee} = « un employé » et \eng{an employer} = « un employeur ». Attention à ne pas inverser les suffixes \cle{-er} (celui qui donne le travail) et \cle{-ee} (celui qui le reçoit).
\item \eng{to trust} = « faire confiance à », sans préposition : \eng{Customers trusted her.} On ne dit pas \emph{trusted on her}.
\end{itemize}
\end{attention}

\courssub{Méthode de lecture en situation d'examen}

Beaucoup d'élèves perdent des points non pas parce qu'ils ne comprennent pas le texte, mais parce qu'ils le lisent \emph{mal}. Voici la démarche professionnelle, en cinq étapes.

\begin{methode}[Comment aborder un texte de compréhension écrite]
\begin{enumerate}
\item \textbf{Lire les questions d'abord.} Avant même de lire le texte, lisez toutes les questions. Elles vous indiquent ce qu'il faut chercher : un chiffre, une date, une opinion, une cause.
\item \textbf{Repérer les mots-clés.} Soulignez dans chaque question deux ou trois mots-clés (noms propres, verbes d'action, chiffres). Exemple : dans \eng{Why did the bank agree to lend Manka money?}, les mots-clés sont \eng{bank}, \eng{agree}, \eng{lend}.
\item \textbf{Première lecture rapide} (skimming) : lisez le texte une fois pour saisir le thème général, le lieu, les personnages. Ne bloquez pas sur un mot inconnu.
\item \textbf{Deuxième lecture attentive} (scanning) : relisez en cherchant les mots-clés des questions. Soulignez dans le texte la phrase qui contient chaque réponse.
\item \textbf{Rédiger les réponses justifiées.} Pour chaque réponse, citez la ligne ou la phrase du texte qui la prouve. Relisez-vous : grammaire, orthographe, majuscules.
\end{enumerate}
\end{methode}

\begin{popfigure}
\begin{center}
\begin{tikzpicture}[
node distance=0.55cm,
box/.style={rectangle, rounded corners=3pt, draw=popblue, fill=popblueL, text width=3.6cm, align=center, font=\small, minimum height=0.9cm},
arr/.style={-{Stealth[length=2.5mm]}, thick, poporange}
]
\node[box] (q) {\textbf{1. Lire les questions} avant le texte};
\node[box, below=of q] (k) {\textbf{2. Repérer les mots-clés} de chaque question};
\node[box, below=of k] (r1) {\textbf{3. Première lecture} : idée générale};
\node[box, below=of r1] (r2) {\textbf{4. Deuxième lecture} : souligner les justifications};
\node[box, below=of r2, fill=popgreenL, draw=popgreen] (a) {\textbf{5. Rédiger} des réponses justifiées};
\draw[arr] (q) -- (k);
\draw[arr] (k) -- (r1);
\draw[arr] (r1) -- (r2);
\draw[arr] (r2) -- (a);
\end{tikzpicture}
\end{center}
\legende{Organigramme de la méthode de lecture en cinq étapes}
\end{popfigure}

\courssub{Questions de compréhension (type examen)}

\begin{exercice}[Comprehension check — 10 points]
Répondez en anglais, en vous appuyant sur le texte « From Hawker to Shop Owner ».

\textbf{A. True or False? Justify with a line from the text.} (3 pts)
\begin{enumerate}
\item Manka's friends encouraged her to save money.
\item The bank refused to give Manka a loan.
\item Manka now gives work to other people.
\end{enumerate}

\textbf{B. Choose the correct answer.} (3 pts)
\begin{enumerate}
\item At the beginning, Manka sold: a) food products in a shop; b) vegetables in the streets; c) soap at a stall.
\item She kept her savings: a) in a bank; b) at the market; c) in a wooden box.
\item The bank lent her: a) two thousand francs; b) five hundred thousand francs; c) three million francs.
\end{enumerate}

\textbf{C. Answer in complete sentences.} (3 pts)
\begin{enumerate}
\item Why did Manka's business grow quickly?
\item What does Manka sell in her shop today?
\item According to the last paragraph, what three qualities lead to success?
\end{enumerate}

\textbf{D. Vocabulary in context.} (1 pt)
What does the expression \eng{“small drops of water make a river”} mean in the text? Explain in your own words.
\end{exercice}

\begin{corrige}[Comprehension check]
\textbf{A. True or False?}
\begin{enumerate}
\item \eng{False. Her friends did not encourage her: “My friends laughed at me… They said the money was too little.”}
\item \eng{False. The bank agreed to lend her the money: “The bank manager studied her savings record and agreed to lend her five hundred thousand francs.”}
\item \eng{True. She employs two assistants: “Today, Manka employs two assistants and pays their salaries every month.”}
\end{enumerate}
\textbf{B. QCM.} 1. b) ; 2. c) ; 3. b).
\par\textbf{C. Questions ouvertes.}
\begin{enumerate}
\item \eng{Her business grew quickly because she was honest and her prices were fair.}
\item \eng{Today she sells food products, soap and cooking oil.}
\item \eng{According to the text, success requires patience, hard work and good management.}
\end{enumerate}
\textbf{D.} \eng{It means that small amounts of money saved regularly can become a large sum over time.} (« De petites sommes épargnées régulièrement finissent par former un grand capital. »)
\end{corrige}

\courssub{Exemple traité en classe : corriger un vrai/faux ensemble}

\begin{exoresolu}[True or False — correction collective]
\textbf{Énoncé.} \eng{Manka's friends encouraged her to save money.} True or False? Justify your answer with a quotation from the text.
\tcblower
\textbf{Étape 1 — Repérage des mots-clés.} Les mots-clés sont \eng{friends} et \eng{encouraged}. On cherche donc le passage du texte qui parle des amis de Manka.

\textbf{Étape 2 — Localisation dans le texte.} Au deuxième paragraphe, on lit : \eng{“My friends laughed at me,” she remembers. “They said the money was too little.”}

\textbf{Étape 3 — Comparaison.} « Se moquer d'elle » (\eng{laughed at me}) est le contraire de « l'encourager ». La phrase est donc \textbf{fausse}.

\textbf{Étape 4 — Rédaction de la réponse type examen.}

\eng{False. Manka's friends did not encourage her. Paragraph 2: “My friends laughed at me… They said the money was too little.”}

(« Faux. Les amis de Manka ne l'ont pas encouragée. Paragraphe 2 : “Mes amis se moquaient de moi… Ils disaient que l'argent était trop peu.” »)

\textbf{Remarque de méthode.} La réponse comporte trois éléments : le verdict (\eng{False}), une reformulation, et la citation exacte avec sa localisation. C'est ce format complet qui rapporte le maximum de points.
\end{exoresolu}

\begin{methode}[Répondre à un vrai/faux : la règle d'or]
\begin{enumerate}
\item Écrivez d'abord \eng{True} ou \eng{False}, clairement.
\item Reformulez la bonne information en une phrase simple (ne recopiez pas tout le paragraphe).
\item Citez \textbf{toujours} le passage du texte qui justifie votre réponse, entre guillemets anglais : \eng{Paragraph 2: “she saved a small part of her profit”}.
\item Une réponse \textbf{sans justification perd des points}, même si le verdict est correct : la justification vaut souvent la moitié de la note de la question.
\end{enumerate}
\end{methode}

\begin{attention}[QCM : la stratégie de l'élimination]
Au QCM, ne choisissez jamais au hasard. \textbf{Éliminez d'abord les réponses impossibles}, puis comparez les restantes avec le texte. Exemple : pour la question sur le montant du prêt, le texte dit \eng{five hundred thousand francs}. Éliminez a) (\eng{two thousand} = le profit quotidien, piège !) et c) (chiffre absent du texte). Il reste b). Les distracteurs d'un QCM sont presque toujours des éléments \emph{présents dans le texte mais déplacés} : méfiez-vous des chiffres et des noms que vous reconnaissez trop facilement.
\end{attention}

\begin{savaistu}[Pourquoi les correcteurs adorent les justifications]
Aux épreuves d'anglais des examens camerounais, les consignes de correction valorisent explicitement la justification citée du texte. Une réponse comme \eng{True — “she saved a small part of her profit”} montre que vous n'avez pas deviné : vous avez \emph{prouvé}. C'est la même logique dans les examens internationaux (IELTS, TOEFL) : la compétence évaluée n'est pas seulement la compréhension, mais la capacité à \cle{localiser une information précise} dans un document — une compétence utile dans toute la vie professionnelle.
\end{savaistu}

\courssub{Listening practice : « At the job interview »}

\textbf{Consigne de classe.} L'enseignant lit le script deux fois (ou le fait lire par deux élèves). Première écoute : notez l'idée générale et les personnages. Deuxième écoute : prenez des notes sur les détails (expérience, salaire, disponibilité). Ensuite, répondez aux questions sans regarder le script.

\begin{textebox}[At the job interview — listening script]
Manager: Good morning, Mr Tabi. Please sit down. Thank you for coming.

Candidate: Good morning, sir. Thank you for the invitation.

Manager: So, you are applying for the position of shop assistant. Can you tell me about your work experience?

Candidate: Yes, of course. I worked for two years in a supermarket in Douala. I served customers, arranged products on the shelves and managed the cash register.

Manager: I see. And why did you leave that job?

Candidate: My family moved to Buea, so I had to leave. But I learnt a lot there, especially how to deal with difficult customers politely.

Manager: Good. This job also requires basic computer skills. Can you use a computer?

Candidate: Yes, I can. I have a certificate in computer studies, and I can use spreadsheets to record sales and stock.

Manager: Excellent. Now, let's talk about salary. What salary do you expect?

Candidate: In my previous job, I earned eighty thousand francs a month. Given my experience, I would expect around ninety thousand francs.

Manager: That is reasonable. Our starting salary is eighty-five thousand francs, with an increase after six months. Is that acceptable?

Candidate: Yes, sir, it is.

Manager: Fine. When can you start?

Candidate: I can start immediately, sir.

Manager: Very good, Mr Tabi. We will call you on Friday to confirm. Have a nice day.

Candidate: Thank you very much, sir. Goodbye.
\end{textebox}

\textbf{Glossaire de l'écoute.}\par

\begin{center}
\begin{tabularx}{\linewidth}{|l|l|X|}
\hline
\rowcolor{popgreenL} English & Français & Remarque \\
\hline
\eng{to apply for a position} & postuler à un emploi & \eng{apply for} + le poste ; \eng{apply to} + l'entreprise \\
\hline
\eng{shop assistant} & vendeur / vendeuse (en magasin) & se prononce « chope a-ssis-tent' » \\
\hline
\eng{cash register} & caisse enregistreuse & \\
\hline
\eng{to deal with customers} & s'occuper des clients & \eng{deal — dealt — dealt} (irrégulier) \\
\hline
\eng{salary} & salaire (mensuel) & ne pas confondre avec \eng{wages} (salaire à l'heure ou à la semaine) \\
\hline
\eng{to expect} & s'attendre à, espérer & faux ami : « attendre quelqu'un » = \eng{to wait for somebody} \\
\hline
\eng{starting salary} & salaire de départ & \\
\hline
\end{tabularx}
\end{center}

\begin{exercice}[Listening comprehension — 5 points]
Écoutez le dialogue deux fois, puis répondez en anglais.
\begin{enumerate}
\item What position is Mr Tabi applying for? (1 pt)
\item How long did he work in the supermarket in Douala? (1 pt)
\item Mention two tasks he did in his previous job. (1 pt)
\item True or False? Justify: \eng{The manager offers exactly the salary the candidate expected.} (1 pt)
\item When will the company call the candidate? (1 pt)
\end{enumerate}
\end{exercice}

\begin{corrige}[Listening comprehension]
\begin{enumerate}
\item \eng{He is applying for the position of shop assistant.} (« Il postule au poste de vendeur. »)
\item \eng{He worked there for two years.} (« Il y a travaillé pendant deux ans. »)
\item \eng{He served customers and managed the cash register.} (Autre réponse possible : \eng{arranged products on the shelves}.)
\item \eng{False. The candidate expected around ninety thousand francs, but the manager offers eighty-five thousand: “Our starting salary is eighty-five thousand francs.”} Le piège : les deux chiffres figurent dans le dialogue ; il faut les comparer.
\item \eng{They will call him on Friday.} (« Ils l'appelleront vendredi. »)
\end{enumerate}
\end{corrige}

\begin{attention}[Écoute : les pièges classiques des francophones]
\begin{itemize}
\item \textbf{Les chiffres.} Distinguez bien \eng{eighty-five} (85) et \eng{ninety} (90) ; à l'oral, la fin du mot est décisive. Entraînez-vous à noter les chiffres en chiffres arabes, pas en lettres.
\item \textbf{Les faux amis.} \eng{to expect} ne veut pas dire « attendre » : \eng{What salary do you expect?} = « Quel salaire espérez-vous ? ».
\item \textbf{Le rythme.} Ne vous arrêtez jamais sur un mot incompris : continuez à écouter, le sens global suffit souvent à répondre. La deuxième écoute sert à vérifier, pas à découvrir.
\end{itemize}
\end{attention}

\begin{experience}[Auto-évaluation rapide]
Après avoir terminé les deux épreuves, remplissez cette grille à deux colonnes. Dans la colonne de gauche, notez votre score à chaque partie (vrai/faux, QCM, questions ouvertes, vocabulaire, écoute). Dans la colonne de droite, écrivez la cause de chaque erreur : mot inconnu, justification oubliée, mauvaise élimination au QCM, chiffre mal entendu. Cette analyse prépare la remédiation de la section 6 : on ne progresse que si l'on sait \emph{pourquoi} l'on s'est trompé.
\end{experience}

Vous disposez désormais d'une méthode complète pour la lecture et l'écoute en conditions d'examen. La section suivante vous permettra de transférer ces mêmes acquis à l'expression orale et écrite, avant l'évaluation finale du chapitre.

\coursec{Speaking \& writing consolidation}

Après avoir vérifié vos compétences de compréhension (reading et listening) à travers des exercices de type examen, il est temps de passer à la \cle{production} : c'est-à-dire de mobiliser tout le vocabulaire et toute la grammaire du chapitre 2 pour \textbf{parler} et \textbf{écrire}. Dans cette section, vous allez d'abord observer un dialogue modèle sur les métiers, puis apprendre à rédiger une courte lettre de candidature (\eng{application letter}), compétence très demandée à l'examen.

\courssub{Speaking : talking about jobs}

\textbf{Observation.} Lisez d'abord ce dialogue entre deux élèves de Première, Amina et Peter, qui comparent les métiers de leurs parents. Repérez les comparatifs et le vocabulaire des occupations.

\begin{textebox}[Model dialogue — Talking about jobs]
\textbf{Amina:} Hi Peter! What does your father do for a living?

\textbf{Peter:} He is a taxi driver in Douala. He works harder than a civil servant, but he earns less. And what about your mother?

\textbf{Amina:} My mother is a nurse at the Buea Regional Hospital. Her job is more tiring than my father's, I think. She works at night sometimes.

\textbf{Peter:} Really? A nurse earns more money than a taxi driver, doesn't she?

\textbf{Amina:} Yes, she does, but her work is more stressful. What would you like to become later?

\textbf{Peter:} I'd like to become an engineer, because engineering is more interesting than driving. And you?

\textbf{Amina:} I want to be an entrepreneur. I think it is the most exciting job of all!
\end{textebox}

\textbf{Glossaire.} \emph{taxi driver} (n.) : chauffeur de taxi ; \emph{civil servant} (n.) : fonctionnaire ; \emph{to earn} (v.) : gagner (de l'argent) ; \emph{nurse} (n.) : infirmier / infirmière ; \emph{tiring} (adj.) : fatigant ; \emph{stressful} (adj.) : stressant ; \emph{entrepreneur} (n.) : entrepreneur.

\begin{definition}[Formules fonctionnelles pour parler des métiers]
Voici les expressions-clés à mémoriser pour toute conversation sur les occupations :
\begin{itemize}
\item \eng{“What do you do for a living?”} « Que fais-tu dans la vie ? » (question sur le métier)
\item \eng{“What does your father do?”} « Que fait ton père (comme métier) ? »
\item \eng{“I'd like to become an engineer.”} « Je voudrais devenir ingénieur. »
\item \eng{“This job is more interesting than that one.”} « Ce métier est plus intéressant que celui-là. »
\item \eng{“Is your job well paid?”} « Est-ce que ton métier est bien payé ? » (préférable à \eng{How much do you earn?})
\item \eng{“She works harder than her brother.”} « Elle travaille plus dur que son frère. »
\end{itemize}
\end{definition}

\begin{attention}[Faux pas culturel et linguistique]
Dans les pays anglophones, demander \eng{“How much do you earn?”} à une personne que l'on connaît peu est considéré comme \textbf{indiscret}. Préférez : \eng{“Is your job well paid?”} « Est-ce que ton métier est bien payé ? ». Attention aussi : \eng{salary} signifie « salaire », pas « salut » !
\end{attention}

\begin{propriete}[Rappel : les comparatifs dans le dialogue]
Le dialogue modèle illustre les deux formes du comparatif étudiées dans ce chapitre :
\begin{center}
\begin{tabularx}{\linewidth}{|X|X|X|}
\hline
\rowcolor{popblueL} Type d'adjectif & Règle & Exemple \\
\hline
Court (1 syllabe) & adjectif + \textbf{-er} + than & \eng{harder than}, \eng{faster than} \\
\hline
Long (2+ syllabes) & \textbf{more} + adjectif + than & \eng{more tiring than}, \eng{more interesting than} \\
\hline
Irrégulier & forme spéciale & \eng{good $\rightarrow$ better}, \eng{bad $\rightarrow$ worse}, \eng{little $\rightarrow$ less} \\
\hline
Superlatif & the + -est / the most & \eng{the most exciting job of all} \\
\hline
\end{tabularx}
\end{center}
\end{propriete}

\courssub{Pronunciation and intonation}

\begin{propriete}[Prononciation des mots-clés du chapitre]
Prononcez ces mots en suivant la transcription en lettres françaises (accent tonique en \textbf{gras}) :
\begin{itemize}
\item \eng{salary} : « \textbf{sa}-la-ri » (accent sur la première syllabe)
\item \eng{entrepreneur} : « on-truh-pruh-\textbf{neur} » (accent sur la dernière syllabe)
\item \eng{interview} : « \textbf{in}-teu-vyou » (accent sur la première syllabe)
\item \eng{economy} : « i-\textbf{ko}-no-mi » — attention, on accentue la \textbf{deuxième} syllabe : \eng{EConomy} est incorrect ; dites « i-KO-no-mi ».
\item \eng{job} : « djobe » ; \eng{worker} : « weu-keur »
\end{itemize}
\end{propriete}

\begin{propriete}[L'intonation des questions]
En anglais, la mélodie de la phrase change le sens :
\begin{itemize}
\item \textbf{Voix qui monte ↗} pour les questions fermées (réponse yes/no) : \eng{Do you like your job? ↗} ; \eng{Does she earn more? ↗}
\item \textbf{Voix qui descend ↘} pour les questions en Wh- : \eng{What do you do? ↘} ; \eng{How much does he earn? ↘}
\end{itemize}
Contrairement au français, où la voix monte presque toujours en fin de question, l'anglais fait \emph{descendre} la voix dans les questions en Wh-. Entraînez-vous à voix haute !
\end{propriete}

\begin{methode}[Réussir l'interaction orale en 5 étapes]
\begin{enumerate}
\item \textbf{Saluez} votre partenaire : \eng{“Hello! How are you today?”}
\item \textbf{Parlez lentement} et articulez : la clarté compte plus que la vitesse.
\item \textbf{Utilisez au moins un comparatif} : \eng{“My father's job is harder than my uncle's.”}
\item \textbf{Ajoutez un connecteur} : \eng{“However, he earns less.”} / \eng{“Therefore, I want to study hard.”}
\item \textbf{Concluez poliment} : \eng{“Thank you very much. It was nice talking to you!”}
\end{enumerate}
\end{methode}

\begin{experience}[Role-play : Job interview in pairs]
Travaillez à deux. L'élève A est journaliste pour une radio de Yaoundé ; l'élève B est un commerçant, un chauffeur de bendskin ou une infirmière. A pose au moins cinq questions (\eng{“What do you do? How long have you worked here? Do you like your job? ↗ ...”}). B répond avec au moins deux comparatifs et un connecteur. Puis échangez les rôles. La classe écoute et note les comparatifs entendus.
\end{experience}

\courssub{Writing : a short application letter}

\textbf{Observation.} Voici le modèle d'une lettre de candidature courte, écrite par une élève qui postule pour un emploi de vacances.

\begin{textebox}[Model application letter]
15, Avenue Kennedy\\
Yaoundé, Cameroon\\
15th May 2024

The Manager\\
Mfoundi Supermarket\\
Yaoundé

Dear Sir,

I am writing to apply for the position of shop assistant in your supermarket. I am a seventeen-year-old student at Government High School Yaoundé.

First, I speak English and French fluently. Moreover, I worked in my uncle's shop last year, so I have some experience with customers. However, I am still a student, therefore I can only work during the holidays.

I am a hard-working and honest person. I would be very happy to come for an interview. In conclusion, I believe I am the right person for this job.

Yours faithfully,

Ngo Bell Sandrine
\end{textebox}

\textbf{Glossaire.} \emph{to apply for} (v.) : postuler à ; \emph{shop assistant} (n.) : vendeur / vendeuse ; \emph{fluently} (adv.) : couramment ; \emph{hard-working} (adj.) : travailleur ; \emph{faithfully} (adv.) : fidèlement (formule de politesse).

\begin{propriete}[Structure de la lettre formelle courte]
Une lettre de candidature (\eng{application letter}) respecte toujours cet ordre :
\begin{enumerate}
\item \textbf{Sender's address} (adresse de l'expéditeur) en haut à gauche ou à droite ;
\item \textbf{Date} : \eng{15th May 2024} (format britannique : jour mois année) ;
\item \textbf{Receiver's address} (adresse du destinataire) ;
\item \textbf{Salutation} : \eng{Dear Sir,} / \eng{Dear Madam,} ;
\item \textbf{Body} en trois paragraphes : présentation et motif $\rightarrow$ compétences et expérience $\rightarrow$ conclusion et disponibilité ;
\item \textbf{Closing} : \eng{Yours faithfully,} (si on ne connaît pas le nom) ou \eng{Yours sincerely,} (si on connaît le nom) ;
\item \textbf{Signature} et nom complet.
\end{enumerate}
\end{propriete}

\begin{popfigure}
\begin{center}
\begin{tikzpicture}[node distance=0.25cm]
\node[draw=popblue, fill=popblueL, rounded corners, minimum width=8.5cm, minimum height=0.75cm] (a) {\textbf{Sender's address} (adresse de l'expéditeur)};
\node[draw=popblue, fill=popblueL, rounded corners, minimum width=8.5cm, minimum height=0.75cm, below=of a] (b) {\textbf{Date} : 15th May 2024};
\node[draw=popgreen, fill=popgreenL, rounded corners, minimum width=8.5cm, minimum height=0.75cm, below=of b] (c) {\textbf{Receiver's address} (destinataire)};
\node[draw=popgreen, fill=popgreenL, rounded corners, minimum width=8.5cm, minimum height=0.75cm, below=of c] (d) {\textbf{Salutation} : Dear Sir, / Dear Madam,};
\node[draw=poporange, fill=poporangeL, rounded corners, minimum width=8.5cm, minimum height=1.6cm, align=center, below=of d] (e) {\textbf{Body} : 1. présentation + motif \\ 2. compétences + expérience \quad 3. conclusion};
\node[draw=poppurple, fill=poppurpleL, rounded corners, minimum width=8.5cm, minimum height=0.75cm, below=of e] (f) {\textbf{Closing} : Yours faithfully, / Yours sincerely,};
\node[draw=poppurple, fill=poppurpleL, rounded corners, minimum width=8.5cm, minimum height=0.75cm, below=of f] (g) {\textbf{Signature} + nom complet};
\draw[-{Stealth}, thick, popdark] (a) -- (b);
\draw[-{Stealth}, thick, popdark] (b) -- (c);
\draw[-{Stealth}, thick, popdark] (c) -- (d);
\draw[-{Stealth}, thick, popdark] (d) -- (e);
\draw[-{Stealth}, thick, popdark] (e) -- (f);
\draw[-{Stealth}, thick, popdark] (f) -- (g);
\end{tikzpicture}
\end{center}
\legende{Les sept blocs de la lettre formelle anglaise, dans l'ordre obligatoire.}
\end{popfigure}

\begin{definition}[Les connecteurs utiles (linking words)]
Les connecteurs donnent du rythme et de la logique à votre texte :
\begin{center}
\begin{tabularx}{\linewidth}{|X|X|X|}
\hline
\rowcolor{popblueL} Connecteur & Traduction & Emploi \\
\hline
\eng{first} & d'abord, premièrement & introduire le premier argument \\
\hline
\eng{moreover} & de plus, en outre & ajouter une idée \\
\hline
\eng{however} & cependant & opposer deux idées \\
\hline
\eng{therefore} & donc, par conséquent & exprimer une conséquence \\
\hline
\eng{in conclusion} & en conclusion & terminer le texte \\
\hline
\end{tabularx}
\end{center}
Exemple : \eng{“First, I speak two languages. Moreover, I have experience. However, I am still a student; therefore, I can work only in July. In conclusion, I am ready to start immediately.”} « D'abord, je parle deux langues. De plus, j'ai de l'expérience. Cependant, je suis encore élève ; par conséquent, je ne peux travailler qu'en juillet. En conclusion, je suis prêt à commencer immédiatement. »
\end{definition}

\begin{attention}[Les trois erreurs fatales des francophones]
\begin{enumerate}
\item \textbf{Oubli du « s » à la 3e personne du singulier} : \eng{She work in a bank} ✗ $\rightarrow$ \eng{She work\textbf{s} in a bank} ✓.
\item \textbf{\eng{I am agree}} ✗ $\rightarrow$ \eng{I agree with you} ✓. \emph{Agree} est un \textbf{verbe}, pas un adjectif : on ne met jamais \emph{be} devant. De même : \eng{I disagree} ✓ (jamais \eng{I am not agree} ✗).
\item \textbf{Mélange des temps} : ne passez pas du présent au passé sans raison. Si vous décrivez votre métier actuel, restez au présent simple : \eng{My father is a driver. He works in Douala. He earns 150,000 francs a month.}
\end{enumerate}
\end{attention}

\begin{exoresolu}[Corriger et compléter]
\textbf{Énoncé.} (a) Corrigez les trois erreurs dans cette phrase : \eng{My mother she work in a bank and I am agree that her job is more good than my father's job.} (b) Complétez avec le bon connecteur (first, however, therefore) : \eng{I want to become a pilot. \ldots, it is a well-paid job. \ldots, the training is very expensive; \ldots, I must save money now.}
\tcblower
\textbf{Solution.} (a) Trois erreurs : le sujet double « My mother she », l'absence du « s » à « work », et « I am agree » + « more good ». Correction : \eng{My mother works in a bank and I agree that her job is better than my father's job.} « Ma mère travaille dans une banque et je reconnais que son métier est meilleur que celui de mon père. » (b) \eng{I want to become a pilot. \textbf{First}, it is a well-paid job. \textbf{However}, the training is very expensive; \textbf{therefore}, I must save money now.} « Je veux devenir pilote. D'abord, c'est un métier bien payé. Cependant, la formation coûte très cher ; par conséquent, je dois économiser dès maintenant. »
\end{exoresolu}

\begin{exercice}[Your turn to write]
Rédigez une courte lettre de candidature (80 à 100 mots) pour un emploi de vacances dans une banque, un hôtel ou un supermarché de votre ville. Respectez les sept blocs de la lettre formelle, utilisez au moins trois connecteurs (\eng{first, moreover, however, therefore, in conclusion}) et un comparatif. Vérifiez le « s » de la 3e personne et n'écrivez jamais \eng{I am agree}.
\end{exercice}

\begin{corrige}[Exercice — Your turn to write]
Modèle de production attendue : adresse de l'expéditeur en haut, date, adresse du destinataire (\eng{The Manager, SGBC Bank, Douala}), salutation \eng{Dear Sir,}. Paragraphe 1 : \eng{I am writing to apply for the position of cashier during the holidays.} Paragraphe 2 : \eng{First, I am good at mathematics. Moreover, I speak English and French. However, I am a student; therefore, I can work only in August.} Paragraphe 3 : \eng{In conclusion, I am available for an interview at any time.} Clôture : \eng{Yours faithfully,} + signature. Barème indicatif : structure (5 pts), connecteurs (3 pts), grammaire et 3e personne (2 pts).
\end{corrige}

\begin{aretenir}[L'essentiel de la section]
Pour parler des métiers : \eng{“What do you do for a living?”} — \eng{“I'd like to become\ldots”} — comparatifs avec \emph{-er/more\ldots than}. La voix monte ↗ pour les questions yes/no et descend ↘ pour les questions en Wh-. Pour écrire une lettre de candidature : sept blocs dans l'ordre, trois paragraphes, connecteurs (\eng{first, moreover, however, therefore, in conclusion}), formule \eng{Yours faithfully}. Et surtout : jamais \eng{I am agree}, jamais de verbe sans « s » à la 3e personne !
\end{aretenir}

\coursec{Evaluation, self-assessment \& remediation}

Après avoir consolidé l'expression orale et écrite dans la section précédente, il est temps de vérifier ce que chaque élève a réellement retenu du chapitre \eng{Economic Life and Occupations}. L'évaluation n'est pas une punition : c'est une photographie de vos compétences à un instant donné. Et la remédiation est le traitement qui suit le diagnostic. Cette dernière section vous donne une épreuve blanche complète de type examen, une méthode d'auto-évaluation honnête et trois parcours de remédiation différenciés.

\courssub{Mock exam : a 30-minute timed test}

\begin{methode}[How to sit a mock exam]
\begin{enumerate}
\item \textbf{Read all the instructions first.} Ne commencez jamais à écrire avant d'avoir lu l'ensemble de l'épreuve : vous saurez ainsi répartir votre temps.
\item \textbf{Budget your time.} Trente minutes pour trente points : comptez environ une minute par point, et gardez trois minutes à la fin pour relire.
\item \textbf{Answer what you know first.} Une question bloquante ? Entourez-la, passez à la suivante, revenez-y ensuite.
\item \textbf{Write in full sentences} dans les parties grammaire et expression écrite : un mot isolé sans contexte perd souvent le point.
\item \textbf{Never leave a blank.} En compréhension, une réponse probable vaut mieux qu'une case vide.
\end{enumerate}
\end{methode}

\begin{savaistu}[Time management at the Baccalauréat]
Au Baccalauréat, la gestion du temps compte autant que les connaissances. La règle d'or des correcteurs : environ \cle{une minute par point}. Une épreuve de 20 points en 2 heures vous laisse même du temps de relecture — utilisez-le ! Beaucoup de copies perdent des points non pas par ignorance, mais parce que la dernière question, souvent la plus rentable, n'a pas été traitée. Entraînez-vous toujours avec une montre.
\end{savaistu}

\textbf{Part A — Reading comprehension (10 marks, 10 minutes).} Lisez le texte suivant, puis répondez aux questions.

\begin{textebox}[A young entrepreneur in Bafoussam]
When Divine Ngassa finished secondary school in Bafoussam, she did not wait for an office job. She noticed that many families in her neighbourhood bought vegetables at the central market every morning, often queuing for a long time. With a small loan from her aunt, she started a delivery service: customers send her a message the evening before, and she brings fresh tomatoes, onions and plantains to their doors before seven o'clock.

At first, business was slower than she expected. Her bicycle was old, and the rainy season made the roads difficult. But Divine was more determined than most people. She saved money for six months and bought a second-hand motorbike. Now her deliveries are faster and cheaper than those of the big shops, and she employs two of her former classmates.

“My plan is simple,” she says. “I want to open a small shop near the market and hire five more young people. Unemployment is a serious problem here, but young people are more creative than they think. If you work harder than the others and you are honest with your customers, you will succeed.”

Divine's story shows that an occupation is not always something you find — sometimes it is something you create.
\end{textebox}

\textbf{Glossaire} : \emph{loan} (n.) : prêt ; \emph{to queue} (v.) : faire la queue ; \emph{second-hand} (adj.) : d'occasion ; \emph{to hire} (v.) : embaucher ; \emph{unemployment} (n.) : chômage ; \emph{honest} (adj.) : honnête.

\textbf{Questions.}
\begin{enumerate}
\item Where does Divine live and work? (1 mark)
\item What problem did she notice in her neighbourhood? (2 marks)
\item How did she finance her business at the beginning? (2 marks)
\item Find two comparatives in the text and quote them. (2 marks)
\item What are Divine's plans for the future? (2 marks)
\item In your own words, what lesson does the text teach? (1 mark)
\end{enumerate}

\textbf{Part B — Grammar and vocabulary (10 marks, 10 minutes).}
\begin{enumerate}
\item Complete with the correct comparative or superlative form: (a) A doctor earns \dots{} (much) money than a farmer. (b) This is the \dots{} (busy) market in the town. (c) My new job is \dots{} (good) than my old one. (3 marks)
\item Match the job to the workplace: teacher, nurse, mechanic, cashier / garage, school, supermarket, hospital. (2 marks)
\item Fill in the gaps with a word from the box (\emph{salary, employer, interview, apply, resign}): (a) Before you get a job, you usually attend an \dots{}. (b) I want to \dots{} for the post of secretary. (c) Her monthly \dots{} is 150 000 francs. (3 marks)
\item Put the words in the correct order: am / I / writing / apply / to / for / this / post. (2 marks)
\end{enumerate}

\textbf{Part C — Short writing (10 marks, 10 minutes).} Write a short paragraph (60--80 words) about your dream job. Include: the name of the job, two reasons for your choice (use at least one comparative), and one quality needed for this occupation.

\courssub{Collective correction : justifying the answers}

\begin{exoresolu}[Correcting Part A with justification]
Énoncé : Question 2 — \eng{What problem did Divine notice in her neighbourhood?}
\tcblower
Solution : \eng{She noticed that many families had to queue for a long time to buy vegetables at the central market.} (« Elle a remarqué que beaucoup de familles devaient faire longtemps la queue pour acheter des légumes au marché central. »)
Justification : la réponse se trouve dans la deuxième phrase du texte. Au tableau, le professeur souligne la phrase source et montre qu'une bonne réponse reprend les mots du texte (\eng{families}, \eng{queuing}) sans les copier mot pour mot. Une réponse comme \eng{She noticed a problem} est trop vague : elle ne vaut qu'un demi-point car elle ne précise pas \emph{quel} problème.
\end{exoresolu}

\begin{center}
\begin{tabularx}{\linewidth}{|X|X|X|}
\hline
\rowcolor{popblueL} Question & Réponse attendue & Erreur fréquente \\
\hline
Q1 & \eng{She lives and works in Bafoussam.} & Répondre seulement \eng{Bafoussam} sans phrase complète \\
\hline
Q3 & \eng{With a small loan from her aunt.} & Confondre \eng{loan} (prêt) et \eng{gift} (cadeau) \\
\hline
Q4 & \eng{slower than, faster, cheaper, more determined} (2 suffisent) & Citer un superlatif au lieu d'un comparatif \\
\hline
Q5 & \eng{She wants to open a shop and hire five more young people.} & Oublier la deuxième partie du projet \\
\hline
Q6 & Réponse personnelle, ex. \eng{We can create our own jobs.} & Recopier une phrase du texte sans la comprendre \\
\hline
\end{tabularx}
\end{center}

Correction de la partie B : 1(a) \eng{more} ; 1(b) \eng{busiest} ; 1(c) \eng{better} ; 2. teacher--school, nurse--hospital, mechanic--garage, cashier--supermarket ; 3(a) \eng{interview} ; 3(b) \eng{apply} ; 3(c) \eng{salary} ; 4. \eng{I am writing to apply for this post.} (« Je vous écris pour postuler à ce poste. ») — c'est la formule canonique d'ouverture d'une lettre de candidature ; toute autre permutation des mots (\eng{I am writing for apply to this post}) est fautive car \eng{to apply for} est un bloc inséparable.

\begin{attention}[Pièges révélés par la correction]
\begin{itemize}
\item \eng{more good} n'existe pas : le comparatif de \eng{good} est \cle{better} (irrégulier), comme \eng{bad} $\rightarrow$ \eng{worse}.
\item Ne confondez pas \eng{apply for a job} (postuler à un emploi) et \eng{apply a rule} (appliquer une règle) : la préposition change le sens.
\item \eng{The most busy} est fautif : les adjectifs courts prennent \eng{-est} : \eng{the busiest} (attention à l'orthographe : \emph{y} devient \emph{i} devant \eng{-est}).
\item \eng{Salary} se prononce « sa-la-ri » ; ne prononcez pas le premier \emph{a} comme dans « sale ».
\end{itemize}
\end{attention}

\courssub{Self-assessment grid : knowing what you can do}

\begin{methode}[How to self-assess honestly]
\begin{enumerate}
\item \textbf{Test yourself without your notes.} Fermez le cahier et essayez la tâche (nommer 10 métiers, conjuguer un comparatif, rédiger une phrase de lettre de candidature).
\item \textbf{Then compare with the lesson.} Rouvrez le cours et vérifiez point par point.
\item \textbf{Measure the gap.} C'est l'écart entre ce que vous avez produit seul et la réponse correcte qui indique le travail à faire — pas votre impression générale.
\item \textbf{Tick honestly.} Cocher \eng{well} alors qu'on a hésité trois fois ne sert à rien : personne ne vous note ici, sauf vous-même.
\end{enumerate}
\end{methode}

\begin{center}
\begin{tabularx}{\linewidth}{|X|c|c|c|}
\hline
\rowcolor{popblueL} \eng{I can...} & \eng{well} & \eng{fairly} & \eng{with difficulty} \\
\hline
\eng{name 10 jobs and their workplaces} & & & \\
\hline
\eng{use comparatives and superlatives correctly} & & & \\
\hline
\eng{talk about salaries, employers and employees} & & & \\
\hline
\eng{write a short application letter} & & & \\
\hline
\eng{understand a text about economic life} & & & \\
\hline
\eng{describe my dream job orally} & & & \\
\hline
\end{tabularx}
\end{center}

\textbf{Bilan individuel.} Après avoir coché la grille, chaque élève complète dans son cahier : \eng{My two strong points are: 1. \dots{} 2. \dots{} / I still need to work on: 1. \dots{} 2. \dots{}} (« Mes deux points forts sont\ldots{} / Je dois encore travailler\ldots{} »). Ce bilan doit être précis : écrire \eng{I am bad at English} est interdit ; écrire \eng{I confuse “better” and “more good”} est utile, car il nomme exactement la difficulté à traiter.

\courssub{Differentiated remediation paths}

\begin{propriete}[The four steps of remediation]
La remédiation suit toujours quatre étapes, dans cet ordre :
\begin{enumerate}
\item \textbf{Identifier l'erreur} : nommer précisément ce qui est faux (par exemple : j'ai écrit \eng{more cheap} au lieu de \eng{cheaper}).
\item \textbf{Comprendre la règle} : relire la règle de grammaire ou la fiche de vocabulaire correspondante.
\item \textbf{S'exercer} : refaire au moins trois exercices ciblés sur ce point précis.
\item \textbf{Se re-tester} : vérifier, sans le cours, que l'erreur a disparu.
\end{enumerate}
Sauter une étape (par exemple s'exercer sans avoir compris la règle) rend la remédiation inefficace.
\end{propriete}

\begin{popfigure}
\begin{center}
\begin{tikzpicture}[node distance=0.55cm, every node/.style={font=\small}]
\node[draw=poporange, fill=poporangeL, rounded corners, align=center, text width=3.1cm] (e) {\textbf{Erreur détectée}\\ \eng{more cheap}};
\node[draw=popblue, fill=popblueL, rounded corners, align=center, text width=3.1cm, right=of e] (i) {\textbf{Identifier la règle}\\ comparatif des adjectifs courts};
\node[draw=poppurple, fill=poppurpleL, rounded corners, align=center, text width=3.1cm, right=of i] (r) {\textbf{Revoir la règle}\\ \eng{cheap $\rightarrow$ cheaper}};
\node[draw=popgreen, fill=popgreenL, rounded corners, align=center, text width=3.1cm, below=of r] (x) {\textbf{Refaire un exercice}\\ 3 phrases ciblées};
\node[draw=popgold, fill=popgoldL, rounded corners, align=center, text width=3.1cm, left=of x] (s) {\textbf{Réussite}\\ test blanc réussi};
\draw[-{Stealth}, thick, popdark] (e) -- (i);
\draw[-{Stealth}, thick, popdark] (i) -- (r);
\draw[-{Stealth}, thick, popdark] (r) -- (x);
\draw[-{Stealth}, thick, popdark] (x) -- (s);
\draw[-{Stealth}, thick, popmuted, dashed] (s.west) .. controls +(-0.8,0.8) .. node[above, align=center, font=\scriptsize]{sinon, on recommence} (e.south);
\end{tikzpicture}
\end{center}
\legende{Organigramme de remédiation : de l'erreur à la réussite en quatre étapes.}
\end{popfigure}

\begin{attention}[Do not revise what you already master]
Beaucoup d'élèves relisent confortablement les leçons qu'ils connaissent déjà : cela donne bonne conscience, mais ne fait pas progresser. La révision efficace \cle{cible les difficultés}. Si vous maîtrisez les métiers mais échouez aux superlatifs, consacrez 80 \% de votre temps aux superlatifs. La grille d'auto-évaluation sert précisément à repérer ces zones fragiles.
\end{attention}

\textbf{Fiche A — Vocabulary (matching game).} Pour les élèves faibles en lexique : associer chaque mot à sa définition, puis rejouer le lendemain sans regarder.

\begin{center}
\begin{tabularx}{\linewidth}{|X|X|}
\hline
\rowcolor{popgreenL} Word & Definition \\
\hline
\eng{employer} & \eng{a person or company that gives work to people} \\
\hline
\eng{employee} & \eng{a person who works for an employer} \\
\hline
\eng{salary} & \eng{the money you receive every month for your work} \\
\hline
\eng{interview} & \eng{a meeting where an employer asks you questions} \\
\hline
\eng{to resign} & \eng{to leave your job voluntarily} \\
\hline
\eng{unemployed} & \eng{without a job} \\
\hline
\end{tabularx}
\end{center}

\textbf{Fiche B — Grammar (targeted exercises).} Pour les élèves faibles en grammaire : dix phrases à compléter sur le point précis qui a échoué (comparatifs, superlatifs ou ordre des mots), avec re-test final de trois phrases sans aide. Exemple de série ciblée : (a) \eng{A motorbike is \dots{} (fast) than a bicycle.} (b) \eng{January is the \dots{} (dry) month of the year here.} (c) \eng{Her new salary is \dots{} (good) than the old one.} Corrigé : (a) \eng{faster} ; (b) \eng{driest} ; (c) \eng{better}.

\textbf{Fiche C — Production (guided rewriting).} Pour les élèves faibles en expression écrite : reprendre son paragraphe de la partie C, corriger les erreurs signalées, puis le réécrire en suivant le plan : phrase 1 = le métier, phrases 2--3 = deux raisons avec un comparatif, phrase 4 = une qualité nécessaire. Modèle : \eng{My dream job is to be a nurse in Yaoundé. This job is more useful than many others because nurses save lives every day. It is also better paid than it was before. A good nurse must be patient and hardworking.} (« Le métier de mes rêves est d'être infirmière à Yaoundé. Ce métier est plus utile que beaucoup d'autres parce que les infirmières sauvent des vies chaque jour. Il est aussi mieux payé qu'avant. Une bonne infirmière doit être patiente et travailleuse. »)

\begin{exercice}[Progress contract]
Sur une fiche que vous conserverez dans votre cahier, écrivez une phrase d'engagement commençant par \eng{Next time, I will\dots{}}. Exemples : \eng{Next time, I will learn five new job words every week.} / \eng{Next time, I will check my comparatives before handing in my paper.} Signez et datez votre contrat : le relire avant le prochain chapitre vous rappellera votre objectif.
\end{exercice}

\begin{corrige}[Contrat de progrès]
Il n'y a pas de réponse unique : le contrat est personnel. Un bon contrat est \textbf{précis} (une action concrète), \textbf{mesurable} (\eng{five new words}, pas \eng{more words}) et \textbf{réaliste}. Le professeur vérifie que chaque engagement correspond bien à un point faible identifié dans la grille d'auto-évaluation : c'est la cohérence entre le diagnostic et le remède qui fait la valeur du contrat.
\end{corrige}

\begin{aretenir}[What to remember from this section]
Évaluer, c'est mesurer ; s'auto-évaluer, c'est se mesurer soi-même, honnêtement ; remédier, c'est réparer exactement là où c'est cassé. Retenez la boucle : \cle{test $\rightarrow$ gap $\rightarrow$ rule $\rightarrow$ practice $\rightarrow$ re-test}. Un élève qui connaît ses points faibles et travaille dessus progresse toujours plus vite qu'un élève qui relit indifféremment tout son cahier. \eng{Know your weaknesses, and they will become your strengths.} (« Connais tes faiblesses, et elles deviendront tes forces. »)
\end{aretenir}

\newpage

\coursec{Activité d'intégration}

\courssub{Objectifs de l'activité}

Cette activité d'intégration clôt le chapitre 2, \eng{Economic Life and Occupations}. Elle mobilise, dans une situation de communication authentique et ludique, l'ensemble des acquis du chapitre :

\begin{itemize}
\item le \cle{vocabulaire} des métiers, des lieux de travail, des salaires et des secteurs économiques (\eng{job, employer, employee, salary, wages, to hire, to apply for, skilled, sector});
\item la \cle{grammaire} du chapitre : le \eng{present simple} pour décrire un métier, les \cle{comparatifs} et \cle{superlatifs} pour comparer des candidats, les questions avec les \eng{wh-} words et les auxiliaires \eng{do/does};
\item l'\cle{expression orale} : se présenter, répondre à un entretien, justifier un choix;
\item l'\cle{expression écrite} : rédiger un mini-CV et une courte lettre de candidature (\eng{application letter});
\item l'\cle{auto-évaluation} : identifier ses réussites et un point à remédier.
\end{itemize}

\courssub{Warm-up : la carte mentale du chapitre}

Avant de commencer, observons la carte mentale qui résume le chapitre 2. Recopie-la dans ton cahier et ajoute au moins deux mots de ton choix dans chaque branche.

\begin{popfigure}
\begin{center}\tcbox[colback=popink,colframe=popink,coltext=white,arc=9pt,boxsep=4pt,left=6pt,right=6pt]{\bfseries\small Chapter 2 Economic Life}\end{center}
\vspace{-2pt}
\begin{tcbraster}[raster columns=3,raster equal height=rows,raster column skip=6pt,raster row skip=6pt,enhanced,blankest]
\begin{tcolorbox}[enhanced,colback=white,colframe=poporange,boxrule=0.9pt,arc=3pt,title={Jobs},fonttitle=\bfseries\scriptsize,coltitle=white,colbacktitle=poporange,left=4pt,right=4pt,top=2pt,bottom=2pt,boxsep=1pt,toptitle=1.5pt,bottomtitle=1.5pt,halign=left]
\begin{itemize}[nosep,leftmargin=*,itemsep=1pt,font=\scriptsize]
\item nurse
\item cashier
\end{itemize}
\end{tcolorbox}
\begin{tcolorbox}[enhanced,colback=white,colframe=popgreen,boxrule=0.9pt,arc=3pt,title={Workplaces},fonttitle=\bfseries\scriptsize,coltitle=white,colbacktitle=popgreen,left=4pt,right=4pt,top=2pt,bottom=2pt,boxsep=1pt,toptitle=1.5pt,bottomtitle=1.5pt,halign=left]
\begin{itemize}[nosep,leftmargin=*,itemsep=1pt,font=\scriptsize]
\item hospital
\item bank
\end{itemize}
\end{tcolorbox}
\begin{tcolorbox}[enhanced,colback=white,colframe=poppurple,boxrule=0.9pt,arc=3pt,title={Money},fonttitle=\bfseries\scriptsize,coltitle=white,colbacktitle=poppurple,left=4pt,right=4pt,top=2pt,bottom=2pt,boxsep=1pt,toptitle=1.5pt,bottomtitle=1.5pt,halign=left]
\begin{itemize}[nosep,leftmargin=*,itemsep=1pt,font=\scriptsize]
\item salary
\item wages
\end{itemize}
\end{tcolorbox}
\begin{tcolorbox}[enhanced,colback=white,colframe=poppink,boxrule=0.9pt,arc=3pt,title={Grammar},fonttitle=\bfseries\scriptsize,coltitle=white,colbacktitle=poppink,left=4pt,right=4pt,top=2pt,bottom=2pt,boxsep=1pt,toptitle=1.5pt,bottomtitle=1.5pt,halign=left]
\begin{itemize}[nosep,leftmargin=*,itemsep=1pt,font=\scriptsize]
\item present simple
\item comparatives
\end{itemize}
\end{tcolorbox}
\begin{tcolorbox}[enhanced,colback=white,colframe=popgold,boxrule=0.9pt,arc=3pt,title={Applying},fonttitle=\bfseries\scriptsize,coltitle=white,colbacktitle=popgold,left=4pt,right=4pt,top=2pt,bottom=2pt,boxsep=1pt,toptitle=1.5pt,bottomtitle=1.5pt,halign=left]
\begin{itemize}[nosep,leftmargin=*,itemsep=1pt,font=\scriptsize]
\item CV
\item interview
\end{itemize}
\end{tcolorbox}
\end{tcbraster}

\legende{Carte mentale de synthèse du chapitre 2 : les cinq familles d'acquis à mobiliser.}
\end{popfigure}

\courssub{Situation}

Le proviseur du lycée de Bafoussam organise, à la fin du deuxième trimestre, un \eng{Job Fair} (salon de l'emploi) pour sensibiliser les élèves de Première au monde du travail. Cinq employeurs de la région de l'Ouest ont accepté de jouer le jeu : une banque, un hôpital, une exploitation agricole, une boutique et une école primaire. Chacun propose un poste. Voici l'annonce publiée dans le journal du lycée :

\begin{textebox}[The Lion's Voice — School Magazine, Bafoussam]
\textbf{JOB FAIR AT SCHOOL — FRIDAY 14 MARCH, 10 a.m.}

Are you looking for a holiday job? Five local employers are coming to our school! Come with your CV and be ready for a short interview in English.

\textbf{1. Bamiléké Commercial Bank} needs a \textbf{cashier}. You must be good with numbers and polite with customers. Salary: 60,000 FCFA per month.

\textbf{2. Regional Hospital} needs a \textbf{receptionist}. You must speak English and French and be friendly. Salary: 55,000 FCFA per month.

\textbf{3. Green Valley Farm} needs a \textbf{farm assistant}. The work is hard but interesting: maize, coffee and vegetables. Salary: 50,000 FCFA per month plus free meals.

\textbf{4. Chez Mami Boutique} needs a \textbf{shop assistant}. You must be honest and smiling. Salary: 40,000 FCFA per month plus a small commission on sales.

\textbf{5. Bright Future Primary School} needs an \textbf{English tutor} for class 6 pupils. You must be patient and speak good English. Salary: 45,000 FCFA per month.

Send your application letter to the school office before Wednesday, or come directly to the Job Fair!
\end{textebox}

\begin{definition}[Glossaire utile]
\begin{itemize}
\item \eng{cashier} (n.) : caissier / caissière
\item \eng{receptionist} (n.) : réceptionniste
\item \eng{farm assistant} (n.) : aide-agriculteur
\item \eng{shop assistant} (n.) : vendeur / vendeuse
\item \eng{tutor} (n.) : répétiteur, tuteur
\item \eng{salary} (n.) : salaire (mensuel, versé aux employés qualifiés) ; \eng{wages} (n. pl.) : salaire (à l'heure ou à la semaine, souvent pour un travail manuel)
\item \eng{commission} (n.) : commission (pourcentage sur les ventes)
\item \eng{to hire} (v.) : embaucher ; \eng{to apply for a job} : postuler à un emploi
\item \eng{customer} (n.) : client ; \eng{advertisement} (n.) : annonce, publicité
\end{itemize}
\end{definition}

\courssub{Consignes}

\begin{enumerate}
\item \textbf{Préparation des rôles (15 min).} La classe est divisée en deux groupes. Le groupe A (les \eng{employers}) reçoit les fiches de poste et prépare trois questions d'entretien. Le groupe B (les \eng{candidates}) rédige un mini-CV de 6 à 8 lignes : nom, âge, études, expérience, qualités, poste souhaité.
\item \textbf{Compréhension des annonces (10 min).} Lis attentivement les cinq annonces. Pour chaque poste, relève : le métier, deux qualités exigées et le salaire. Souligne les mots de vocabulaire du chapitre.
\item \textbf{Les entretiens (25 min).} Les candidats circulent de stand en stand. Chaque entretien dure deux minutes : le candidat se présente (\eng{Good morning, my name is... I would like to apply for the job of...}), l'employeur pose ses questions (\eng{What do you do? Why do you want this job? Are you more experienced than the other candidates?}).
\item \textbf{La comparaison (10 min).} Chaque employeur compare les candidats à voix haute avec des comparatifs et des superlatifs, puis annonce son choix : \eng{Mariam is more patient than Jean, but Paul is the most experienced. I choose Paul.}
\item \textbf{Inversion des rôles (25 min).} Les groupes échangent : le groupe B devient employeur, le groupe A devient candidat. Nouveaux entretiens.
\item \textbf{Production écrite (20 min).} Chaque élève rédige une courte lettre de candidature (10 à 12 lignes) pour le poste de son choix, en suivant le modèle vu en classe : adresse, date, objet, salutation, corps en deux paragraphes, formule de politesse.
\item \textbf{Auto-évaluation (10 min).} Chaque élève remplit la grille ci-dessous et note \textbf{un point à remédier} avec sa stratégie personnelle.
\end{enumerate}

\begin{center}
\begin{tabularx}{\linewidth}{|X|c|c|}
\hline
\rowcolor{popblueL} \textbf{Self-assessment grid} & \textbf{Yes} & \textbf{Not yet} \\
\hline
I can name at least ten jobs and workplaces in English. & & \\
\hline
I can describe a job with the present simple. & & \\
\hline
I can compare two candidates with comparatives. & & \\
\hline
I can use superlatives to justify a choice. & & \\
\hline
I can answer interview questions fluently. & & \\
\hline
I can write a short application letter. & & \\
\hline
\end{tabularx}
\end{center}

\courssub{Ressources à mobiliser}

\begin{propriete}[Rappel 1 : le present simple pour décrire un métier]
On utilise le \eng{present simple} pour les faits et les habitudes professionnelles : \eng{A nurse works in a hospital. She looks after patients.} « Une infirmière travaille dans un hôpital. Elle s'occupe des patients. »

Attention à la 3\textsuperscript{e} personne du singulier : \eng{work\textbf{s}}, \eng{look\textbf{s}}, \eng{teach\textbf{es}}, \eng{stud\textbf{ies}} (le \eng{-y} devient \eng{-ies} après consonne), \eng{go\textbf{es}}, \eng{ha\textbf{s}} (irrégulier).

Forme négative et interrogative avec l'auxiliaire \eng{do/does} : \eng{A cashier does not work at night. Does a tutor earn much money?} « Un caissier ne travaille pas la nuit. Un répétiteur gagne-t-il beaucoup d'argent ? »
\end{propriete}

\begin{propriete}[Rappel 2 : comparatifs et superlatifs]
\begin{center}
\begin{tabular}{|l|l|l|}
\hline
\rowcolor{popblueL} Adjectif & Comparatif & Superlatif \\
\hline
adjectifs courts (tall) & taller than & the tallest \\
\hline
adjectifs en -y (friendly) & friendlier than & the friendliest \\
\hline
adjectifs longs (patient) & more patient than & the most patient \\
\hline
good / bad & better / worse than & the best / the worst \\
\hline
\end{tabular}
\end{center}
Attention aux adjectifs courts terminés par consonne + voyelle + consonne : on double la consonne finale — \eng{big, bigger than, the biggest} ; \eng{thin, thinner than, the thinnest}.

Exemple : \eng{Paul is more experienced than John, but Mariam is the most qualified candidate.} « Paul est plus expérimenté que John, mais Mariam est la candidate la plus qualifiée. »
\end{propriete}

\begin{attention}[Erreurs fréquentes des francophones]
\begin{itemize}
\item Ne dis pas \eng{I am agree} : dis \eng{I agree}. Le verbe \eng{agree} ne prend pas l'auxiliaire \eng{be}.
\item \eng{an economic problem} signifie « un problème économique » (qui concerne l'économie) ; pour « économique » au sens de « qui fait faire des économies », utilise \eng{cheap} ou \eng{economical}.
\item L'ordre des mots dans les questions : \eng{Why do you want this job?} et non \eng{Why you want this job?} — l'auxiliaire \eng{do} est obligatoire.
\item \eng{money} est indénombrable : \eng{I earn a lot of money}, jamais \eng{moneys}.
\item Faux ami : \eng{a formation} n'existe pas pour « une formation » ; dis \eng{training} ou \eng{a training course}.
\item Prononciation : \eng{work} se prononce « weurk », \eng{salary} « sa-la-ri », \eng{receptionist} « ri-sèp-cheu-niste ».
\end{itemize}
\end{attention}

\begin{methode}[Rédiger une lettre de candidature courte]
\begin{enumerate}
\item En haut à droite : ton adresse et la date.
\item À gauche : l'adresse de l'employeur.
\item L'objet : \eng{Application for the post of...}
\item La salutation : \eng{Dear Sir or Madam,} (destinataire inconnu) ou \eng{Dear Mr Ngo,} (destinataire connu).
\item Paragraphe 1 : le poste visé et où tu as vu l'annonce.
\item Paragraphe 2 : tes qualités et ton expérience (comparatifs bienvenus).
\item Formule finale : \eng{Yours faithfully,} (après \eng{Dear Sir or Madam}) ou \eng{Yours sincerely,} (après un nom), puis ta signature.
\end{enumerate}
\end{methode}

\courssub{Entraînement de type examen}

\begin{exoresolu}[Compréhension écrite et grammaire — sujet type Probatoire]
Énoncé : read the Job Fair advertisement again and answer the following questions in complete sentences.

\begin{enumerate}
\item Which job offers the highest salary?
\item What two qualities must the receptionist have?
\item Which employer offers free meals?
\item Complete with the correct comparative or superlative: \eng{The cashier's salary is \dots\ (high) than the shop assistant's, but the tutor's job is \dots\ (interesting) for me.}
\item Put the words in the correct order: \eng{you / why / do / this job / want / ?}
\end{enumerate}

\tcblower

\textbf{Solution détaillée :}
\begin{enumerate}
\item \eng{The cashier's job at Bamiléké Commercial Bank offers the highest salary: 60,000 FCFA per month.} (On reprend le superlatif \eng{the highest} et on cite le chiffre pour justifier.)
\item \eng{The receptionist must speak English and French and be friendly.} (Deux exigences relevées dans l'annonce : la langue et la qualité personnelle.)
\item \eng{Green Valley Farm offers free meals to the farm assistant.}
\item \eng{The cashier's salary is \textbf{higher} than the shop assistant's, but the tutor's job is \textbf{the most interesting} for me.} (\eng{high} est un adjectif court : \eng{higher} ; \eng{interesting} est un adjectif long : \eng{the most interesting}.)
\item \eng{Why do you want this job?} (Mot interrogatif + auxiliaire \eng{do} + sujet + verbe.)
\end{enumerate}
\end{exoresolu}

\begin{exercice}[À toi de jouer — remédiation]
\begin{enumerate}
\item Conjugue au present simple : \eng{A farmer (grow) maize. He (not work) in an office. (he / earn) a high salary?}
\item Écris le comparatif et le superlatif de : \eng{hard, honest, good, bad, qualified}.
\item Rédige trois questions d'entretien avec \eng{what}, \eng{where} et \eng{why}.
\item Corrige les erreurs : \eng{She work in a bank and she earn many moneys.}
\end{enumerate}
\end{exercice}

\begin{corrige}[Exercice de remédiation]
\begin{enumerate}
\item \eng{A farmer \textbf{grows} maize. He \textbf{does not work} in an office. \textbf{Does he earn} a high salary?} (3\textsuperscript{e} personne : \eng{grows} ; négation et question avec \eng{does} + base verbale.)
\item \eng{hard, harder than, the hardest ; honest, more honest than, the most honest ; good, better than, the best ; bad, worse than, the worst ; qualified, more qualified than, the most qualified.}
\item Exemples : \eng{What do you do? Where do you work? Why do you want this job?}
\item \eng{She \textbf{works} in a bank and she \textbf{earns a lot of money}.} (\eng{-s} à la 3\textsuperscript{e} personne ; \eng{money} est indénombrable : \eng{a lot of money}, jamais \eng{many moneys}.)
\end{enumerate}
\end{corrige}

\courssub{Proposition de production et corrigé commenté}

\begin{exoresolu}[Modèle de lettre de candidature]
Énoncé : write a short application letter (10--12 lines) for the post of English tutor at Bright Future Primary School.

\tcblower

\textbf{Modèle de production :}

\begin{center}
\begin{minipage}{0.9\linewidth}
\emph{Government High School, P.O. Box 45, Bafoussam}\\\emph{14 March 2024}

\emph{The Headmaster,}\\\emph{Bright Future Primary School,}\\\emph{Bafoussam}

\emph{\textbf{Application for the post of English tutor}}

\emph{Dear Sir,}

\emph{I saw your advertisement in our school magazine and I would like to apply for the post of English tutor for class 6 pupils.}

\emph{I am a Form Five student at Government High School Bafoussam. I am eighteen years old and I speak English and French fluently. I am patient and I like working with children. Last year, I helped my younger brother with his English homework every evening, and his marks became much better. I think I am more experienced than many other candidates because I have already taught small groups of pupils during the holidays.}

\emph{I am available every afternoon after 3 p.m. I hope you will consider my application.}

\emph{Yours faithfully,}\\\emph{Njoya Larissa}
\end{minipage}
\end{center}

\textbf{Commentaire des choix de langue :}
\begin{itemize}
\item Structure de la lettre respectée : adresses, date, objet, salutation, deux paragraphes, formule finale (\eng{Yours faithfully} car on ne connaît pas le destinataire).
\item \eng{I would like to apply for the post of...} : formule standard et polie de candidature.
\item Present simple pour les faits : \eng{I am a Form Five student}, \eng{I speak English and French}.
\item Past simple pour l'expérience passée : \eng{I helped my younger brother}, \eng{his marks became better} (verbe irrégulier \eng{become, became, become}).
\item Comparatif mobilisé comme demandé : \eng{I am more experienced than many other candidates}.
\item Vocabulaire du chapitre : \eng{advertisement, to apply for, post, available, application}.
\end{itemize}
\end{exoresolu}

\courssub{Bilan de l'activité}

\begin{aretenir}[Ce que le Job Fair m'a appris]
À l'issue de cette activité, je suis capable de :
\begin{itemize}
\item comprendre une annonce d'emploi en anglais et en dégager les informations essentielles (poste, qualités, salaire);
\item me présenter et répondre à un entretien court avec le \eng{present simple} et les questions en \eng{wh-};
\item comparer des candidats avec les comparatifs et justifier un choix avec les superlatifs (\eng{better than, the most experienced});
\item rédiger un mini-CV et une lettre de candidature courte et correcte;
\item m'auto-évaluer : je coche mes réussites dans la grille et j'écris \textbf{un point à remédier} (par exemple : \eng{I must revise the spelling of jobs} ou \eng{I must practise the third person -s}) avec une stratégie concrète (relire la fiche de grammaire, refaire deux exercices, interroger un camarade).
\end{itemize}
\end{aretenir}

\begin{savaistu}[Le monde du travail au Cameroun]
Au Cameroun, pays officiellement bilingue, maîtriser l'anglais est un véritable atout professionnel : de nombreuses entreprises internationales installées à Douala et à Yaoundé exigent des candidats capables de passer un entretien en anglais. Savoir rédiger une \eng{application letter} correcte est donc une compétence directement utile pour ton avenir — et c'est aussi une tâche classique des épreuves d'anglais au Probatoire et au Baccalauréat !
\end{savaistu}

\newpage

\coursec{Méthodes et exercices résolus}

\coursec{Integration Activities / Evaluation / Remediation (Chapter 2)}

\courssub{Warm-up : Chapter 2 Mind-map}

\begin{methode}[Activer ses connaissances : la carte mentale du chapitre]
Avant de commencer la révision, construisez (ou complétez) cette carte mentale. Elle structure les 4 piliers du chapitre : \cle{lexique}, \cle{grammaire}, \cle{compétences} et \cle{contexte camerounais}. Utilisez-la comme fiche de révision unique.
\end{methode}

\begin{popfigure}
\begin{center}\tcbox[colback=popink,colframe=popink,coltext=white,arc=9pt,boxsep=4pt,left=6pt,right=6pt]{\bfseries\small Economic Life \& Occupations}\end{center}
\vspace{-2pt}
\begin{tcbraster}[raster columns=2,raster equal height=rows,raster column skip=6pt,raster row skip=6pt,enhanced,blankest]
\begin{tcolorbox}[enhanced,colback=white,colframe=popgreen,boxrule=0.9pt,arc=3pt,title={Vocabulary},fonttitle=\bfseries\scriptsize,coltitle=white,colbacktitle=popgreen,left=4pt,right=4pt,top=2pt,bottom=2pt,boxsep=1pt,toptitle=1.5pt,bottomtitle=1.5pt,halign=left]
\begin{itemize}[nosep,leftmargin=*,itemsep=1pt,font=\scriptsize]
\item Jobs: farmer, trader, tailor, driver, nurse, engineer
\item Places: market, bank, workshop, office, factory, cooperative
\item Actions: earn, save, borrow, hire, produce, trade, apply for
\item Concepts: salary, profit, debt, loan, unemployment, taxes
\end{itemize}
\end{tcolorbox}
\begin{tcolorbox}[enhanced,colback=white,colframe=poporange,boxrule=0.9pt,arc=3pt,title={Grammar},fonttitle=\bfseries\scriptsize,coltitle=white,colbacktitle=poporange,left=4pt,right=4pt,top=2pt,bottom=2pt,boxsep=1pt,toptitle=1.5pt,bottomtitle=1.5pt,halign=left]
\begin{itemize}[nosep,leftmargin=*,itemsep=1pt,font=\scriptsize]
\item Present Simple: habits, facts (\eng{works})
\item Present Continuous: now (\eng{are closing})
\item Present Perfect: experience, since/for (\eng{has worked})
\item Past Simple: dated past (\eng{sold yesterday})
\item Future: \eng{will} vs \eng{going to}
\end{itemize}
\end{tcolorbox}
\begin{tcolorbox}[enhanced,colback=white,colframe=poppurple,boxrule=0.9pt,arc=3pt,title={Skills},fonttitle=\bfseries\scriptsize,coltitle=white,colbacktitle=poppurple,left=4pt,right=4pt,top=2pt,bottom=2pt,boxsep=1pt,toptitle=1.5pt,bottomtitle=1.5pt,halign=left]
\begin{itemize}[nosep,leftmargin=*,itemsep=1pt,font=\scriptsize]
\item Reading: skimming, scanning, vocab in context
\item Listening: job interview, radio report
\item Speaking: career plans, role-play interview
\item Writing: application letter, composition
\end{itemize}
\end{tcolorbox}
\begin{tcolorbox}[enhanced,colback=white,colframe=poppink,boxrule=0.9pt,arc=3pt,title={Cameroon Context},fonttitle=\bfseries\scriptsize,coltitle=white,colbacktitle=poppink,left=4pt,right=4pt,top=2pt,bottom=2pt,boxsep=1pt,toptitle=1.5pt,bottomtitle=1.5pt,halign=left]
\begin{itemize}[nosep,leftmargin=*,itemsep=1pt,font=\scriptsize]
\item Informal sector: bendskin, call-box, street vendors
\item Agriculture: cocoa, coffee, plantain, cooperative
\item Youth: GCE, vocational training, entrepreneurship
\item Institutions: MINCOMMERCE, MINJEC, GIC, banks
\end{itemize}
\end{tcolorbox}
\end{tcbraster}

\legende{Carte mentale de synthèse — Chapter 2 : Economic Life and Occupations}
\end{popfigure}

\courssub{Mobiliser ses acquis dans une situation d'intégration}

\begin{methode}[Réussir une tâche d'intégration sur la vie économique]
Une tâche d'intégration vous demande de réutiliser \emph{en même temps} le vocabulaire, la grammaire et les compétences du chapitre dans une situation réelle (lettre de candidature, article sur l'économie locale, dialogue avec un employeur).
\begin{enumerate}
\item \textbf{Lire la consigne deux fois} et souligner : le \cle{type de texte} attendu (letter, article, dialogue), le \cle{destinataire}, le \cle{nombre de mots}.
\item \textbf{Faire un mini-plan} : lister 5 à 8 mots-clés du chapitre à placer obligatoirement : \eng{unemployed}, \eng{to apply for a job}, \eng{salary}, \eng{trade}, \eng{customer}, \eng{to earn a living}, \eng{qualifications}.
\item \textbf{Choisir les temps avant d'écrire} : présent simple pour les habitudes et vérités générales (\eng{Farmers sell their crops at the market.}), present perfect pour l'expérience (\eng{I have worked as a shop assistant for two years.}), futur avec \eng{will} ou \eng{going to} pour les projets.
\item \textbf{Rédiger} en paragraphes courts : introduction, développement (2 arguments), conclusion.
\item \textbf{Relire} en traquant trois fautes : la 3\up{e} personne du singulier (\eng{he works}), les prépositions (\eng{to apply \textbf{for}}, \eng{to depend \textbf{on}}), l'ordre des mots (sujet + verbe + complément).
\end{enumerate}
\end{methode}

\begin{exoresolu}[Integration task : a letter of application]
\textbf{Énoncé.} You saw an advertisement in \emph{The Post} newspaper: “Shop assistant wanted at Bonaberi Market, Douala. Good English required.” Write a letter of application (80--100 words) to the manager, Mr. Ekane. Mention your studies, your experience and your qualities.
\tcblower
\textbf{Solution détaillée.}
\begin{enumerate}
\item \textbf{Analyse de la consigne} : lettre formelle, destinataire = Mr. Ekane, 80--100 mots, trois informations obligatoires (studies, experience, qualities).
\item \textbf{Vocabulaire mobilisé} : \eng{to apply for the post of}, \eng{qualifications}, \eng{experience}, \eng{hard-working}, \eng{I look forward to hearing from you}.
\item \textbf{Temps choisis} : present perfect pour l'expérience, présent simple pour les qualités.
\end{enumerate}
\textbf{Réponse rédigée :}

\eng{Dear Mr. Ekane,}

\eng{I am writing to apply for the post of shop assistant advertised in The Post. I passed my GCE Ordinary Level last year, and I speak English and French fluently. I have worked in my uncle's shop in Douala for two years, so I know how to welcome customers and handle money. I am hard-working, honest and punctual. I would be grateful for an interview. I look forward to hearing from you.}

\eng{Yours sincerely,}

\eng{Alima Ngo Bell}

(95 mots.) \textbf{Justifications} : \eng{to apply \textbf{for}} (préposition obligatoire) ; \eng{I \textbf{have worked}} = present perfect car l'expérience a un lien avec le présent ; \eng{I \textbf{look forward to hearing}} : après \eng{to look forward to}, le verbe prend la forme en -\emph{ing} ; formule de politesse \eng{Yours sincerely} car le nom du destinataire est connu.
\textbf{Conclusion} : toutes les exigences de la consigne sont couvertes, le ton est formel, aucune faute de préposition ni de conjugaison.
\end{exoresolu}

\courssub{Synthèse de grammaire : les temps du chapitre}

\begin{propriete}[Tableau récapitulatif : formes, emplois, marqueurs temporels]
Le tableau ci-dessous synthétise les quatre temps principaux travaillés au Chapter 2. Mémorisez la colonne « Marqueurs » : elle est la clé pour choisir le bon temps à l'examen.
\end{propriete}

\begin{center}
\begin{tabularx}{\linewidth}{|X|X|X|X|}
\hline
\rowcolor{popblueL}
\textbf{Temps} & \textbf{Forme (affirmative)} & \textbf{Emploi principal (Chapter 2)} & \textbf{Marqueurs temporels typiques} \\
\hline
Present Simple & Sujet + V (+\textbf{s} 3\up{e} pers. sg.) & Habitudes, vérités générales, emplois permanents & \eng{every day, usually, always, often, generally} \\
\hline
Present Continuous & Sujet + \textbf{am/is/are} + V-\textbf{ing} & Action en cours \emph{now}, situation temporaire & \eng{now, at the moment, currently, Look! Listen!} \\
\hline
Present Perfect & Sujet + \textbf{have/has} + V-\textbf{ed/pp} & Expérience de vie, durée depuis le passé (\eng{since/for}), résultat présent & \eng{already, just, yet, ever, never, since 2019, for 2 years} \\
\hline
Past Simple & Sujet + V-\textbf{ed/pp} (irréguliers) & Action terminée à un moment passé \emph{daté} & \eng{yesterday, last week/month/year, in 2019, ... ago} \\
\hline
\end{tabularx}
\end{center}

\begin{attention}[Piège majeur : Present Perfect vs Past Simple]
\eng{Since}, \eng{for}, \eng{already}, \eng{just}, \eng{never} $\rightarrow$ \textbf{Present Perfect} (lien avec le présent).
\eng{Yesterday}, \eng{last year}, \eng{in 2019}, \eng{two days ago} $\rightarrow$ \textbf{Past Simple} (coupé du présent).
\textbf{Jamais} : \eng{I have seen him yesterday.} $\rightarrow$ \eng{I \textbf{saw} him yesterday.}
\end{attention}

\courssub{S'auto-évaluer et remédier aux difficultés de grammaire}

\begin{methode}[Réviser les temps clés du chapitre (auto-évaluation)]
\begin{enumerate}
\item \textbf{Identifier l'indice de temps} dans la phrase : \eng{every day}, \eng{usually} $\rightarrow$ présent simple ; \eng{now}, \eng{at the moment} $\rightarrow$ présent continu ; \eng{already}, \eng{just}, \eng{since}, \eng{for} $\rightarrow$ present perfect ; \eng{yesterday}, \eng{last year}, \eng{ago} $\rightarrow$ prétérit simple.
\item \textbf{Vérifier la marque de la 3\up{e} personne} : \eng{he/she/it} + verbe en -\emph{s} au présent simple (\eng{She sells tomatoes.}).
\item \textbf{Ne jamais mélanger} : le present perfect est incompatible avec un moment passé précis. On dit \eng{I have sold my car.} mais \eng{I sold my car \textbf{yesterday}.}
\item \textbf{Se tester} : refaire les exercices sans regarder le corrigé, compter ses erreurs, puis relire uniquement la règle correspondant à l'erreur commise (remédiation ciblée).
\end{enumerate}
\end{methode}

\begin{exoresolu}[Grammar : choose the correct tense]
\textbf{Énoncé.} Complete with the correct form of the verb in brackets.

1. My father \dots\ (work) as a carpenter in Bafoussam. He \dots\ (have) his own workshop since 2015.

2. Look! The traders \dots\ (close) their shops because it is raining.

3. Last month, our cooperative \dots\ (sell) two tons of cocoa to a company in Douala.

4. I \dots\ (never / meet) a bank manager before.
\tcblower
\textbf{Solution détaillée.}

1. \eng{My father \textbf{works} as a carpenter in Bafoussam. He \textbf{has had} his own workshop since 2015.} — Première phrase : vérité permanente $\rightarrow$ présent simple avec -\emph{s} à la 3\up{e} personne. Deuxième phrase : l'indice \eng{since 2015} impose le present perfect (\eng{has} + participe passé \eng{had}).

2. \eng{Look! The traders \textbf{are closing} their shops because it is raining.} — L'indice \eng{Look!} signale une action en cours $\rightarrow$ présent continu : \eng{are} + \eng{closing} (attention : \eng{close} perd son \emph{e} muet devant -\emph{ing}).

3. \eng{Last month, our cooperative \textbf{sold} two tons of cocoa to a company in Douala.} — \eng{Last month} = moment passé terminé $\rightarrow$ prétérit simple ; \eng{sell} est irrégulier : \emph{sell -- sold -- sold}.

4. \eng{I \textbf{have never met} a bank manager before.} — \eng{never\ldots before} exprime l'expérience de vie $\rightarrow$ present perfect ; \eng{meet} est irrégulier : \emph{meet -- met -- met}.

\textbf{Conclusion} : la méthode « indice de temps $\rightarrow$ temps verbal » permet de résoudre chaque item sans hésitation.
\end{exoresolu}

\courssub{Consolider le vocabulaire de la vie économique}

\begin{methode}[Retenir et réutiliser le lexique d'un champ sémantique]
\begin{enumerate}
\item \textbf{Classer le vocabulaire en familles} : métiers (\eng{farmer, trader, teacher, driver, tailor, nurse}), lieux (\eng{market, bank, workshop, office, factory}), actions (\eng{to buy, to sell, to earn, to save, to borrow, to hire, to produce}), concepts (\eng{salary, profit, debt, unemployment, trade}).
\item \textbf{Apprendre les collocations} (mots qui vont ensemble) : \eng{to earn money}, \eng{to apply for a job}, \eng{to run a business}, \eng{to pay taxes}, \eng{to look for work}.
\item \textbf{Attention aux faux amis} : \eng{to earn} = gagner (de l'argent), pas « gagner » un match (\eng{to win}) ; \eng{a salary} = un salaire ; \eng{actually} = en fait, pas « actuellement » (\eng{currently}).
\item \textbf{Se tester par la traduction} : traduire 10 mots du français vers l'anglais, puis vérifier l'orthographe (\eng{business}, \eng{entrepreneur}, \eng{government}).
\end{enumerate}
\end{methode}

\begin{exoresolu}[Vocabulary : gap-fill and word families]
\textbf{Énoncé.} A. Complete the sentences with: \emph{salary, customers, profit, unemployed, loan}.

1. A bank can give you a \dots\ to start a small business.

2. Since the factory closed, many workers have been \dots .

3. A good shopkeeper is always polite to his \dots .

4. At the end of the month, she receives her \dots\ from the company.

5. After paying all his expenses, the trader made a \dots\ of 50,000 francs.

B. Translate into English: « Mon frère gagne sa vie comme chauffeur de taxi à Yaoundé. »
\tcblower
\textbf{Solution détaillée.}

A. 1. \eng{loan} — une banque \emph{prête} de l'argent ; \eng{a loan} = un prêt. 2. \eng{unemployed} — adjectif signifiant « au chômage » ; on dit \eng{to be unemployed} (jamais \eng{to have unemployment}). 3. \eng{customers} — les clients d'un commerçant ; ne pas confondre avec \eng{customs} (la douane). 4. \eng{salary} — le salaire mensuel ; \eng{wages} désigne plutôt une paie à l'heure ou à la semaine. 5. \eng{profit} — le bénéfice ; collocation : \eng{to make a profit} (son contraire : \eng{to make a loss}).

B. \eng{My brother earns his living as a taxi driver in Yaoundé.} — Justifications : « gagner sa vie » se traduit par la collocation \eng{to earn one's living} (et non \eng{to win his life}, calque du français) ; « comme » devant un métier = \eng{as} ; pas d'article devant le métier après \eng{as}.

\textbf{Conclusion} : les collocations et les faux amis sont les deux points où le correcteur vérifie la maîtrise réelle du lexique.
\end{exoresolu}

\courssub{Synthèse de vocabulaire : familles de mots et collocations}

\begin{propriete}[Tableau de dérivation et collocations clés (à connaître par cœur)]
\end{propriete}

\begin{center}
\begin{tabularx}{\linewidth}{|l|X|X|}
\hline
\rowcolor{popblueL}
\textbf{Nom (Noun)} & \textbf{Verbe (Verb)} & \textbf{Collocation forte (à apprendre par bloc)} \\
\hline
\eng{employment / unemployment} & \eng{to employ / to be unemployed} & \eng{to look for employment}, \eng{unemployment rate} \\
\hline
\eng{application} & \eng{to apply for} & \eng{to apply for a job / a post / a loan} \\
\hline
\eng{salary / wages} & \eng{to earn / to pay} & \eng{to earn a salary}, \eng{monthly salary} \\
\hline
\eng{profit / loss} & \eng{to make a profit / a loss} & \eng{to make a profit}, \eng{profit margin} \\
\hline
\eng{business} & \eng{to run a business / to do business} & \eng{small business}, \eng{business partner} \\
\hline
\eng{qualification} & \eng{to qualify} & \eng{academic qualifications}, \eng{to qualify as a teacher} \\
\hline
\eng{trade} & \eng{to trade} & \eng{international trade}, \eng{trade fair} \\
\hline
\eng{loan / debt} & \eng{to borrow / to lend} & \eng{to take out a loan}, \eng{to be in debt} \\
\hline
\eng{savings} & \eng{to save} & \eng{life savings}, \eng{to save money} \\
\hline
\eng{customer / client} & \eng{to serve / to attract} & \eng{regular customer}, \eng{customer service} \\
\hline
\end{tabularx}
\end{center}

\courssub{Réussir les épreuves de compréhension de type examen}

\begin{methode}[Traiter un texte de compréhension écrite (reading)]
\begin{enumerate}
\item \textbf{Lire d'abord les questions}, pas le texte : ainsi vous savez quoi chercher.
\item \textbf{Première lecture rapide} (skimming) pour le sens général : qui ? où ? quel sujet ?
\item \textbf{Lecture ciblée} (scanning) : repérer dans le texte le mot-clé de chaque question ; la réponse se trouve presque toujours dans la phrase qui l'entoure.
\item \textbf{Répondre avec ses propres mots} en phrase complète, sans recopier un bloc entier du texte.
\item \textbf{Pour les questions de vocabulaire} (\eng{Find in the text a word meaning\ldots}), chercher un synonyme dans le paragraphe indiqué et vérifier la nature grammaticale demandée (nom, verbe, adjectif).
\item \textbf{Vérifier} que chaque réponse est en anglais correct : sujet + verbe conjugué.
\end{enumerate}
\end{methode}

\begin{exoresolu}[Reading comprehension (exam-type)]
\textbf{Énoncé.} Read the text and answer the questions.

\begin{textebox}[Extrait adapté d'un article économique camerounais]
“Mama Fotso left her village near Bafoussam ten years ago. She first sold vegetables in a small market. With her savings, she opened a restaurant in Douala in 2019. Today she employs six people and supplies three schools with meals. ‘Hard work and honesty are the secrets of my success,' she says. She now wants to open a second restaurant in Buea.”
\end{textebox}

1. Where did Mama Fotso start her business? (2 marks)

2. How did she get the money to open her restaurant? (2 marks)

3. Find in the text a word meaning “workers”. (2 marks)

4. What are her plans for the future? (2 marks)

5. True or False? Justify: “Mama Fotso sells food to schools.” (2 marks)
\tcblower
\textbf{Solution détaillée.}

1. \eng{She started her business in a small market near Bafoussam, where she sold vegetables.} — Réponse en phrase complète ; prétérit simple car l'action est passée et datée.

2. \eng{She used her savings, that is the money she had saved from selling vegetables.} — Le texte dit \eng{with her savings} ; on reformule au lieu de recopier.

3. \eng{The text does not contain the noun “workers” or “employees”. The phrase “six people” refers to them. Acceptable answer: “people” (in the phrase “she employs six people”).} — \textbf{Note pédagogique} : en examen, si le nom exact n'est pas dans le texte, on cite le groupe nominal qui y fait référence (\eng{six people}) ou le verbe \eng{employs} si la nature n'est pas imposée. Ici, « a word meaning workers » attend un nom ; \eng{people} est la meilleure réponse contextuelle.

4. \eng{She plans to open a second restaurant in Buea.} — Reformulation de \eng{she wants to open\ldots} ; \eng{to plan to do something} = projeter de faire quelque chose.

5. \eng{True. The text says she “supplies three schools with meals”, which means she sells food to schools.} — Une justification doit toujours citer le texte entre guillemets.

\textbf{Conclusion} : chaque réponse est une phrase complète, justifiée par le texte, sans recopie excessive.
\end{exoresolu}

\courssub{Compréhension de l'oral (Listening) : script et questions}

\begin{methode}[Stratégies pour l'écoute (listening) en examen]
\begin{enumerate}
\item \textbf{Lire les questions} pendant le temps de préparation (souligner mots-clés : \emph{who, where, how much, why}).
\item \textbf{Première écoute} : compréhension globale (sujet, locuteurs, ton).
\item \textbf{Deuxième écoute} : prise de notes ciblées (chiffres, noms, verbes d'action).
\item \textbf{Troisième écoute} : vérification et rédaction des réponses en phrases complètes.
\item \textbf{Ne pas paniquer} si un mot échappe : le contexte suffit souvent.
\end{enumerate}
\end{methode}

\begin{textebox}[Script d'écoute — Enregistrement simulé : Job Fair at Yaoundé Congress Centre]
\eng{Presenter:} Good morning. Welcome to the National Youth Employment Fair here at the Yaoundé Congress Centre. Today, over fifty companies are recruiting. We speak with Mr. Atangana, HR manager at a cocoa processing factory in Douala.

\eng{Journalist:} Mr. Atangana, what profiles are you looking for?

\eng{Mr. Atangana:} We need young technicians. We have ten vacancies for machine operators and five for quality controllers. Candidates must have a BTS or HND in food technology or mechanical engineering. Experience is a plus, but we accept fresh graduates. We offer a starting salary of 180,000 FCFA per month, plus transport and health insurance.

\eng{Journalist:} How can they apply?

\eng{Mr. Atangana:} They can drop their CV and cover letter at our stand, number 14, before 4 p.m. today. Or send them by email to recruitment@cocoa-pro.cm. Interviews start next Monday.

\eng{Presenter:} Thank you, Mr. Atangana. Good luck to all job seekers!
\end{textebox}

\begin{exercice}[Listening Comprehension Questions]
Answer in complete sentences.
1. Where does the job fair take place?
2. How many vacancies are there for machine operators?
3. What diploma is required for the jobs?
4. What is the starting salary offered?
5. What are the two ways to apply mentioned by Mr. Atangana?
6. When do the interviews start?
\end{exercice}

\begin{corrige}[Exercice Listening]
1. \eng{The job fair takes place at the Yaoundé Congress Centre.}
2. \eng{There are ten vacancies for machine operators.}
3. \eng{Candidates must have a BTS or HND in food technology or mechanical engineering.}
4. \eng{The starting salary is 180,000 FCFA per month.}
5. \eng{They can drop their CV and cover letter at stand 14 or send them by email to recruitment@cocoa-pro.cm.}
6. \eng{The interviews start next Monday.}
\end{corrige}

\courssub{Consolider l'expression orale et écrite}

\begin{methode}[Parler de son projet professionnel (speaking)]
\begin{enumerate}
\item \textbf{Structurer en trois temps} : présentation (\eng{My name is\ldots I am a student at\ldots}), projet (\eng{In the future, I would like to become\ldots / I am going to\ldots}), justification (\eng{because I am good at\ldots / because this job is useful for my community}).
\item \textbf{Utiliser les bons futurs} : \eng{going to} pour un projet déjà décidé (\eng{I am going to study accounting.}), \eng{will} pour une prédiction ou une promesse (\eng{I will work hard.}).
\item \textbf{Après \eng{to be good at}, \eng{to be interested in}, \eng{to dream of}} : verbe en -\emph{ing} : \eng{I am interested in trading.}
\item \textbf{Soigner la prononciation} des mots du chapitre : \eng{business} se prononce « biz-nèss », \eng{entrepreneur} « on-trö-prö-neur », \eng{economy} « i-ko-no-mi ».
\item \textbf{Conclure} avec une phrase de bilan : \eng{That is why this occupation is my dream job.}
\end{enumerate}
\end{methode}

\begin{experience}[Jeu de rôle : Entretien d'embauche (Job Interview)]
\textbf{Consigne.} Par binôme. L'un est le recruteur (manager d'une entreprise agro-alimentaire à Douala), l'autre est le candidat (titulaire d'un BTS en transformation alimentaire).
\textbf{Recruteur :} Accueillir, demander le parcours (études, stages), les qualités (\eng{hard-working, punctual, team spirit}), la prétention salariale, conclure.
\textbf{Candidat :} Se présenter, décrire son expérience (present perfect), exprimer sa motivation (\eng{interested in}, \eng{good at}), poser une question sur la formation continue.
\textbf{Durée :} 3 minutes. Changement de rôles. L'enseignant note : utilisation du vocabulaire chapitre, justesse des temps (present perfect pour l'expérience), fluidité.
\end{experience}

\begin{exoresolu}[Writing : a short composition]
\textbf{Énoncé.} Write a composition of about 100 words on: “The occupation I would like to have in the future and why.”
\tcblower
\textbf{Solution détaillée.}

\textbf{Plan} : 1) nom du métier ; 2) deux raisons (goût personnel, utilité sociale) ; 3) les moyens pour y arriver ; 4) conclusion.

\textbf{Réponse rédigée :}

\eng{In the future, I would like to become an agricultural engineer. First, I have always been interested in farming because my grandparents are cocoa farmers in the Centre Region. Secondly, this occupation is very useful: Cameroon needs experts who can help farmers produce more food and earn more money. To achieve my goal, I am going to study science seriously and then attend a university of agriculture. I know the road will be long, but I will work hard. One day, I hope to create jobs for young people in my village. That is why agricultural engineering is my dream occupation.}

(99 mots.) \textbf{Justifications grammaticales} : \eng{I would like to become} = conditionnel de politesse pour exprimer un souhait ; \eng{interested \textbf{in} farming} = préposition \eng{in} + -\emph{ing} ; \eng{I am \textbf{going to study}} = projet décidé ; \eng{I \textbf{will work} hard} = promesse ; \eng{who can help} = pronom relatif pour les personnes.

\textbf{Conclusion} : le texte respecte le nombre de mots, mobilise le lexique économique (\eng{occupation, farmers, earn money, create jobs}) et utilise trois futurs correctement différenciés.
\end{exoresolu}

\courssub{Évaluation, auto-évaluation et remédiation}

\begin{exercice}[Examen blanc — Chapter 2 (Durée : 45 min)]
\textbf{Section A : Grammar (10 marks)} \\
Complete with the correct tense (Present Simple, Present Continuous, Present Perfect, Past Simple).
1. My uncle \_\_\_\_\_\_ (run) a small business in Bamenda. He \_\_\_\_\_\_ (start) it in 2010.
2. Look! The customers \_\_\_\_\_\_ (queue) at the bank counter.
3. I \_\_\_\_\_\_ (never / be) unemployed because I always \_\_\_\_\_\_ (find) work quickly.
4. Last year, the government \_\_\_\_\_\_ (create) 5,000 new jobs in the agricultural sector.
5. How long \_\_\_\_\_\_ you \_\_\_\_\_\_ (know) the manager of this company?

\textbf{Section B : Vocabulary (10 marks)} \\
A. Match the words (1--5) with their definitions (a--e).
1. Loan \quad 2. Profit \quad 3. Unemployed \quad 4. Salary \quad 5. Customer
a) Money you borrow from a bank. \quad b) A person who buys goods. \quad c) Money earned after expenses. \quad d) Without a job. \quad e) Monthly pay.

B. Translate into English.
6. « Elle a postulé pour un poste de secrétaire. »
7. « Nous gagnons notre vie en vendant du cacao. »
8. « Actuellement, il cherche du travail. »
9. « L'entreprise a fait un bénéfice l'an dernier. »
10. « Ce jeune entrepreneur dirige sa propre affaire. »

\textbf{Section C : Writing (10 marks)} \\
Choose ONE topic. Write 100--120 words.
\textbf{Topic 1:} You are the manager of a cooperative in Buea. Write a formal letter to the Minister of Commerce asking for a loan to buy a new truck. Explain why you need it and how you will repay.
\textbf{Topic 2:} “Starting a small business is better than looking for a salaried job.” Do you agree? Give two reasons.
\end{exercice}

\begin{corrige}[Examen blanc — Chapter 2]
\textbf{Section A : Grammar}
1. \eng{runs} (Present Simple, 3\up{e} pers. sg., habitude) / \eng{started} (Past Simple, date passée \eng{in 2010}).
2. \eng{are queuing} (Present Continuous, \eng{Look!} action en cours).
3. \eng{have never been} (Present Perfect, \eng{never} expérience de vie) / \eng{find} (Present Simple, habitude \eng{always}).
4. \eng{created} (Past Simple, \eng{Last year} moment daté).
5. \eng{have / known} (Present Perfect, \eng{How long} + \eng{know} verbe d'état = durée depuis le passé).

\textbf{Section B : Vocabulary}
A. 1-a, 2-c, 3-d, 4-e, 5-b.
B. 6. \eng{She applied for the post of secretary.} (ou \eng{a secretary position})
7. \eng{We earn our living by selling cocoa.} (ou \eng{selling cocoa})
8. \eng{Currently, he is looking for a job.} (Attention : \eng{actually} = en fait ; \eng{currently} = actuellement)
9. \eng{The company made a profit last year.} (\eng{make a profit}, Past Simple \eng{last year})
10. \eng{This young entrepreneur runs his own business.} (\eng{run a business}, \eng{entrepreneur} orthographe)

\textbf{Section C : Writing (Grilles indicatives)}
\textbf{Topic 1 (Lettre formelle) :} Structure (en-tête, objet, formules \eng{Dear Sir/Madam}, \eng{Yours faithfully}) / Contenu (besoin camion, augmentation production, remboursement via bénéfices) / Langue (Present Perfect pour expérience coopérative, Future \eng{will repay}, conditionnel \eng{would be grateful}, vocabulaire \eng{loan, truck, cooperative, repay, profit}).
\textbf{Topic 2 (Opinion) :} Structure (Intro avis, 2 arguments, conclusion) / Arguments ex: autonomie, création emplois, risque échec vs sécurité salaire, congés payés / Langue (Connecteurs \eng{Firstly, However, In conclusion}, modaux \eng{can, should}, vocabulaire \eng{business, salary, risk, independence, employee}).
\end{corrige}

\begin{aretenir}[Grille d'auto-évaluation — Chapter 2]
Cochez \eng{✓} si maîtrisé, \eng{~} si à réviser, \eng{✗} si non acquis. Pour chaque \eng{~} ou \eng{✗}, relisez la fiche correspondante et refaites l'exercice type.

\begin{center}
\begin{tabularx}{\linewidth}{|X|c|c|c|}
\hline
\rowcolor{popblueL}
\textbf{Compétence / Notion} & \eng{✓} & \eng{~} & \eng{✗} \\
\hline
Vocabulaire : nommer 10 métiers, 5 lieux, 5 actions, 5 concepts économiques & & & \\
\hline
Collocations : \eng{apply for, earn a living, make a profit, run a business, look for work} & & & \\
\hline
Faux amis : \eng{earn/win, actually/currently, salary/wages, customer/customs} & & & \\
\hline
Grammaire : choisir le bon temps (indices \eng{since/for} vs \eng{yesterday/last}) & & & \\
\hline
Grammaire : 3\up{e} personne \eng{-s} présent simple (affirmatif, négatif, interrogatif) & & & \\
\hline
Grammaire : \eng{going to} (projet) vs \eng{will} (prédiction/promesse) & & & \\
\hline
Reading : technique \emph{skimming/scanning}, répondre phrase complète, justifier True/False & & & \\
\hline
Listening : noter chiffres, noms, diplômes, consignes (stand, email, deadline) & & & \\
\hline
Speaking : présenter projet pro (3 temps), prononcer \eng{business, entrepreneur, economy} & & & \\
\hline
Writing : lettre candidature (formule, 80-100 mots), composition avis (structure, connecteurs) & & & \\
\hline
\end{tabularx}
\end{center}
\end{aretenir}

\begin{savaistu}[Le secteur informel au Cameroun : moteur de l'emploi]
Au Cameroun, près de \textbf{90\% de la population active} travaille dans le secteur informel (source : INS, enquêtes récentes). Ce secteur comprend les \eng{street vendors} (vendeurs à la sauvette), les \eng{moto-taxi drivers} (\eng{bendskin} drivers), les artisans (\eng{tailors, carpenters, mechanics}), les petits commerçants de marché (\eng{buyam-sellam}). Bien que souvent précaires (pas de \eng{social security}, pas de \eng{paid leave}), ces activités permettent à des millions de jeunes de \eng{earn a living} en attendant un emploi formel. Le gouvernement encourage la \eng{formalisation} via la carte d'identification unique et l'accès au microcrédit (\eng{microfinance institutions} comme les MC\textsuperscript{2}). Connaître ce vocabulaire (\eng{informal sector, self-employed, hustle, daily income}) est un atout pour l'oral et l'écrit du Baccalauréat.
\end{savaistu}

\begin{attention}[Les 5 erreurs qui coûtent des points au Bac]
\begin{enumerate}
\item \textbf{Oublier le -\emph{s} de la 3\up{e} personne} au présent simple : écrire \eng{She work in a bank} au lieu de \eng{She \textbf{works} in a bank}. C'est la faute la plus sanctionnée.
\item \textbf{Calquer le français} : \eng{to win money} pour « gagner de l'argent » (correct : \eng{to \textbf{earn} money}) ; \eng{to make his studies} pour « faire ses études » (correct : \eng{to \textbf{study}} ou \eng{to \textbf{do} one's studies}) ; \eng{actually} pour « actuellement » (correct : \eng{currently}).
\item \textbf{Mélanger present perfect et passé daté} : \eng{I have seen him yesterday} est impossible. Moment précis (\eng{yesterday, last year, in 2019}) $\rightarrow$ prétérit simple : \eng{I \textbf{saw} him yesterday.}
\item \textbf{Se tromper de préposition} : on dit \eng{to apply \textbf{for} a job}, \eng{to depend \textbf{on}}, \eng{to be interested \textbf{in}}, \eng{to be good \textbf{at}}, \eng{to look \textbf{for} work}. Apprenez chaque verbe \emph{avec} sa préposition.
\item \textbf{Fautes d'orthographe dans le lexique économique} : \eng{buisness} au lieu de \eng{\textbf{business}}, \eng{goverment} au lieu de \eng{\textbf{government}}, \eng{entreprenor} au lieu de \eng{\textbf{entrepreneur}}, \eng{unemployement} au lieu de \eng{\textbf{unemployment}}. Une dictée de dix mots du chapitre avant l'examen suffit à les éviter.
\end{enumerate}
\end{attention}

\newpage

\coursec{Exercices}

\courssub{Je vérifie mes connaissances}

\begin{exercice}[QCM de grammaire : la vie économique]
Entoure la bonne réponse. Un seul choix est correct.
\begin{enumerate}
\item My father \textbf{---} in a bank in Douala. \quad a) work \quad b) works \quad c) is work \quad d) working
\item Listen! The manager \textbf{---} to the workers right now. \quad a) speaks \quad b) is speaking \quad c) speak \quad d) speaking
\item A doctor earns \textbf{---} money than a shopkeeper. \quad a) much \quad b) more \quad c) most \quad d) many
\item This is the \textbf{---} factory in the whole region. \quad a) bigger \quad b) big \quad c) biggest \quad d) most big
\item There are too \textbf{---} unemployed young people in our town. \quad a) much \quad b) many \quad c) a lot \quad d) little
\item She has \textbf{---} experience in teaching. \quad a) many \quad b) much \quad c) a \quad d) few
\item The cocoa harvest this year is \textbf{---} than last year. \quad a) good \quad b) better \quad c) best \quad d) more good
\item My uncle \textbf{---} his own business since 2015. \quad a) has \quad b) have \quad c) is having \quad d) having
\item \textbf{---} information do you need about the job? \quad a) How many \quad b) How much \quad c) How long \quad d) How far
\item A taxi driver works \textbf{---} hours than an office clerk. \quad a) long \quad b) longer \quad c) longest \quad d) more long
\end{enumerate}
\end{exercice}

\begin{exercice}[Vrai ou faux, justifié]
Lis les phrases suivantes. Écris \textbf{True} ou \textbf{False}, puis justifie ta réponse en une courte phrase en anglais.
\begin{enumerate}
\item “Money” is a countable noun in English.
\item We use the present continuous for actions happening at the moment of speaking.
\item The comparative of “good” is “gooder”.
\item We say “much people” to talk about a large number of persons.
\item The superlative of “far” can be “the farthest”.
\item An employer is a person who looks for a job.
\end{enumerate}
\end{exercice}

\begin{exercice}[Vocabulaire : métiers et vie économique]
Complète chaque phrase avec le mot anglais qui convient, formé à partir de l'indice entre parenthèses.
\begin{enumerate}
\item A person who buys and sells goods in a market is a \textbf{---} (to trade).
\item The money a worker receives every month is his \textbf{---} (to pay, nom).
\item A person who has no job is \textbf{---} (employ, préfixe négatif).
\item The person who gives you a job is your \textbf{---} (to employ).
\item A meeting between an employer and a candidate is a job \textbf{---} (to view, préfixe).
\item Someone who starts and runs his own business is an \textbf{---} (entreprise).
\end{enumerate}
\end{exercice}

\begin{exercice}[Conjugaison à compléter]
Mets le verbe entre parenthèses au temps qui convient (present simple ou present continuous).
\begin{enumerate}
\item Every morning, traders \textbf{---} (open) their shops at 7 a.m. in Mokolo market.
\item Look! The bendskin driver \textbf{---} (carry) three passengers!
\item My mother \textbf{---} (sell) vegetables at the market, but today she \textbf{---} (rest) at home.
\item The government \textbf{---} (build) new roads this year to help farmers.
\item Banks in Cameroon usually \textbf{---} (close) at 3.30 p.m.
\item I \textbf{---} (not / understand) this economic report. Can you explain it?
\end{enumerate}
\end{exercice}

\courssub{Je m'entraîne}

\begin{exercice}[Traduction]
Traduis les phrases suivantes en anglais correct.
\begin{enumerate}
\item Mon père gagne plus d'argent que mon oncle.
\item Il y a trop de chômeurs dans notre ville.
\item Elle cherche du travail depuis six mois.
\item Le cacao est le produit d'exportation le plus important de la région.
\item En ce moment, les ouvriers construisent un nouveau marché à Bafoussam.
\end{enumerate}
\end{exercice}

\begin{exercice}[Reformulation : comparatif et superlatif]
Transforme chaque phrase comme indiqué, sans changer le sens.
\begin{enumerate}
\item No other bank in Yaoundé is bigger than Afriland. (Use the superlative.)
\item A nurse is busier than a secretary. (Use “less … than”.)
\item This is the most expensive shop in town. (Use the comparative with “than”.)
\item Farming is harder than office work. (Use “not as … as”.)
\item No job is more dangerous than mining. (Use the comparative.)
\end{enumerate}
\end{exercice}

\begin{exercice}[Texte à trous : lettre de candidature]
Complète la lettre suivante avec les mots de la banque ci-dessous. Chaque mot ne peut être utilisé qu'une seule fois.

\textbf{Banque de mots :} salary — apply — experience — interview — employer — position — skills — advertisement

\begin{textebox}[A letter of application]
Dear Sir,

I am writing to (1) \textbf{---} for the (2) \textbf{---} of shop assistant which I saw in your (3) \textbf{---} in the \emph{Cameroon Tribune} of Monday.

I have two years of (4) \textbf{---} in a supermarket in Douala. I believe my (5) \textbf{---} in sales and customer service would be useful to your company. I would be happy to attend an (6) \textbf{---} at any time convenient to you.

I hope the (7) \textbf{---} and working conditions can be discussed during our meeting. I am sure I can be a good worker for any serious (8) \textbf{---}.

Yours faithfully,

Amina Bello
\end{textebox}
\end{exercice}

\begin{exercice}[Appariements : le mot et sa définition]
Associe chaque mot de la colonne A à sa définition de la colonne B. Écris les paires (ex. 1–c).

\begin{center}
\begin{tabularx}{\linewidth}{|l|X|l|X|}
\hline
\rowcolor{popblueL} \multicolumn{2}{|c|}{\textbf{Column A}} & \multicolumn{2}{c|}{\textbf{Column B}} \\
\hline
1 & a customer & a & a person who works for a company \\
\hline
2 & an employee & b & money paid for work, usually every month \\
\hline
3 & a salary & c & a person who buys goods or services \\
\hline
4 & a farmer & d & a place where goods are made in large quantities \\
\hline
5 & a factory & e & a person who grows crops or keeps animals \\
\hline
6 & unemployment & f & the situation of people who have no job \\
\hline
\end{tabularx}
\end{center}
\end{exercice}

\courssub{Je me prépare au Bac}

\begin{exercice}[Compréhension écrite type Bac]
Lis le texte suivant, puis réponds aux questions.

\begin{textebox}[A young entrepreneur from Bamenda]
Ngong Blessing is twenty-three years old. She lives in Bamenda, in the North-West Region of Cameroon. After finishing secondary school, she could not find a job, so she decided to create her own business. With a small loan from her aunt, she bought a sewing machine and started making clothes for women in her neighbourhood.

At the beginning, business was very difficult. “I worked longer hours than any of my friends, and I earned less money,” she remembers. But Blessing did not give up. Little by little, her clothes became famous because they were cheaper and more beautiful than those in the big shops. Today, she employs four young girls and is training two apprentices.

Blessing believes that young Cameroonians should not wait for the government to give them jobs. “The best way to fight unemployment is to create your own work,” she says. She is now planning to open a second workshop in Bafoussam, and she hopes to export her clothes to Nigeria one day. Her story shows that with courage and hard work, young people can build a better future for themselves and for their country.
\end{textebox}

\textbf{A. True or False. Justify each answer with a quotation from the text.} (5 marks)
\begin{enumerate}
\item Blessing created her business because she had a lot of money.
\item At the start, she earned more money than her friends.
\item Her clothes were successful because of their price and quality.
\item Blessing thinks young people should wait for government jobs.
\item She wants to sell her clothes outside Cameroon in the future.
\end{enumerate}

\textbf{B. Answer these questions in your own words.} (4 marks)
\begin{enumerate}
\item Where does Blessing live?
\item Who gave her the money to start her business?
\item How many people work with her today?
\item What lesson can young people learn from Blessing's story?
\end{enumerate}

\textbf{C. Language work.} (3 marks)
\begin{enumerate}
\item Find in the text the comparative of “long” and the superlative of “good”.
\item Find in the text a synonym of “well-known”.
\item Turn into the negative form: “She employs four young girls.”
\end{enumerate}
\end{exercice}

\begin{exercice}[Traduction type Bac]
Traduis en anglais le passage suivant. (5 marks)

\begin{textebox}[Texte à traduire]
Mon frère est chauffeur de taxi à Douala. Il travaille plus dur que beaucoup de gens, mais il gagne peu d'argent. En ce moment, il économise pour acheter sa propre voiture. Je pense que c'est le métier le plus difficile de notre famille, mais il ne se plaint jamais.
\end{textebox}
\end{exercice}

\begin{exercice}[Expression écrite type Bac]
\textbf{Topic :} Write a paragraph of about \textbf{120 words} on the following subject:

\begin{center}
\emph{“The job I would like to do in the future, and why.”}
\end{center}

In your paragraph, you should:
\begin{itemize}
\item name the job and say where you would like to do it (in Cameroon or abroad);
\item give at least \textbf{two reasons} for your choice;
\item use at least \textbf{one comparative} and \textbf{one superlative};
\item use link words such as \eng{first}, \eng{moreover}, \eng{however}, \eng{in conclusion}.
\end{itemize}

\textbf{Grille d'évaluation (10 marks) :}
\begin{center}
\begin{tabularx}{\linewidth}{|X|c|}
\hline
\rowcolor{popblueL} Critère & Points \\
\hline
Contenu : le sujet est traité, les deux raisons sont données & 3 \\
\hline
Vocabulaire : lexique des métiers et de la vie économique & 2 \\
\hline
Grammaire : temps corrects, comparatif et superlatif & 3 \\
\hline
Connecteurs et organisation du paragraphe & 2 \\
\hline
\end{tabularx}
\end{center}
\end{exercice}

\newpage

\coursec{Corrigés des exercices}

\begin{corrige}[Exercice 1 — QCM de grammaire : la vie économique]
\begin{enumerate}
\item \textbf{b) works} — \eng{My father works in a bank in Douala.} Present simple, 3\textsuperscript{e} personne du singulier : le verbe prend un \textbf{-s}.
\item \textbf{b) is speaking} — \eng{Listen! The manager is speaking to the workers right now.} Le marqueur \eng{Listen!} et \eng{right now} imposent le present continuous (be + V-ing).
\item \textbf{b) more} — \eng{A doctor earns more money than a shopkeeper.} Comparatif de supériorité devant un nom : \eng{more + nom + than}.
\item \textbf{c) biggest} — \eng{This is the biggest factory in the whole region.} Superlatif d'un adjectif court : \textbf{the + adj + -est} ; \eng{big} double sa consonne finale (big--bigger--biggest).
\item \textbf{b) many} — \eng{There are too many unemployed young people in our town.} \eng{People} est comptable pluriel : on utilise \eng{many} (et non \eng{much}).
\item \textbf{b) much} — \eng{She has much experience in teaching.} \eng{Experience} (expérience professionnelle) est ici indénombrable : \eng{much}.
\item \textbf{b) better} — \eng{The cocoa harvest this year is better than last year.} Comparatif irrégulier : good--better--best. Jamais \eng{more good}.
\item \textbf{a) has} — \eng{My uncle has his own business since 2015.} 3\textsuperscript{e} personne du singulier : \eng{has}. (On accepte aussi l'idée qu'avec \eng{since 2015}, le present perfect \eng{has had} serait idéal, mais parmi les choix proposés, seul \eng{has} est grammatical.)
\item \textbf{b) How much} — \eng{How much information do you need about the job?} \eng{Information} est indénombrable : \eng{How much} (et non \eng{How many}).
\item \textbf{b) longer} — \eng{A taxi driver works longer hours than an office clerk.} Comparatif d'un adjectif court : \eng{longer than}. Jamais \eng{more long}.
\end{enumerate}
\textbf{Points de vigilance :} ne pas oublier le \textbf{-s} de la 3\textsuperscript{e} personne ; distinguer \eng{much} (indénombrable) et \eng{many} (dénombrable) ; mémoriser les irréguliers good--better--best et far--further/farther--furthest/farthest.
\end{corrige}

\begin{corrige}[Exercice 2 — Vrai ou faux, justifié]
\begin{enumerate}
\item \textbf{False.} \eng{“Money” is an uncountable noun; we say “much money”, not “many money”.} (« Money » est indénombrable en anglais.)
\item \textbf{True.} \eng{We use the present continuous for actions happening now, for example: “The manager is speaking.”}
\item \textbf{False.} \eng{The comparative of “good” is irregular: “better”.} (Jamais \eng{gooder}.)
\item \textbf{False.} \eng{We say “many people” because “people” is countable.} (On dit \eng{many people}, jamais \eng{much people}.)
\item \textbf{True.} \eng{“Far” has two superlatives: “the farthest” (distance) and “the furthest”.}
\item \textbf{False.} \eng{An employer gives jobs; a job seeker (or an unemployed person) looks for a job.} (Attention au faux ami : \eng{employer} = employeur, celui qui embauche.)
\end{enumerate}
\end{corrige}

\begin{corrige}[Exercice 3 — Vocabulaire : métiers et vie économique]
\begin{enumerate}
\item \textbf{trader} — \eng{A person who buys and sells goods in a market is a trader.} (verbe \eng{to trade} + suffixe \textbf{-er}).
\item \textbf{salary} — \eng{The money a worker receives every month is his salary.} (nom dérivé de la famille de \eng{to pay} ; ne pas confondre avec \eng{wages}, payés à l'heure ou à la semaine).
\item \textbf{unemployed} — \eng{A person who has no job is unemployed.} (préfixe négatif \textbf{un-} + participe passé employé comme adjectif).
\item \textbf{employer} — \eng{The person who gives you a job is your employer.} (attention : \eng{employee} = employé, celui qui reçoit le travail).
\item \textbf{interview} — \eng{A meeting between an employer and a candidate is a job interview.} (préfixe \textbf{inter-} + \eng{view}).
\item \textbf{entrepreneur} — \eng{Someone who starts and runs his own business is an entrepreneur.} (mot emprunté au français ; prononciation « on-trö-prö-neur »).
\end{enumerate}
\textbf{Points de vigilance :} bien mémoriser les suffixes de personnes \textbf{-er / -or / -ee} (\eng{employer / employee}, \eng{trainer / trainee}) et l'article \eng{an} devant voyelle : \eng{an entrepreneur}, \eng{an interview}.
\end{corrige}

\begin{corrige}[Exercice 4 — Conjugaison à compléter]
\begin{enumerate}
\item \textbf{open} — \eng{Every morning, traders open their shops at 7 a.m. in Mokolo market.} Habitude (\eng{every morning}) : present simple ; sujet pluriel, pas de \textbf{-s}.
\item \textbf{is carrying} — \eng{Look! The bendskin driver is carrying three passengers!} Action en cours (\eng{Look!}) : present continuous ; sujet singulier : \eng{is}.
\item \textbf{sells — is resting} — \eng{My mother sells vegetables at the market, but today she is resting at home.} Habitude (present simple, 3\textsuperscript{e} pers. : \eng{sells}) contre situation temporaire d'aujourd'hui (present continuous : \eng{is resting}).
\item \textbf{is building} — \eng{The government is building new roads this year to help farmers.} Projet en cours sur une période limitée (\eng{this year}) : present continuous.
\item \textbf{close} — \eng{Banks in Cameroon usually close at 3.30 p.m.} Horaire habituel (\eng{usually}) : present simple ; sujet pluriel.
\item \textbf{do not understand} (ou \eng{don't understand}) — \eng{I do not understand this economic report. Can you explain it?} \eng{Understand} est un verbe d'état : il ne se met presque jamais au present continuous. Négation : \eng{do not + base verbale}.
\end{enumerate}
\textbf{Points de vigilance :} les verbes d'état (\eng{understand, know, want, like, believe}) refusent la forme en \textbf{-ing} ; repérer les marqueurs : \eng{every morning, usually} (simple) vs \eng{look, listen, now, today} (continuous).
\end{corrige}

\begin{corrige}[Exercice 5 — Traduction]
\begin{enumerate}
\item \eng{My father earns more money than my uncle.} (\eng{to earn money} = gagner de l'argent ; comparatif \eng{more ... than}.)
\item \eng{There are too many unemployed people in our town.} (« Il y a » = \eng{there are} ; « trop de » + nom pluriel = \eng{too many}.)
\item \eng{She has been looking for a job for six months.} (Durée commencée dans le passé et qui continue : present perfect continuous + \eng{for} ; « chercher » = \eng{to look for}, ne jamais oublier la préposition \eng{for}.)
\item \eng{Cocoa is the most important export product in the region.} (Superlatif d'un adjectif long : \eng{the most important} ; « produit d'exportation » = \eng{export product}.)
\item \eng{At the moment, the workers are building a new market in Bafoussam.} (« En ce moment » = \eng{at the moment} : present continuous.)
\end{enumerate}
\textbf{Points de vigilance :} ne pas traduire « il y a » par \eng{it has} ; « depuis six mois » avec un verbe d'action se rend par \eng{for six months} au present perfect ; « peu d'argent » = \eng{little money} (indénombrable).
\end{corrige}

\begin{corrige}[Exercice 6 — Reformulation : comparatif et superlatif]
\begin{enumerate}
\item \eng{Afriland is the biggest bank in Yaoundé.} (Le superlatif reprend l'adjectif du comparatif : bigger $\rightarrow$ the biggest.)
\item \eng{A secretary is less busy than a nurse.} (On inverse les deux termes avec \eng{less ... than} ; \eng{busier} devient \eng{less busy}.)
\item \eng{This shop is more expensive than any other shop in town.} (Du superlatif au comparatif : \eng{more ... than any other + nom singulier}.)
\item \eng{Office work is not as hard as farming.} (On inverse les termes ; structure \eng{not as + adjectif + as}.)
\item \eng{Mining is more dangerous than any other job.} (Ou : \eng{Mining is the most dangerous job.} — ici le comparatif demandé : \eng{more dangerous than any other job}.)
\end{enumerate}
\textbf{Points de vigilance :} avec \eng{less ... than} et \eng{not as ... as}, il faut inverser l'ordre des éléments pour garder le sens ; après \eng{any other}, le nom est au singulier.
\end{corrige}

\begin{corrige}[Exercice 7 — Texte à trous : lettre de candidature]
\begin{enumerate}
\item \textbf{apply} — \eng{I am writing to apply for the position...} (\eng{to apply for a job} = postuler ; la préposition \eng{for} est obligatoire.)
\item \textbf{position} — \eng{...for the position of shop assistant...} (\eng{position} = poste.)
\item \textbf{advertisement} — \eng{...which I saw in your advertisement in the \emph{Cameroon Tribune}...} (\eng{advertisement} = annonce ; souvent abrégé \eng{ad / advert}.)
\item \textbf{experience} — \eng{I have two years of experience in a supermarket in Douala.} (Indénombrable ici : pas de \textbf{-s}.)
\item \textbf{skills} — \eng{...my skills in sales and customer service...} (\eng{skills} = compétences.)
\item \textbf{interview} — \eng{I would be happy to attend an interview...} (\eng{to attend an interview} = se présenter à un entretien.)
\item \textbf{salary} — \eng{I hope the salary and working conditions can be discussed...}
\item \textbf{employer} — \eng{...a good worker for any serious employer.}
\end{enumerate}
\textbf{Points de vigilance :} dans une lettre formelle, ouverture \eng{Dear Sir,} et clôture \eng{Yours faithfully,} ; ne pas écrire \eng{experiences} au sens d'expérience professionnelle.
\end{corrige}

\begin{corrige}[Exercice 8 — Appariements : le mot et sa définition]
\begin{center}
\begin{tabular}{|c|c|l|}
\hline
\rowcolor{popblueL} Paire & Mot & Définition \\
\hline
1--c & a customer & a person who buys goods or services \\
\hline
2--a & an employee & a person who works for a company \\
\hline
3--b & a salary & money paid for work, usually every month \\
\hline
4--e & a farmer & a person who grows crops or keeps animals \\
\hline
5--d & a factory & a place where goods are made in large quantities \\
\hline
6--f & unemployment & the situation of people who have no job \\
\hline
\end{tabular}
\end{center}
\textbf{Points de vigilance :} distinguer \eng{customer} (client) et \eng{customs} (douane) ; \eng{unemployment} est indénombrable : \eng{Unemployment is rising}, jamais \eng{unemployments}.
\end{corrige}

\begin{corrige}[Exercice 9 — Compréhension écrite type Bac]
\textbf{A. True or False (5 marks — 1 mark par réponse justifiée)}
\begin{enumerate}
\item \textbf{False.} \eng{“With a small loan from her aunt, she bought a sewing machine.”} (Elle a emprunté de l'argent ; elle n'était pas riche.)
\item \textbf{False.} \eng{“I worked longer hours than any of my friends, and I earned less money.”} (Elle gagnait \emph{moins}, pas plus.)
\item \textbf{True.} \eng{“Her clothes became famous because they were cheaper and more beautiful than those in the big shops.”}
\item \textbf{False.} \eng{“Young Cameroonians should not wait for the government to give them jobs.”}
\item \textbf{True.} \eng{“She hopes to export her clothes to Nigeria one day.”}
\end{enumerate}
\textbf{B. Questions (4 marks — 1 mark par réponse)}
\begin{enumerate}
\item \eng{She lives in Bamenda, in the North-West Region of Cameroon.}
\item \eng{Her aunt gave her the money (a small loan) to start her business.}
\item \eng{Six people work with her today: four young girls and two apprentices.}
\item \eng{Young people can learn that with courage and hard work, they can create their own jobs and build a better future.}
\end{enumerate}
\textbf{C. Language work (3 marks — 1 mark par réponse)}
\begin{enumerate}
\item Comparative of \eng{long} : \eng{longer} (\eng{“I worked longer hours”}) ; superlative of \eng{good} : \eng{best} (\eng{“The best way to fight unemployment”}).
\item Synonym of \eng{well-known} : \eng{famous}.
\item \eng{She does not employ four young girls.} (Négation au present simple : \eng{does not + base verbale} ; le \textbf{-s} disparaît car \eng{does} le porte déjà.)
\end{enumerate}
\textbf{Points de vigilance :} en partie A, une réponse sans citation du texte ne vaut pas le point complet ; en partie B, répondre avec ses propres mots et des phrases complètes ; attention à la négation : jamais \eng{she doesn't employs}.
\end{corrige}

\begin{corrige}[Exercice 10 — Traduction type Bac]
\textbf{Proposition de traduction (5 marks) :}

\eng{My brother is a taxi driver in Douala. He works harder than many people, but he earns little money. At the moment, he is saving to buy his own car. I think it is the most difficult job in our family, but he never complains.}

\textbf{Barème indicatif :} 1 point par segment correctement traduit (5 segments) ; une demi-point peut être retiré par erreur de grammaire majeure (temps, comparatif, article).
\begin{itemize}
\item « chauffeur de taxi » : \eng{a taxi driver} (article obligatoire devant un métier : jamais \eng{my brother is taxi driver}).
\item « plus dur que » : \eng{harder than} (comparatif régulier, pas \eng{more hard}).
\item « il gagne peu d'argent » : \eng{he earns little money} (\eng{money} indénombrable : \eng{little}, pas \eng{few}).
\item « il économise » : \eng{he is saving} (action en cours : present continuous ; \eng{to save money} = économiser).
\item « le métier le plus difficile » : \eng{the most difficult job} (adjectif long : \eng{the most}).
\item « il ne se plaint jamais » : \eng{he never complains} (place de \eng{never} : avant le verbe ; pas de double négation).
\end{itemize}
\textbf{Points de vigilance :} faux ami : « économiser » = \eng{to save}, pas \eng{to economize} au sens courant ; « se plaindre » = \eng{to complain}, sans préposition.
\end{corrige}

\begin{corrige}[Exercice 11 — Expression écrite type Bac]
\textbf{Proposition de rédaction modèle (environ 120 mots) :}

\begin{textebox}[Model paragraph]
The job I would like to do in the future is nursing, and I would love to work in a hospital in Yaoundé. First, I have always wanted to help sick people, because a nurse can save lives and bring comfort to families. Moreover, nursing is more useful to our community than many other jobs, since Cameroon needs more health workers in both towns and villages. However, I know it is one of the hardest professions: nurses work longer hours than office workers and their salary is not always the best. In conclusion, I believe nursing is the most rewarding job in the world, and I am ready to study hard to achieve my dream.
\end{textebox}

\textbf{Traduction du modèle :} « Le métier que je voudrais exercer plus tard est celui d'infirmier/infirmière, et j'aimerais travailler dans un hôpital de Yaoundé. D'abord, j'ai toujours voulu aider les malades, car un infirmier peut sauver des vies et réconforter les familles. De plus, ce métier est plus utile à notre communauté que beaucoup d'autres, car le Cameroun a besoin de plus de personnels de santé en ville comme au village. Cependant, je sais que c'est l'un des métiers les plus durs : les infirmiers travaillent plus longtemps que les employés de bureau et leur salaire n'est pas toujours le meilleur. En conclusion, je pense que c'est le métier le plus gratifiant du monde, et je suis prêt(e) à étudier dur pour réaliser mon rêve. »

\textbf{Barème indicatif (10 marks) :}
\begin{itemize}
\item Contenu (3 pts) : métier nommé, lieu indiqué, deux raisons claires.
\item Vocabulaire (2 pts) : lexique des métiers (\eng{nurse, salary, hospital, worker...}) correctement utilisé.
\item Grammaire (3 pts) : temps corrects ; au moins un comparatif (\eng{more useful than}, \eng{longer hours than}) et un superlatif (\eng{the most rewarding}, \eng{the best}).
\item Connecteurs et organisation (2 pts) : \eng{first, moreover, however, in conclusion} présents et bien placés ; paragraphe unique et cohérent.
\end{itemize}
\textbf{Points de vigilance :} respecter la longueur (environ 120 mots) ; ne pas écrire \eng{the most hardest} (double marque du superlatif) ; accorder le verbe avec le sujet (\eng{a nurse can}, \eng{nurses work}) ; éviter les phrases trop longues sans connecteur.
\end{corrige}

\newpage

\coursec{Fiche bilan}

\begin{popfigure}
\begin{center}\tcbox[colback=popblue,colframe=popblue,coltext=white,arc=9pt,boxsep=4pt,left=6pt,right=6pt]{\bfseries\small Chapter 2 Economic Life and Occupations}\end{center}
\vspace{-2pt}
\begin{tcbraster}[raster columns=3,raster equal height=rows,raster column skip=6pt,raster row skip=6pt,enhanced,blankest]
\begin{tcolorbox}[enhanced,colback=poporange,colframe=poporange,coltext=white,boxrule=0.9pt,arc=3pt,left=4pt,right=4pt,top=3pt,bottom=3pt,boxsep=1pt,halign=left]\bfseries\scriptsize Vocabulary jobs, trades, money, business\end{tcolorbox}
\begin{tcolorbox}[enhanced,colback=popgreen,colframe=popgreen,coltext=white,boxrule=0.9pt,arc=3pt,left=4pt,right=4pt,top=3pt,bottom=3pt,boxsep=1pt,halign=left]\bfseries\scriptsize Grammar present simple vs continuous\end{tcolorbox}
\begin{tcolorbox}[enhanced,colback=poppurple,colframe=poppurple,coltext=white,boxrule=0.9pt,arc=3pt,left=4pt,right=4pt,top=3pt,bottom=3pt,boxsep=1pt,halign=left]\bfseries\scriptsize Grammar present perfect already, yet, since\end{tcolorbox}
\begin{tcolorbox}[enhanced,colback=poppink,colframe=poppink,coltext=white,boxrule=0.9pt,arc=3pt,left=4pt,right=4pt,top=3pt,bottom=3pt,boxsep=1pt,halign=left]\bfseries\scriptsize Grammar comparatives and superlatives\end{tcolorbox}
\begin{tcolorbox}[enhanced,colback=popgold,colframe=popgold,coltext=white,boxrule=0.9pt,arc=3pt,left=4pt,right=4pt,top=3pt,bottom=3pt,boxsep=1pt,halign=left]\bfseries\scriptsize Reading exam-type comprehension\end{tcolorbox}
\begin{tcolorbox}[enhanced,colback=popblue!70,colframe=popblue!70,coltext=white,boxrule=0.9pt,arc=3pt,left=4pt,right=4pt,top=3pt,bottom=3pt,boxsep=1pt,halign=left]\bfseries\scriptsize Speaking job interview role-play\end{tcolorbox}
\begin{tcolorbox}[enhanced,colback=popmuted,colframe=popmuted,coltext=white,boxrule=0.9pt,arc=3pt,left=4pt,right=4pt,top=3pt,bottom=3pt,boxsep=1pt,halign=left]\bfseries\scriptsize Writing letter of application\end{tcolorbox}
\end{tcbraster}

\legende{Carte mentale de synthèse : les sept compétences-clés du chapitre 2.}
\end{popfigure}

\begin{bilan}
\textbf{1. Lexique essentiel du chapitre (à connaître par cœur)}

\begin{itemize}
  \item \eng{a job / an occupation} : un métier, une profession ; \eng{a trade} : un métier manuel ; \eng{a career} : une carrière.
  \item \eng{an employer / an employee} : un employeur / un employé ; \eng{a worker} : un travailleur ; \eng{staff} : le personnel.
  \item \eng{to hire / to employ} : embaucher ; \eng{to fire / to dismiss} : licencier ; \eng{to resign} : démissionner ; \eng{to retire} : prendre sa retraite.
  \item \eng{a salary / wages} : un salaire (mensuel / à l'heure) ; \eng{income} : le revenu ; \eng{to earn money} : gagner de l'argent.
  \item \eng{a customer} : un client ; \eng{a seller / a trader} : un vendeur / un commerçant ; \eng{goods} : des marchandises ; \eng{a market} : un marché.
  \item \eng{a bank account} : un compte en banque ; \eng{to save / to spend} : épargner / dépenser ; \eng{a loan} : un prêt ; \eng{a debt} : une dette.
  \item \eng{a farmer} : un agriculteur ; \eng{crops} : les cultures ; \eng{to harvest cocoa} : récolter le cacao ; \eng{a carpenter} : un menuisier ; \eng{a tailor} : un tailleur.
\end{itemize}

\textbf{2. Règles de grammaire à maîtriser}

\begin{center}
\begin{tabularx}{\linewidth}{|l|X|X|}
\hline
\rowcolor{popblueL} Point & Règle & Exemple \\
\hline
Present simple & Habitudes, faits permanents ; \textbf{-s} à la 3\textsuperscript{e} personne du singulier. & \eng{She works in a bank.} « Elle travaille dans une banque. » \\
\hline
Present continuous & Action en cours : \textbf{be + V-ing}. & \eng{He is selling fish now.} « Il vend du poisson en ce moment. » \\
\hline
Present perfect & Bilan, expérience, lien avec le présent : \textbf{have/has + participe passé} ; \eng{already, yet, since, for, never}. & \eng{I have worked here since 2020.} « Je travaille ici depuis 2020. » \\
\hline
Comparatifs & Court : \textbf{-er + than} ; long : \textbf{more + adj. + than}. & \eng{Farming is harder than trading.} « L'agriculture est plus dure que le commerce. » \\
\hline
Superlatifs & Court : \textbf{the + -est} ; long : \textbf{the most + adj.} & \eng{The most important job in the village.} « Le métier le plus important du village. » \\
\hline
Irréguliers & \eng{good / better / the best} ; \eng{bad / worse / the worst} ; \eng{much / more / the most}. & \eng{This shop is the best in town.} « Cette boutique est la meilleure de la ville. » \\
\hline
\end{tabularx}
\end{center}

\textbf{3. Méthodes express}

\begin{itemize}
  \item \textbf{Compréhension de l'écrit} : lire d'abord les questions, repérer les mots-clés dans le texte, justifier chaque réponse par une ligne du texte, répondre en phrases complètes.
  \item \textbf{Compréhension de l'oral} : lire les questions avant l'écoute, noter des mots (pas des phrases), se concentrer sur \eng{who, where, what, how much}.
  \item \textbf{Expression orale (entretien d'embauche)} : se présenter (\eng{My name is..., I am sixteen...}), parler de ses qualités (\eng{I am hard-working and punctual}), poser une question à la fin.
  \item \textbf{Expression écrite (lettre de candidature)} : adresse et date en haut, formule d'appel \eng{Dear Sir or Madam}, 3 paragraphes (motif — qualités — conclusion), formule finale \eng{Yours faithfully}.
\end{itemize}

\textbf{4. Pièges fréquents des francophones}

\begin{itemize}
  \item \eng{He works} et non \emph{he work} : ne pas oublier le \textbf{-s} de la 3\textsuperscript{e} personne.
  \item \eng{to look for a job} : chercher un emploi (ne pas oublier \eng{for}).
  \item \eng{money, information, news} sont indénombrables : jamais de \textbf{-s}.
  \item Faux ami : \eng{actually} signifie « en fait », pas « actuellement » (\eng{currently}).
  \item \eng{since} + point de départ (\eng{since 2019}) ; \eng{for} + durée (\eng{for five years}).
\end{itemize}
\end{bilan}

\begin{aretenir}[Checklist avant le Bac]
Cochez chaque case lorsque vous êtes sûr de maîtriser le point :

\begin{itemize}
  \item[$\square$] Je connais au moins vingt mots de vocabulaire sur les métiers, l'argent et le commerce.
  \item[$\square$] Je distingue \eng{a job} (dénombrable) et \eng{work} (indénombrable).
  \item[$\square$] Je conjugue le present simple sans oublier le \textbf{-s} à la 3\textsuperscript{e} personne du singulier.
  \item[$\square$] Je sais quand employer le present simple (habitude) et le present continuous (action en cours).
  \item[$\square$] Je forme le present perfect avec \eng{have/has} + participe passé et j'utilise correctement \eng{since} et \eng{for}.
  \item[$\square$] Je connais les participes passés irréguliers fréquents : \eng{earned, sold, bought, made, written, been}.
  \item[$\square$] Je forme les comparatifs et superlatifs, y compris \eng{better / the best} et \eng{worse / the worst}.
  \item[$\square$] Je sais répondre à des questions de compréhension en phrases complètes en justifiant par le texte.
  \item[$\square$] Je sais me présenter et parler de mes qualités lors d'un entretien d'embauche en anglais.
  \item[$\square$] Je connais la structure d'une lettre de candidature : adresse, date, \eng{Dear Sir or Madam}, \eng{Yours faithfully}.
  \item[$\square$] J'évite les faux amis : \eng{actually} (en fait), \eng{a salary} (un salaire), \eng{to assist} (assister au sens d'aider).
  \item[$\square$] Je relis toujours ma production pour vérifier les accords, les temps et l'orthographe.
\end{itemize}
\end{aretenir}

\findecours

\end{document}
